\chapter{自由传播、开放系统与量子测量}\label{chap:prop-measure}

前面已经求出电子离开相互作用区时的量子态。实验中的探测器通常位于更远的位置，电子还要传播一段距离才到达探测面。这段自由传播会改变不同动量分量的相对相位，即使能谱没有变化，电子在时间和空间中的分布也可能改变。要由出射态预言探测图样，我们先求传播后的态，再计算探测结果的概率。

记相互作用结束时刻为 $t_+$，其电子密度矩阵元为 $\rho_{\ell\ell'}^{(+)}$。以下自由传播中的 $t$ 是从 $t_+$ 起算的经过时间。若第 $\ell$ 个通道的自由能量为 $\mathcal E_\ell$，态矢振幅获得因子 $\exp(-\ii\mathcal E_\ell t/\hbar)$；密度矩阵由一个振幅和另一个振幅的共轭组成，因而
\begin{equation}
\boxed{
\rho_{\ell\ell'}(t)
=\rho_{\ell\ell'}^{(+)}
\exp\!\left[-\frac{\ii t}{\hbar}
(\mathcal E_\ell-\mathcal E_{\ell'})\right].}
\label{eq:sideband_coherence_free_propagation_master}
\end{equation}
当 $\ell=\ell'$ 时，相位因子等于一，各边带的概率保持不变。当 $\ell\ne\ell'$ 时，相对相位继续演化；在位置表象中把各分量叠加，便会出现随传播距离变化的干涉项。后面讨论的电子脉冲列和时间聚焦都来自这一过程。

传播过程中若还有材料涨落或未探测光子参与，电子就需要用约化密度算符描述。我们将在自由传播之后讨论环境怎样改变相干和涨落，最后引入测量算符，说明能谱、相位读出和条件测量分别能够获得哪些信息。这样，同一个输出信号中的传播变化、环境作用和探测器响应就都有明确的计算位置。

\section{自由传播与相空间}

自由粒子的每个动量分量只积累相位，因此在动量表象中传播最容易计算。位置密度则由所有分量的干涉决定，需要在传播后作 Fourier 叠加。本节先写出这一解，再对窄动量分布展开色散，从而得到波包展宽、纵向聚束和时间透镜的计算公式。

\subsection{自由传播与色散曲率矩阵}

设电子在自由传播开始时的动量振幅为 $f_s(\mathbf p)$，其中 $s$ 是自旋标签，归一化为 $\sum_s\int\dd^3p\,|f_s(\mathbf p)|^2=1$。由能量本征态的时间演化，正能电子波包在时刻 $t$ 为
\begin{equation*}
|\Psi(t)\rangle
=\sum_s\int\dd^3 p\,f_s(\mathbf p)
\ee^{-\ii E(\mathbf p)t/\hbar}|\mathbf p,s\rangle,
\qquad
E(\mathbf p)=\sqrt{\me^2c^4+c^2p^2}.
\end{equation*}
当实际谱支持足够窄时，局域传播方程由这一精确相位的受控展开得到。

设中心动量 $\mathbf p_0=p_0\mathbf e_z$，记 $\delta\mathbf p=\mathbf p-\mathbf p_0$。对 $E(\mathbf p)$ 作 Taylor 展开。一次导数给出波包中心的群速度；二次导数组成的矩阵称为 Hessian 矩阵，它描述不同动量分量的速度怎样偏离中心速度。保留到二阶，得
\begin{equation*}
E_2(\mathbf p)
=E_0+\mathbf v_0\cdot\delta\mathbf p
+\frac{\delta p_x^2+\delta p_y^2}{2m_\perp}
+\frac{\delta p_z^2}{2m_\parallel},
\end{equation*}
其中
\begin{equation*}
\mathbf v_0=\frac{c^2\mathbf p_0}{E_0},
\qquad
m_\perp=\frac{E_0}{c^2}=\gamma_0\me,
\qquad
m_\parallel=\gamma_0^3\me.
\end{equation*}
$m_\perp$ 与 $m_\parallel$ 分别是同一相对论色散沿横向和纵向方向的局部曲率尺度。传播误差由被舍弃的谱相位而不是单独的带宽控制。对实际动量支持 $\mathcal W_p$ 定义
\begin{equation}
\boxed{
\epsilon_{\rm disp}^{(\ge3)}(t)
\equiv
\sup_{\mathbf p\in\mathcal W_p}
\left|\frac{t}{\hbar}[E(\mathbf p)-E_2(\mathbf p)]\right|.}
\label{eq:dispersion_truncation_control}
\end{equation}
只有在 $\epsilon_{\rm disp}^{(\ge3)}\ll1$ 时，进入共动坐标 $\zeta=z-v_0t$ 后才得到二阶局域传播方程
\begin{equation}
\boxed{
\ii\hbar\partial_t\psi_2
=-\frac{\hbar^2}{2m_\perp}\nabla_\perp^2\psi_2
-\frac{\hbar^2}{2m_\parallel}\partial_\zeta^2\psi_2.}
\label{eq:quadratic_dispersion_local_propagation}
\end{equation}
令 $P_{\mathcal W_p}$ 投影到实际谱支持，精确与二阶传播子的差满足
\begin{equation}
\boxed{
\left\|[U_{\rm ex}(t)-U_2(t)]P_{\mathcal W_p}\right\|
\le\epsilon_{\rm disp}^{(\ge3)}(t).}
\label{eq:quadratic_dispersion_propagator_error_bound}
\end{equation}
该近似的误差由高阶色散相位随传播时间的累积控制，因此必须用无量纲相位而非孤立的哈密顿量系数判断失效。

固定总能量、以 $z$ 作为传播坐标时，同一个质量壳给出
\begin{equation*}
p_z(\mathbf p_\perp)=\sqrt{p_0^2-p_\perp^2}.
\end{equation*}
近轴条件首先要求 $p_\perp^2/p_0^2\ll1$，但长距离传播还要求四阶及以上纵向相位不可分辨。直接定义
\begin{equation}
\boxed{
\epsilon_{\rm par}^{(\ge4)}(z)
\equiv
\sup_{\mathbf p_\perp\in\mathcal W_\perp}
\left|\frac{z}{\hbar}
\left[\sqrt{p_0^2-p_\perp^2}-p_0+\frac{p_\perp^2}{2p_0}\right]\right|.}
\label{eq:paraxial_phase_control_parameter}
\end{equation}
在 $\epsilon_{\rm par}^{(\ge4)}\ll1$ 时，抽去载波 $\ee^{\ii p_0z/\hbar}$ 得到
\begin{equation}
\boxed{
\ii\hbar\partial_z\psi_{\rm par}
=-\frac{\hbar^2}{2p_0}\nabla_\perp^2\psi_{\rm par}.}
\label{eq:relativistic_paraxial_from_mass_shell}
\end{equation}
并有传播子误差界
\begin{equation}
\boxed{
\left\|[U_z^{\rm ex}-U_z^{\rm par}]P_{\mathcal W_\perp}\right\|
\le\epsilon_{\rm par}^{(\ge4)}(z).}
\label{eq:paraxial_propagator_error_bound}
\end{equation}
沿中心轨迹取 $z\simeq v_0t$ 时，$p_0/v_0=\gamma_0\me=m_\perp$，因此时间域横向 Hessian 矩阵与定态近轴方程是同一精确色散在不同边值问题下的二阶极限。

\begin{derivation}[从\eqnrefs{eq:quadratic_dispersion_local_propagation}到\eqnrefs{eq:quadratic_dispersion_propagator_error_bound}]
其适用条件是初态动量主要支撑在给定带宽内、$E(p)$ 在该区间足够光滑，并且显式传播相位误差小于一。 Taylor 展开写成

\begin{align*}
E(\mathbf p_0+\delta\mathbf p)
={}&E_0+DE(\mathbf p_0)[\delta\mathbf p]
+\frac12D^2E(\mathbf p_0)[\delta\mathbf p^{\otimes2}]\\
&+\frac1{3!}D^3E(\mathbf p_0)[\delta\mathbf p^{\otimes3}]+R_4.
\end{align*}
对 $E=\sqrt{\me^2c^4+c^2p^2}$，
\begin{equation*}
\partial_{p_i}E=\frac{c^2p_i}{E},
\qquad
\partial_{p_i}\partial_{p_j}E
=\frac{c^2}{E}\delta_{ij}-\frac{c^4p_ip_j}{E^3}.
\end{equation*}
在 $\mathbf p_0=p_0\mathbf e_z$ 上，横向 Hessian 元素为 $1/m_\perp=c^2/E_0$，纵向元素为
\begin{align*}
\frac1{m_\parallel}
&=\frac{c^2}{E_0}-\frac{c^4p_0^2}{E_0^3}\\
&=\frac{c^2}{E_0}\left(1-\frac{v_0^2}{c^2}\right)
=\frac{1}{\gamma_0^3\me}.
\end{align*}
把载波和群延迟从精确 Fourier 波包中提出。令 $\zeta=z-v_0t$，并写
\begin{equation*}
\Psi(\mathbf r,t)
=\ee^{\ii(p_0z-E_0t)/\hbar}
\psi_2(\mathbf r_\perp,\zeta,t).
\end{equation*}
在二阶色散下，包络为
\begin{equation*}
\psi_2
=\int\frac{\dd^3(\delta p)}{(2\pi\hbar)^3}
A(\delta\mathbf p)
\exp\!\left[
\frac{\ii}{\hbar}(\delta\mathbf p_\perp\cdot\mathbf r_\perp+\delta p_z\zeta)
-\frac{\ii t}{\hbar}
\left(\frac{\delta p_\perp^2}{2m_\perp}
+\frac{\delta p_z^2}{2m_\parallel}\right)\right].
\end{equation*}
对 $t$ 求导，群速度项已经被 $\zeta=z-v_0t$ 吸收，因此
\begin{align*}
\ii\hbar\partial_t\psi_2
&=\int\frac{\dd^3(\delta p)}{(2\pi\hbar)^3}
A(\delta\mathbf p)
\left(\frac{\delta p_\perp^2}{2m_\perp}
+\frac{\delta p_z^2}{2m_\parallel}\right)\ee^{\cdots}\\
&=-\frac{\hbar^2}{2m_\perp}\nabla_\perp^2\psi_2
-\frac{\hbar^2}{2m_\parallel}\partial_\zeta^2\psi_2,
\end{align*}
其中第二行使用 $\delta p_j \ee^{\ii\delta p_jx_j/\hbar}=-\ii\hbar\partial_{x_j}\ee^{\ii\delta p_jx_j/\hbar}$。这就是\eqnrefs{eq:quadratic_dispersion_local_propagation}。精确传播子与二阶传播子在动量表象都是乘法算符。定义余项相位
\begin{equation*}
\delta\varphi(\mathbf p,t)
=\frac{t}{\hbar}[E(\mathbf p)-E_2(\mathbf p)].
\end{equation*}
对任意归一化且支撑在 $\mathcal W_p$ 的波包，
\begin{align*}
\|[U_{\rm ex}(t)-U_2(t)]\psi\|^2
&=\int_{\mathcal W_p}\dd^3 p\,|A(\mathbf p)|^2
|\ee^{-\ii\delta\varphi}-1|^2\\
&\le\sup_{\mathbf p\in\mathcal W_p}|\ee^{-\ii\delta\varphi}-1|^2\\
&\le\sup_{\mathbf p\in\mathcal W_p}|\delta\varphi|^2.
\end{align*}
取平方根后得到\eqnrefs{eq:quadratic_dispersion_propagator_error_bound}。若已知三阶导数上界 $M_3=\sup_{\mathcal W_p}\|D^3E\|$，Taylor 余项还给出
\begin{equation*}
\epsilon_{\rm disp}^{(\ge3)}(t)
\le\frac{|t|M_3}{6\hbar}
\sup_{\mathcal W_p}|\delta\mathbf p|^3+O(|\delta\mathbf p|^4),
\end{equation*}
从而把“带宽小”转化为传播时间相关的相位误差界。
\end{derivation}
\begin{derivation}[从\eqnrefs{eq:relativistic_paraxial_from_mass_shell}到\eqnrefs{eq:paraxial_propagator_error_bound}]
设电子束以 $+z$ 为中心传播方向，并且横向动量与纵向失谐都小于中心动量。固定总能量时有 $p_z^2+p_\perp^2=p_0^2$；取沿 $+z$ 传播的正根，并把高阶角谱项保留为显式余项，

\begin{align*}
p_z
&=p_0\sqrt{1-\frac{p_\perp^2}{p_0^2}}\\
&=p_0\left[1-\frac12\frac{p_\perp^2}{p_0^2}
-\frac18\frac{p_\perp^4}{p_0^4}
+O\!\left(\frac{p_\perp^6}{p_0^6}\right)\right]\\
&=p_0-\frac{p_\perp^2}{2p_0}
-\frac{p_\perp^4}{8p_0^3}+O(p_\perp^6/p_0^5).
\end{align*}
把共同纵向载波抽出，写
\begin{equation*}
\Psi(\mathbf r_\perp,z)
=\ee^{\ii p_0z/\hbar}\psi_{\rm env}(\mathbf r_\perp,z),
\end{equation*}
其中包络的精确角谱表示为
\begin{equation*}
\psi_{\rm env}(\mathbf r_\perp,z)
=\int\frac{\dd^2 p_\perp}{(2\pi\hbar)^2}
\,A(\mathbf p_\perp)
\exp\!\left[\frac{\ii}{\hbar}
\mathbf p_\perp\!\cdot\!\mathbf r_\perp
+\frac{\ii}{\hbar}(p_z-p_0)z\right].
\end{equation*}
对 $z$ 求导时每个角谱分量给出
\begin{align*}
\ii\hbar\,\partial_z\psi_{\rm env}
&=\int\frac{\dd^2 p_\perp}{(2\pi\hbar)^2}
A(\mathbf p_\perp)
\bigl[p_0-p_z(\mathbf p_\perp)\bigr]
\ee^{\ii\mathbf p_\perp\cdot\mathbf r_\perp/\hbar}
\ee^{\ii(p_z-p_0)z/\hbar}\\
&=-\frac{\hbar^2}{2p_0}\nabla_\perp^2\psi_{\rm env}
+\mathcal R_{\rm par}(\mathbf r_\perp,z),
\end{align*}
其中残差被精确定义为
\begin{equation*}
\mathcal R_{\rm par}
\equiv
\int\frac{\dd^2 p_\perp}{(2\pi\hbar)^2}
A(\mathbf p_\perp)
\left[p_0-p_z(\mathbf p_\perp)-\frac{p_\perp^2}{2p_0}\right]
\ee^{\ii\mathbf p_\perp\cdot\mathbf r_\perp/\hbar}
\ee^{\ii(p_z-p_0)z/\hbar}.
\end{equation*}
当 $\mathcal R_{\rm par}$ 相对于保留项低于目标精度时，得到抛物传播方程\eqnrefs{eq:relativistic_paraxial_from_mass_shell}。其适用范围可以直接由精确传播子与抛物传播子的范数差量化。定义
\begin{equation*}
\delta\varphi(\mathbf p_\perp;z)
=\frac{z}{\hbar}
\left[p_z-p_0+\frac{p_\perp^2}{2p_0}\right].
\end{equation*}
两个传播子在角谱表象都为乘法算符，因此对任意支撑在 $\mathcal W_\perp$ 的归一化态，
\begin{align*}
\|[U_z^{\rm ex}-U_z^{\rm par}]\psi\|^2
&=\int_{\mathcal W_\perp}\dd^2 p_\perp\,
|A(\mathbf p_\perp)|^2
|\ee^{\ii\delta\varphi}-1|^2\\
&\le \sup_{\mathcal W_\perp}|\ee^{\ii\delta\varphi}-1|^2
\le \sup_{\mathcal W_\perp}|\delta\varphi|^2,
\end{align*}
其中用了 $|\ee^{\ii x}-1|=2|\sin(x/2)|\le|x|$。取平方根即得到\eqnrefs{eq:paraxial_propagator_error_bound}。由上面的展开还可读出主导误差
\begin{equation*}
\delta\varphi
=-\frac{z\,p_\perp^4}{8\hbar p_0^3}
+O\!\left(\frac{z p_\perp^6}{\hbar p_0^5}\right),
\end{equation*}
所以几何小角条件 $p_\perp/p_0\ll1$ 与长距离相位条件 $z p_\perp^4/(8\hbar p_0^3)\ll1$ 必须同时满足。
\end{derivation}

\subsection{能量格相干到时间密度的定量映射}\label{subsec:sideband-to-temporal-density-dp46}

前面的自由传播一般表达式还可以直接化成实验可读取的纵向密度。对\chapref{chap:external-pinem}定义的单频交换格，在中心同步条件 $v_0\kappa_z=\omega$ 下，相邻电子通道在共动坐标 $\zeta=z-v_0t$ 中相差纵向相位 $\kappa_z\zeta$。若只保留二阶纵向色散，则传播后的约化密度核给出
\begin{equation}
\boxed{
 \mathfrak n_e(\zeta,t)
 =\sum_{\ell,\ell'}
 \rho_{\ell\ell'}^{(+)}
 \exp\!\left[\ii(\ell-\ell')\kappa_z\zeta
 -\ii(\ell^2-\ell'^2)\omega_{\rm r}t\right].}
\label{eq:sideband_to_longitudinal_density_dp46}
\end{equation}
这里 $\mathfrak n_e$ 是把慢变纵向包络的共同归一化因子抽出后得到的无量纲周期密度；若需要绝对线密度，只需乘入入射包络线密度。\eqnrefs{eq:sideband_to_longitudinal_density_dp46} 显示：对角能谱 $\rho_{\ell\ell}$ 只给常数背景，所有时间结构都来自能量格非对角相干。

把谐波阶数定义为 $\Delta\ell=\ell-\ell'$，可将密度写成
\begin{equation}
 \mathfrak n_e(\zeta,t)
 =\sum_{\Delta\ell\in\mathbb Z}B_{\Delta\ell}(t)\ee^{\ii \Delta\ell\kappa_z\zeta},
 \qquad B_{-\Delta\ell}=B_{\Delta\ell}^*,
\label{eq:electron_density_harmonic_expansion_dp46}
\end{equation}
其中
\begin{equation}
\boxed{
 B_{\Delta\ell}(t)
 =\sum_{\ell}
 \rho_{\ell,\ell-\Delta\ell}^{(+)}
 \exp\!\left[-\ii(2\ell \Delta\ell-\Delta\ell^2)\omega_{\rm r}t\right].}
\label{eq:electron_bunching_harmonic_dp46}
\end{equation}
$B_{\Delta\ell}$ 是第 $\Delta\ell$ 阶电子聚束因子，也就是归一化密度在该空间频率上的复 Fourier 系数。$|B_{\Delta\ell}|$ 给出这一谐波的幅度，$\arg B_{\Delta\ell}$ 给出它相对于光学相位基准的位置；若密度只含常数项与一对 $\pm\Delta\ell$ 谐波，则对应余弦项的峰值幅度为 $2|B_{\Delta\ell}|$。尤其
\begin{equation*}
 B_0=\sum_\ell\rho_{\ell\ell}^{(+)}=1,
\end{equation*}
因此归一化由密度矩阵迹自动保证。

二阶反冲相位是 $\ell^2\omega_{\rm r}t$，故所有整数格点相位在
\begin{equation}
\boxed{
 T_{\rm T}=\frac{2\pi}{\omega_{\rm r}},
 \qquad
 L_{\rm T}=v_0T_{\rm T}
 =\frac{2\pi v_0}{\omega_{\rm r}}}
\label{eq:electron_talbot_time_length_dp46}
\end{equation}
后同时复现。$T_{\rm T}$ 与 $L_{\rm T}$ 分别是电子能量格的 Talbot 时间和沿中心轨迹的 Talbot 长度；二者均由同一个纵向色散质量
$m_\parallel=\gamma_0^3\me$ 和动量步长 $\hbar \kappa_z$ 决定：
\begin{equation*}
 L_{\rm T}
 =\frac{4\pi m_\parallel v_0}{\hbar \kappa_z^2}.
\end{equation*}
若使用同步关系 $\kappa_z=\omega/v_0$，还可写成
\begin{equation*}
 L_{\rm T}
 =\frac{4\pi m_\parallel v_0^3}{\hbar\omega^2}.
\end{equation*}
Talbot 复现还可以直接从整数格的数论相位看出。半个 Talbot 时间 $t=T_{\rm T}/2$ 时，
\begin{equation*}
 \ee^{-\ii\ell^2\omega_{\rm r}T_{\rm T}/2}
 =\ee^{-\ii\pi\ell^2}
 =\ee^{-\ii\pi\ell},
\end{equation*}
因为整数满足 $\ell^2-\ell=\ell(\ell-1)\in2\mathbb Z$。因此每个格点多出的因子 $\ee^{-\ii\pi\ell}$ 可吸收到空间相位中：
\begin{equation*}
 \ee^{\ii\ell\kappa_z\zeta}\ee^{-\ii\pi\ell}
 =\ee^{\ii\ell\kappa_z(\zeta-\pi/\kappa_z)}.
\end{equation*}
所以半 Talbot 时间并非原位复现，而是把整个周期密度平移半个空间周期 $\pi/\kappa_z$；到 $T_{\rm T}$ 时该平移再消失，得到原位复现。这个结果来自离散二次谱相位本身，与 PINEM 初始边带振幅的具体形状无关。

实际电子只占据有限边带，并且精确相对论色散还含三阶及更高阶谱相位。若有效占据范围为 $|\ell|\le\ell_{\max}$，把一个 Talbot 长度内所有三阶以上附加相位的最大值定义为
\begin{equation}
 \boxed{
 \epsilon_T^{(\ge3)}
 \equiv
 \max_{|\ell|\le\ell_{\max}}
 \left|\varphi_\ell^{(\ge3)}(L_T)\right|.}
 \label{eq:electron-talbot-higher-order-control-dp76}
\end{equation}
只有 $\epsilon_T^{(\ge3)}\ll1$ 时，二次谱相位给出的整数复现才保持高对比度；否则必须回到完整相对论传播相位逐个边带计算。

由 $m_\parallel=\gamma_0^3\me$ 可直接读出：提高中心 Lorentz 因子会增大纵向色散质量，因而在固定 $\kappa_z$ 下减小 $\omega_{\rm r}=\hbar\kappa_z^2/(2m_\parallel)$；形成同一二次谱相位所需的漂移时间与距离随之增加。

对理想经典 PINEM 输出
\begin{equation*}
 c_\ell^{(+)}
 =J_\ell(2|\beta|)
 \ee^{\ii\ell\varphi_\beta},
 \qquad
 \rho_{\ell\ell'}^{(+)}=c_\ell^{(+)}c_{\ell'}^{(+)*},
\end{equation*}
在理想双边 PINEM 输出并采用前述二阶纵向色散时，\eqnrefs{eq:electron_bunching_harmonic_dp46} 变为完全可计算的闭式级数
\begin{equation}
\boxed{
 B_{\Delta\ell}^{\rm PINEM}(t)
 =\ee^{\ii \Delta\ell\varphi_\beta}
 \sum_{\ell=-\infty}^{\infty}
 J_\ell(2|\beta|)J_{\ell-\Delta\ell}(2|\beta|)
 \ee^{-\ii(2\ell \Delta\ell-\Delta\ell^2)\omega_{\rm r}t}.}
\label{eq:pinem_bunching_exact_series_dp46}
\end{equation}
它给出“耦合强度 $|\beta|$、漂移距离 $L\simeq v_0t$、反冲频率 $\omega_{\rm r}$”到任意密度谐波的直接前向映射，不再需要用“压缩点”“最佳聚束”作定性判断。若实验只读取第一谐波，可先定义第一谐波的相对调制幅度
\begin{equation}
\boxed{\mathcal M_1(t)=2|B_1(t)|,}
\label{eq:first_harmonic_visibility_dp46}
\end{equation}
它一般不等于限制在 $[0,1]$ 的 Michelson 可见度；只有高阶谐波可忽略时二者才一致。完整密度对比度应直接由\eqnrefs{eq:sideband_to_longitudinal_density_dp46} 的最大、最小值计算。

若存在纯退相干而不改变边带占据，可把第 $\Delta\ell$ 条副对角线写成
\begin{equation*}
 \rho_{\ell,\ell-\Delta\ell}^{(+)}
 \longrightarrow
 \mathsf D_{\Delta\ell}^{\rm coh}\rho_{\ell,\ell-\Delta\ell}^{(+)},
 \qquad 0\le|\mathsf D_{\Delta\ell}^{\rm coh}|\le1.
\end{equation*}
于是
\begin{equation}
\boxed{B_{\Delta\ell}(t)\longrightarrow\mathsf D_{\Delta\ell}^{\rm coh} B_{\Delta\ell}(t),}
\label{eq:bunching_dephasing_transfer_dp46}
\end{equation}
而 $B_0$ 不变。这样开放系统中的相干损失直接变成可测密度谐波的乘性衰减，从而把退相干、层析和 Fisher 信息都连接到同一可测密度谐波。

\begin{derivation}[从\eqnrefs{eq:sideband_to_longitudinal_density_dp46}到\eqnrefs{eq:pinem_bunching_exact_series_dp46}]
先把共同的慢变纵向包络抽出，只处理一个同步周期内的离散能量格。取
\[
L_\zeta\equiv\frac{2\pi}{\kappa_z},
\qquad
\langle\zeta|\ell\rangle
=\frac{1}{\sqrt{L_\zeta}}
\ee^{+\ii\ell\kappa_z\zeta},
\qquad 0\le\zeta<L_\zeta,
\]
则 $\{|\ell\rangle\}_{\ell\in\mathbb Z}$ 在该周期盒中正交归一：
\begin{align*}
\langle\ell'|\ell\rangle
&=\frac1{L_\zeta}
\int_0^{L_\zeta}\dd\zeta\,
\ee^{+\ii(\ell-\ell')\kappa_z\zeta}\\
&=\delta_{\ell\ell'}.
\end{align*}
盒归一化的位置密度为
$\langle\zeta|\rho(t)|\zeta\rangle$；正文使用的无量纲周期密度定义为
\[
\mathfrak n_e(\zeta,t)
\equiv L_\zeta\langle\zeta|\rho(t)|\zeta\rangle,
\]
所以共同的 $L_\zeta^{-1}$ 因子不会进入最终谐波系数。

在二阶纵向色散下，共动表象已经移除了 $E_0+\ell\hbar\omega$ 的常数与线性相位，只留下反冲项 $\hbar\omega_{\rm r}\ell^2$。因此正文\eqnrefs{eq:sideband_coherence_free_propagation_master} 给
\[
\rho_{\ell\ell'}(t)
=\rho_{\ell\ell'}^{(+)}
\ee^{-\ii(\ell^2-\ell'^2)\omega_{\rm r}t}.
\]
在位置密度中插入两次能量格完备关系：
\begin{align*}
\mathfrak n_e(\zeta,t)
&=L_\zeta
\sum_{\ell,\ell'}
\langle\zeta|\ell\rangle
\rho_{\ell\ell'}(t)
\langle\ell'|\zeta\rangle\\
&=\sum_{\ell,\ell'}
\rho_{\ell\ell'}^{(+)}
\ee^{+\ii(\ell-\ell')\kappa_z\zeta}
\ee^{-\ii(\ell^2-\ell'^2)\omega_{\rm r}t},
\end{align*}
这就是正文\eqnrefs{eq:sideband_to_longitudinal_density_dp46}。

为了把双重和式按可测空间谐波重新分组，定义
\[
r\equiv\Delta\ell=\ell-\ell',
\qquad
\ell'=\ell-r.
\]
二次相位差逐项化为
\begin{align*}
\ell^2-\ell'^2
&=\ell^2-(\ell-r)^2\\
&=\ell^2-(\ell^2-2\ell r+r^2)\\
&=2\ell r-r^2.
\end{align*}
因此
\begin{align*}
\mathfrak n_e(\zeta,t)
&=\sum_{r\in\mathbb Z}
\ee^{+\ii r\kappa_z\zeta}
\sum_{\ell\in\mathbb Z}
\rho_{\ell,\ell-r}^{(+)}
\ee^{-\ii(2\ell r-r^2)\omega_{\rm r}t}\\
&\equiv
\sum_{r\in\mathbb Z}B_r(t)
\ee^{+\ii r\kappa_z\zeta},
\end{align*}
从而
\[
B_r(t)
=\sum_\ell
\rho_{\ell,\ell-r}^{(+)}
\ee^{-\ii(2\ell r-r^2)\omega_{\rm r}t},
\]
即正文\eqnrefs{eq:electron_bunching_harmonic_dp46}。密度矩阵厄米性
$\rho_{ab}=\rho_{ba}^*$ 再给
\begin{align*}
B_{-r}(t)
&=\sum_\ell
\rho_{\ell,\ell+r}^{(+)}
\ee^{+\ii(2\ell r+r^2)\omega_{\rm r}t}\\
&=\left[
\sum_j
\rho_{j,j-r}^{(+)}
\ee^{-\ii(2jr-r^2)\omega_{\rm r}t}
\right]^*
=B_r^*(t),
\end{align*}
其中第二行取 $j=\ell+r$。所以 $\mathfrak n_e$ 自动为实函数；而
\[
B_0=\sum_\ell\rho_{\ell\ell}^{(+)}
=\operatorname{Tr}\rho^{(+)}=1
\]
保证周期平均密度归一化。

最后代入理想经典 PINEM 纯态振幅
\[
c_\ell^{(+)}
=J_\ell(2|\beta|)\ee^{+\ii\ell\varphi_\beta},
\]
有
\begin{align*}
\rho_{\ell,\ell-r}^{(+)}
&=c_\ell^{(+)}c_{\ell-r}^{(+)*}\\
&=J_\ell(2|\beta|)J_{\ell-r}(2|\beta|)
\ee^{+\ii r\varphi_\beta}.
\end{align*}
代入 $B_r(t)$ 并把与 $\ell$ 无关的相位提出求和号，得到
\[
\boxed{
B_r^{\rm PINEM}(t)
=\ee^{+\ii r\varphi_\beta}
\sum_{\ell=-\infty}^{\infty}
J_\ell(2|\beta|)J_{\ell-r}(2|\beta|)
\ee^{-\ii(2\ell r-r^2)\omega_{\rm r}t},}
\]
即正文\eqnrefs{eq:pinem_bunching_exact_series_dp46}。
\end{derivation}

\subsection{从能谱相干算出电子脉冲列}
\label{sec:electron-temporal-bunching-worked}

在能谱上看到一串边带以后，还不能仅凭各峰的高度确定电子的时间分布；不同边带的相位决定它们在何时相长干涉。下面把上一节的密度矩阵传播式用于 PINEM 输出，求出周期密度及其 Fourier 谐波。令 $\theta_\zeta=\kappa_z\zeta$ 表示一个空间周期内的位置，$s_{\rm d}=\omega_{\rm r}t\simeq\omega_{\rm r}L/v_0$ 表示传播中积累的反冲相位；两者均无量纲。相对周期密度为
\begin{equation}
 \mathfrak n_e(\theta_\zeta,s_{\rm d})
 =\sum_{\ell m}\rho_{\ell m}^{(+)}
 \ee^{\ii(\ell-m)\theta_\zeta-\ii(\ell^2-m^2)s_{\rm d}},\qquad
 \frac1{2\pi}\int_0^{2\pi}\mathfrak n_e\dd\theta_\zeta=1.
 \label{eq:electron-bunching-normalized-density}
\end{equation}
$\mathfrak n_e$ 是相对于慢包络背景的无量纲密度，不是单位体积内的电子数。在固定探测面，相对于中心电子的到达延迟为 $\tau_{\rm d}=t-L/v_0$，所以 $\theta_\zeta=-\omega\tau_{\rm d}$；正的到达延迟对应负的共动坐标。把 $s_{\rm d}$ 在一个光学周期内视为常数，还要求实际占据阶数上 $\ell^2\omega_{\rm r}|\tau_{\rm d}|\ll1$。

对任意混态，$\rho_{\ell m}^{(+)}$ 的对角部分给常背景，非对角部分才给随 $\theta_\zeta$ 变化的项。若能谱为 $P_\ell$ 而相位完全随机，则 $\rho_{\ell m}=P_\ell\delta_{\ell m}$，无论传播多远都有 $\mathfrak n_e=1$。相反，纯 PINEM 态的 $c_\ell=J_\ell(2|\beta|)\ee^{\ii\ell\varphi_\beta}$ 保留确定的交叉项。刚离开相互作用区时
\begin{equation*}
 \sum_\ell c_\ell\ee^{\ii\ell\theta_\zeta}
 =\exp[2\ii|\beta|\sin(\theta_\zeta+\varphi_\beta)],
\end{equation*}
模平方仍为一；后续二次色散才把相位调制变为密度调制。产生能谱边带与产生时间脉冲是不同步骤。

将周期密度写为 $\mathfrak n_e=\sum_d B_d\ee^{\ii d\theta_\zeta}$，其 Fourier 系数为
\[
B_d=\frac1{2\pi}\int_0^{2\pi}\dd\theta_\zeta\,
\mathfrak n_e(\theta_\zeta,s_{\rm d})\ee^{-\ii d\theta_\zeta}.
\]
把密度的双重求和代入，角积分只留下 $\ell-m=d$ 的项。这解释了为什么第 $d$ 阶密度谐波需要所有相隔 $d$ 条边带的相干共同贡献。对上述纯 PINEM 输入，Bessel 和式进一步给出
\begin{equation}
 B_d(s_{\rm d})=(-\ii)^d\ee^{\ii d\varphi_\beta}
 J_d[4|\beta|\sin(ds_{\rm d})],\qquad d\in\mathbb Z.
 \label{eq:electron-pinem-bunching-closed}
\end{equation}
$d$ 是被读取的密度谐波阶数，不是某一个末态边带编号；$B_d$ 汇总所有相隔 $d$ 阶的相干。该式保持 $B_0=1$、$B_{-d}=B_d^*$，并在 $s_{\rm d}=0$ 时使所有非零阶谐波消失。

\begin{exbox}[第一谐波最强的距离：弱调制与强调制]
给定 $b=|\beta|>0$，求第一谐波首次达到最大值的漂移距离。由闭式，问题化为在 $0\le s_{\rm d}\le\pi/2$ 内最大化
$|B_1|=|J_1(4b\sin s_{\rm d})|$。

令 $x_*$ 为 $J_1'(x)=0$ 的第一个正根。递推关系给
$J_1'(x)=[J_0(x)-J_2(x)]/2$，所以该常数也由 $J_0(x_*)=J_2(x_*)$ 定义；$x_*\simeq1.84118$。$J_1$ 的首个正极大值也是其绝对最大值。现在只需检查驱动能否达到这个参数：
\begin{equation*}
 s_{\rm d,*}=\begin{cases}
 \pi/2,&4b<x_*,\\
 \arcsin[x_*/(4b)],&4b\ge x_*,
 \end{cases}\qquad
 L_* =\frac{v_0s_{\rm d,*}}{\omega_{\rm r}}.
\end{equation*}
弱调制时可用 $J_1(x)\simeq x/2$，故 $|B_1|\simeq2b|\sin s_{\rm d}|$，首个最优位置是 $L_T/4$。强调制时展开 $\arcsin x\simeq x$，得
$L_*\simeq v_0x_*/(4b\omega_{\rm r})$，最优距离随场幅近似反比缩短。

达到 $x_*$ 后，第一谐波上限为 $|B_1|=J_1(x_*)\simeq0.581865$；对应 $\mathcal M_1=2|B_1|>1$。这不违反密度正性，因为 $\mathcal M_1$ 是 Fourier 调制幅度。若只保留 $1+2\Re(B_1\ee^{\ii\theta_\zeta})$，反而会出现负值，说明必须保留更高谐波。第一谐波最大、完整脉冲最窄和 Talbot 复现是三个不同条件。
\end{exbox}

到达时间抖动会改变不同次电子相对于驱动的相位。设逐发误差 $\delta t$ 为零均值高斯变量，标准差为 $\sigma_{t,\rm jit}$，且仅引起相位平移、不改变脉冲包络耦合，则
\begin{equation}
 \overline B_d=B_d\,\langle\ee^{\ii d\omega\delta t}\rangle
 =B_d\exp[-d^2\omega^2\sigma_{t,\rm jit}^2/2].
 \label{eq:electron-bunching-jitter-transfer}
\end{equation}
这就是为何高阶时间结构比总能谱更容易消失：抖动不改变 $P_\ell$，却按 $d^2$ 更快抹去高阶相干。探测器自身的高斯时间响应也给相同形式的谐波衰减，但它只模糊记录，不能据此断言到达探测器前的量子态已经退相干。

\begin{exbox}[从所需谐波反推时间抖动容差]
要求第 $d\ne0$ 阶谐波至少保留理想幅度的比例 $0<\eta_d<1$。由高斯衰减因子取对数，直接得到
\begin{equation*}
 \exp[-d^2\omega^2\sigma_{t,\rm jit}^2/2]\ge\eta_d
 \ \Longleftrightarrow\ 
 \sigma_{t,\rm jit}\le\frac{\sqrt{2\ln(1/\eta_d)}}{|d|\omega}.
\end{equation*}
因此分辨越高的谐波，允许的抖动按 $1/|d|$ 缩小。若给定同一保留阈值 $\eta$，可用谐波阶数的上限为
\begin{equation*}
 d_{\max}=\left\lfloor
 \frac{\sqrt{2\ln(1/\eta)}}{\omega\sigma_{t,\rm jit}}
 \right\rfloor,
\end{equation*}
其中向下取整表示取不超过括号内实数的最大整数。

若探测器另有独立高斯时间响应宽度 $\sigma_{t,\rm det}$，两个特征函数相乘，测得谐波为
\begin{equation*}
 B_d^{\rm obs}=B_d\exp\!\left[-\frac{d^2\omega^2}{2}
 (\sigma_{t,\rm jit}^2+\sigma_{t,\rm det}^2)\right].
\end{equation*}
因此仅凭谐波衰减，不能区分束流到达抖动与探测时间模糊。若 $B_d\ne0$，不同 $d$ 的 $-\ln|B_d^{\rm obs}/B_d|/d^2$ 应为同一常数；偏离这个关系说明单一高斯模糊模型不充分。这给出不依赖某组特定时间数值的实验检查。
\end{exbox}

\figref{fig:electron-bunching-worked} 比较相干态、相位混合以及有限抖动后的计算结果。
\begin{figure}[htbp]
 \centering
 \includegraphics[width=0.92\linewidth]{electron_bunching_worked.pdf}
 \caption{能谱到时间密度需要非对角相干。数值图取 $|\beta|=1.20$、$\SI{200}{\kilo\electronvolt}$ 电子及 $\SI{800}{\nano\metre}$ 同步光学步长。左图用相同的 PINEM 占据比较纯态、完全相位混合和有限到达抖动在 $L_*$ 的周期密度；右图给出前三阶聚束因子随漂移的变化。实空间相位与探测面到达延迟满足 $\theta_\zeta=-\omega\tau_{\rm d}$。}
 \label{fig:electron-bunching-worked}
\end{figure}

\begin{derivation}[从副对角线级数到 Bessel 闭式]
令 $x=2|\beta|$。先把\eqnrefs{eq:pinem_bunching_exact_series_dp46} 中的阶差固定为 $d$，而不截断其中的边带和：
\begin{equation*}
 B_d=\ee^{\ii d\varphi_\beta+\ii d^2s_{\rm d}}
 \sum_{\ell\in\mathbb Z}J_\ell(x)J_{\ell-d}(x)
 \ee^{-2\ii\ell ds_{\rm d}}.
\end{equation*}
对两个实 Bessel 系数分别使用互为共轭的 Fourier 表示，
\begin{equation*}
 J_\ell(x)=\frac1{2\pi}\int_0^{2\pi}\ee^{\ii x\sin u-\ii\ell u}\dd u,
 \qquad
 J_{\ell-d}(x)=\frac1{2\pi}\int_0^{2\pi}
 \ee^{-\ii x\sin v+\ii(\ell-d)v}\dd v.
\end{equation*}
此时所有 $\ell$ 依赖都集中在
$\sum_\ell\exp[-\ii\ell(u-v+2ds_{\rm d})]$。
它是 $2\pi$ 倍的周期 Dirac $\delta$，要求 $v=u+2ds_{\rm d}$ 模 $2\pi$；该等式按周期分布理解，也可先加 $\ee^{-\epsilon\ell^2}$ 收敛因子再令 $\epsilon\to0^+$。因此双角积分只剩一个。使用
\begin{equation*}
 \sin u-\sin(u+2ds_{\rm d})
 =-2\sin(ds_{\rm d})\cos(u+ds_{\rm d}),
\end{equation*}
并把剩余角平移为 $u+ds_{\rm d}$，外部的 $\ee^{\ii d^2s_{\rm d}}$ 恰与平移产生的相位抵消，得到
\begin{equation*}
 B_d=\frac{\ee^{\ii d\varphi_\beta}}{2\pi}
 \int_0^{2\pi}\dd u\,\ee^{-\ii du}
 \exp[-2\ii x\sin(ds_{\rm d})\cos u].
\end{equation*}
使用 $\exp(-\ii y\cos u)=\sum_r(-\ii)^rJ_r(y)\ee^{\ii ru}$，角积分只保留 $r=d$，便得到\eqnrefs{eq:electron-pinem-bunching-closed}。最后对抖动相位的高斯特征函数取平均，给\eqnrefs{eq:electron-bunching-jitter-transfer}。这两个步骤分别使用确定性自由传播和不同实验重复之间的统计平均，不能把后者当作额外的确定性传播相位。
\end{derivation}


\subsection{Wigner--Moyal 输运}

在位置表象中容易看出波包有多宽，在动量表象中容易看出它含有哪些速度分量。若要理解一个带初始啁啾的波包为何先压缩再展宽，就需要同时保留位置与动量的相关。Wigner 表示把密度算符改写为位置 $z$ 和动量 $p$ 的函数，适合描述这种相关。它保留全部密度矩阵信息，但不必处处非负，因此需要从密度矩阵定义，而不能直接假设它是经典联合概率。

在纵向坐标表象中定义密度核
\[
\rho(z_1,z_2)\equiv\langle z_1|\hat\rho|z_2\rangle,
\qquad
\operatorname{Tr}\hat\rho=\int\dd z\,\rho(z,z)=1.
\]
对角元给位置概率，非对角元携带空间相干。为了同时保留平均位置与两点相干距离，引入中心坐标 $z=(z_1+z_2)/2$ 和相对坐标 $s=z_1-z_2$，并对 $s$ 做傅里叶变换：
\begin{equation}
\boxed{
W(z,p)
\equiv\frac{1}{2\pi\hbar}
\int_{-\infty}^{\infty}\dd s\,
\ee^{-\ii ps/\hbar}
\rho\!\left(z+\frac{s}{2},z-\frac{s}{2}\right).}
\label{eq:electron_wigner}
\end{equation}
其中 $W$ 的 SI 量纲是作用量$^{-1}$，使 $W\,\dd z\,\dd p$ 无量纲。它不是经典联合概率：相干交叉项可以使 $W$ 局部为负。但它的两个边缘严格恢复可测概率，
\begin{equation}
\boxed{
\int\dd p\,W(z,p)=\langle z|\hat\rho|z\rangle,
\qquad
\int\dd z\,W(z,p)=\langle p|\hat\rho|p\rangle.}
\label{eq:wigner_marginals}
\end{equation}
因此 Wigner 负值并不意味着“负概率事件”；它记录的是不能由正经典相空间分布表示的非对角相干。

波函数改用 Wigner 表示以后，Hamiltonian 和其他算符也需要相应的函数表示。对算符 $\hat A$，定义它的\emph{Weyl 符号}
\[
A_W(z,p)=\int\dd s\,\ee^{-\ii ps/\hbar}
\left\langle z+\frac{s}{2}\right|\hat A\left|z-\frac{s}{2}\right\rangle.
\]
注意此处不含 $1/(2\pi\hbar)$；密度算符的 Weyl 符号因此为 $\rho_W=2\pi\hbar W$。相应地，期待值为 $\langle\hat A\rangle=\int\dd z\,\dd p\,A_W W$。位置和动量算符的符号分别就是 $z$ 和 $p$。

两个算符的乘积一般不能用两个符号的普通乘积代替，因为 $\hat z\hat p$ 与 $\hat p\hat z$ 不同，而普通函数 $zp$ 与 $pz$ 相同。正确的对应乘法称为\emph{Moyal 星积}，定义为
\begin{equation}
\boxed{
(\hat A\hat B)_W=A_W\star B_W
=A_W\exp\!\left[
\frac{\ii\hbar}{2}
\left(
\overleftarrow\partial_z\overrightarrow\partial_p
-\overleftarrow\partial_p\overrightarrow\partial_z
\right)\right]B_W.}
\label{eq:weyl_star_product_derived}
\end{equation}
左箭头导数作用于左侧函数，右箭头导数作用于右侧函数。例如，$z\star p=zp+\ii\hbar/2$，而 $p\star z=pz-\ii\hbar/2$，相减正好恢复 $[\hat z,\hat p]=\ii\hbar$。指数中的高阶导数进一步保留一般算符乘法的量子修正。于是 Liouville--von Neumann 方程
\[
\ii\hbar\partial_t\hat\rho=[\hat H,\hat\rho]
\]
在相空间中精确等价于
\begin{equation}
\boxed{
\partial_tW=\{H_W,W\}_{\rm M}
\equiv\frac{1}{\ii\hbar}(H_W\star W-W\star H_W).}
\label{eq:moyal_general}
\end{equation}
当由 $\hbar$ 设置的相空间单元相对于 $H_W$ 与 $W$ 的变化尺度足够小时，星积展开的最低阶给出 Poisson 括号；经典 Liouville 输运因此是\eqnrefs{eq:moyal_general} 的受控长尺度极限。

对精确自由相对论电子，哈密顿量只依赖动量，
\[
H_W(p)=E(p)=\sqrt{\me^2c^4+c^2p^2},
\qquad
v(p)=E'(p)=\frac{c^2p}{E(p)}.
\]
Moyal 方程于是化成奇数空间导数级数
\begin{equation}
\boxed{
\partial_tW
=-E'(p)\partial_zW
+\frac{\hbar^2}{24}E'''(p)\partial_z^3W
-\frac{\hbar^4}{1920}E^{(5)}(p)\partial_z^5W+\cdots.}
\label{eq:relativistic_free_moyal}
\end{equation}
第一项就是沿精确群速度 $v(p)$ 的经典平流；其余项衡量“非二次色散 + 有限量子相空间分辨率”造成的量子修正。若只保留中心动量 $p_0$ 附近的二阶色散
\[
E(p_0+\delta p)
\simeq E_0+v_0\delta p+\frac{\delta p^2}{2m_\parallel},
\qquad
m_\parallel\equiv\frac{1}{E''(p_0)}=\gamma_0^3\me,
\]
所有三阶以上动量导数都消失，所以 Moyal 方程在这个二次模型内\emph{精确}退化成相空间剪切：
\begin{equation}
\boxed{
\partial_tW+
\left(v_0+\frac{p-p_0}{m_\parallel}\right)\partial_zW=0,
\qquad
W(z,p,t)=
W_0\!\left(z-\left[v_0+\frac{p-p_0}{m_\parallel}\right]t,p\right).}
\label{eq:wigner_shear}
\end{equation}

自由传播不改变每个动量分量本身，却让不同动量以不同群速度前进。于是相空间分布沿 $z$ 方向剪切；所谓“电子啁啾压缩”就是先制备一张倾斜的相空间分布，再让这张分布被自由剪切到位置投影最窄的时刻。



把这个剪切写成矩阵，就可以直接计算位置、动量方差以及它们的协方差怎样变化。

在随中心电子移动的坐标
\[
\zeta\equiv z-v_0t,
\qquad
\delta p\equiv p-p_0
\]
中，二阶自由哈密顿量的线性漂移已经被消去，留下
\[
H_{\rm com}=\frac{\delta p^2}{2m_\parallel}.
\]
因此 Heisenberg 演化为
\begin{equation}
\boxed{
\zeta(t)=\zeta(0)+\frac{t}{m_\parallel}\delta p(0),
\qquad
\delta p(t)=\delta p(0).}
\label{eq:electron-shear-heisenberg}
\end{equation}
这与\eqnrefs{eq:wigner_shear} 是同一剪切在算符与相空间中的两种写法。

定义中心化二阶矩
\[
\sigma_z^2\equiv\langle\zeta^2\rangle,
\qquad
\sigma_p^2\equiv\langle\delta p^2\rangle,
\qquad
C_{zp}\equiv\frac12\langle\zeta\delta p+\delta p\zeta\rangle.
\]
\eqnrefs{eq:electron-shear-heisenberg} 给出
\begin{equation}
\boxed{
\sigma_z^2(t)
=\sigma_z^2(0)+\frac{2t}{m_\parallel}C_{zp}(0)
+\frac{t^2}{m_\parallel^2}\sigma_p^2.}
\label{eq:electron_shear_variance}
\end{equation}
负协方差 $C_{zp}<0$ 表示“前后位置与动量反相关”，自由剪切会先减小位置宽度。最佳时间与最小宽度为
\begin{equation}
\boxed{
 t_*=-m_\parallel\frac{C_{zp}(0)}{\sigma_p^2},
\qquad
\sigma_{z,\min}^2
=\sigma_z^2(0)-\frac{C_{zp}^2(0)}{\sigma_p^2}
\ge\frac{\hbar^2}{4\sigma_p^2}.}
\label{eq:electron_chirp_compression_bound}
\end{equation}
所以啁啾不会压缩相空间面积；它只是把已有的相空间椭圆剪切到位置投影最窄。$m_\parallel$ 决定焦点出现的传播时间，却不改变固定 $\sigma_p$ 下的 Heisenberg 极限。

无初始 $z$--$p$ 相关的最小不确定度高斯波包是最简单的基准。若初始纵向宽度为 $\sigma_{z0}$，则二阶传播保持高斯形状，并有
\[
\sigma_z^2(t)
=\sigma_{z0}^2+
\left(\frac{\hbar t}{2m_\parallel\sigma_{z0}}\right)^2.
\]
横向二次哈密顿量为 $H_\perp=(\delta p_x^2+\delta p_y^2)/(2m_\perp)$。若取其中任一横向方向 $x$，初态同样为无相关最小不确定度高斯，宽度 $\sigma_{x0}$ 与动量宽度满足 $\sigma_{p_x}=\hbar/(2\sigma_{x0})$，则
\[
\sigma_x^2(t)
=\sigma_{x0}^2+
\left(\frac{\hbar t}{2m_\perp\sigma_{x0}}\right)^2.
\]
定义展宽时间为使第二项与初始方差相等的时间，
\[
t_{\rm disp,\parallel}
=\frac{2m_\parallel\sigma_{z0}^2}{\hbar},
\qquad
t_{\rm disp,\perp}
=\frac{2m_\perp\sigma_{x0}^2}{\hbar}.
\]
对相同初始空间宽度 $\sigma_{x0}=\sigma_{z0}$，
\[
\frac{t_{\rm disp,\parallel}}{t_{\rm disp,\perp}}
=\frac{m_\parallel}{m_\perp}=\gamma_0^2.
\]
对高斯初态，前述二阶色散可以解析求解，并直接给出展宽与压缩条件。

二阶模型的适用域由精确色散中未被 $E_2$ 吸收的余项决定。令 $\delta p=p-p_0$。Taylor 定理给出精确分解
\[
E(p_0+\delta p)
=E_0+v_0\delta p+\frac{\delta p^2}{2m_\parallel}
+\frac{E'''_0}{6}\delta p^3+R_4(\delta p),
\qquad
E'''_0=-\frac{3v_0}{\gamma_0^4\me^2c^2},
\]
其中积分余项定义为
\[
R_4(\delta p)
=\frac{\delta p^4}{6}
\int_0^1\dd s\,(1-s)^3
E^{(4)}(p_0+s\delta p).
\]
对一标准差工作窗 $|\delta p|\le\sigma_p$ 定义
\[
\epsilon_3(t)
\equiv\frac{|E'''_0|t\sigma_p^3}{6\hbar},
\qquad
\epsilon_4(t)
\equiv
\frac{t\sigma_p^4}{24\hbar}
\sup_{|p-p_0|\le\sigma_p}|E^{(4)}(p)|.
\]
于是精确色散截断量满足
\begin{equation}
\boxed{
\epsilon_{\rm disp}^{(\ge3)}(t;\sigma_p)
\equiv
\sup_{|\delta p|\le\sigma_p}
\left|\frac{t}{\hbar}[E(p_0+\delta p)-E_2(p_0+\delta p)]\right|
\le\epsilon_3(t)+\epsilon_4(t).}
\label{eq:electron_dispersion_cubic_control}
\end{equation}
因此 $\epsilon_3$ 只有在同时满足 $\epsilon_4\ll\epsilon_3$ 时才可单独作为“三阶主导”的简化指标；二阶传播本身的受控条件是右端总上界远小于 $1$。一旦这一条件失效，应回到精确传播相位 $\exp[-\ii E(p)t/\hbar]$，而不是引入经验“修正质量”。

对 PINEM/QPINEM 产生的离散相干能量梳，同一剪切可在能量--时间傅里叶对偶中读取。设边带能量为 $E_0+\ell\hbar\omega$，并把随中心电子运动的延迟写成
\[
\tau\equiv t-\frac{z}{v_0}.
\]
若共同慢包络为 $F(\tau)$，边带振幅为 $c_\ell$，则
\begin{equation}
\boxed{
\Psi(\tau)=F(\tau)\sum_\ell c_\ell\ee^{-\ii\ell\omega\tau},
\qquad
I(\tau)=|F(\tau)|^2
\sum_{\ell,m}c_\ell c_m^*\ee^{-\ii(\ell-m)\omega\tau}.}
\label{eq:coherent_energy_comb_time_amplitude}
\end{equation}
只有非对角相干 $c_\ell c_m^*$ 才产生光学周期或亚周期密度结构；边带概率很宽却彼此无相位关系时，不会自动形成阿秒脉冲列。若相干支持大致延伸到 $|\ell|\le L$，其能量带宽 $\mathcal B_E^{\rm coh}\sim2L\hbar\omega$ 给出傅里叶尺度
\[
\Delta t\gtrsim\frac{\hbar}{\mathcal B_E^{\rm coh}},
\]
但真正达到这一尺度还要求传播后的二次及更高阶谱相位得到控制。

在固定传播距离 $L_z$ 上，最方便的是展开纵向动量 $p(E)$。由 $\dd p/\dd E=1/v$ 和 $\dd v/\dd p=1/m_\parallel$，中心能量处
\[
\left.\frac{\dd^2 p}{\dd E^2}\right|_{E_0}
=-\frac{1}{m_\parallel v_0^3}.
\]
因此第 $\ell$ 个边带在去掉共同相位与线性时间平移后获得二次谱相位
\begin{equation}
\boxed{
\varphi_\ell^{(2)}(L_z)
=-\ell^2\vartheta_2(L_z),
\qquad
\vartheta_2(L_z)
\equiv\frac{\hbar\omega^2L_z}{2m_\parallel v_0^3}.}
\label{eq:energy_comb_quadratic_phase}
\end{equation}
由 $\dd p/\dd E=1/v$ 与 $\dd v/\dd p=1/m_\parallel$ 逐步展开，二阶 Taylor 系数因此完全由中心群速度和纵向色散质量确定。
因此 $\vartheta_2(L_z)$ 直接衡量离散能量格的二次传播相位，并决定时间聚焦与 Talbot 相位匹配条件。

\begin{figure}[htbp]
\centering
\includegraphics[width=0.76\linewidth]{gaussian_compression.png}
\caption{带负 $z$--$p$ 协方差的高斯电子包在二阶相对论色散下先压缩后展宽。焦点由\eqnrefs{eq:electron_chirp_compression_bound} 决定。}
\label{fig:gaussian-compression}
\end{figure}

Wigner 边缘性质属于定义层面的 Fourier 代数，不另设推导框。对\eqnrefs{eq:electron_wigner} 积分动量并用

\[
\int\dd p\,\ee^{-\ii ps/\hbar}=2\pi\hbar\,\delta(s),
\]
得到
\[
\int\dd p\,W(z,p)
=\rho(z,z).
\]
反过来积分 $z$，作变量变换 $z_1=z+s/2,z_2=z-s/2$，雅可比行列式为一，并使用全文约定
\[
\langle z|p\rangle=(2\pi\hbar)^{-1/2}\ee^{\ii pz/\hbar},
\]
于是
\begin{align*}
\int\dd z\,W(z,p)
&=\frac{1}{2\pi\hbar}
\int\dd z_1\dd z_2\,
\ee^{-\ii p(z_1-z_2)/\hbar}\rho(z_1,z_2)\\
&=\langle p|\hat\rho|p\rangle.
\end{align*}
这同时固定了\eqnrefs{eq:electron_wigner} 前因子 $1/(2\pi\hbar)$。

$\epsilon_3$ 衡量二次模型遗漏的首个谱相位，$\epsilon_4$ 控制继续截断在三阶本身的误差；一般应检查两者之和，而不能只看三阶项。

\begin{derivation}[从\eqnrefs{eq:weyl_star_product_derived}到\eqnrefs{eq:moyal_general}]

先从 Weyl 符号的 Fourier 表示出发。对足够良性的相空间函数 \(A_W(z,p)\)，写成
\[
A_W(z,p)
=
\int\frac{\dd\xi_z\,\dd\xi_p}{(2\pi)^2}\,
\widetilde A(\xi_z,\xi_p)
\ee^{\ii(\xi_z z+\xi_p p)},
\]
相应 Weyl 量子化为
\[
\hat A
=
\int\frac{\dd\xi_z\,\dd\xi_p}{(2\pi)^2}\,
\widetilde A(\xi_z,\xi_p)
\ee^{\ii(\xi_z\hat z+\xi_p\hat p)}.
\]
同理
\[
\hat B
=
\int\frac{\dd\xi_z'\,\dd\xi_p'}{(2\pi)^2}\,
\widetilde B(\xi_z',\xi_p')
\ee^{\ii(\xi_z'\hat z+\xi_p'\hat p)}.
\]

令
\[
X=\ii(\xi_z\hat z+\xi_p\hat p),
\qquad
Y=\ii(\xi_z'\hat z+\xi_p'\hat p).
\]
利用
\[
[\hat z,\hat p]=\ii\hbar
\]
得到
\[
[X,Y]
=
-\ii\hbar
(\xi_z\xi_p'-\xi_p\xi_z')\,I.
\]
该对易子是中心元，所以
\[
[X,[X,Y]]
=
[Y,[X,Y]]
=0,
\]
Baker--Campbell--Hausdorff 级数精确截断为
\[
\ee^X\ee^Y
=
\ee^{X+Y+\frac12[X,Y]}.
\]
于是
\begin{align*}
&\ee^{\ii(\xi_z\hat z+\xi_p\hat p)}
\ee^{\ii(\xi_z'\hat z+\xi_p'\hat p)}\\
&\qquad=
\exp\!\left[
-\frac{\ii\hbar}{2}
(\xi_z\xi_p'-\xi_p\xi_z')
\right]
\ee^{\ii[(\xi_z+\xi_z')\hat z+(\xi_p+\xi_p')\hat p]}.
\end{align*}

把这一乘法代回 \(\hat A\hat B\)，再取 Weyl 符号，得到
\begin{align*}
(\hat A\hat B)_W(z,p)
&=
\int
\frac{\dd\xi_z\,\dd\xi_p\,\dd\xi_z'\,\dd\xi_p'}{(2\pi)^4}\,
\widetilde A(\xi_z,\xi_p)
\widetilde B(\xi_z',\xi_p')\\
&\quad\times
\ee^{-\frac{\ii\hbar}{2}(\xi_z\xi_p'-\xi_p\xi_z')}
\ee^{\ii[(\xi_z+\xi_z')z+(\xi_p+\xi_p')p]}.
\end{align*}
另一方面，对 Fourier 基函数有
\[
\partial_z
\ee^{\ii(\xi_z z+\xi_p p)}
=
\ii\xi_z
\ee^{\ii(\xi_z z+\xi_p p)},
\qquad
\partial_p
\ee^{\ii(\xi_z z+\xi_p p)}
=
\ii\xi_p
\ee^{\ii(\xi_z z+\xi_p p)}.
\]
因此相位因子可以等价写成双向导数算符：
\[
\ee^{-\frac{\ii\hbar}{2}(\xi_z\xi_p'-\xi_p\xi_z')}
\longleftrightarrow
\exp\!\left[
\frac{\ii\hbar}{2}
\left(
\overleftarrow\partial_z\overrightarrow\partial_p
-
\overleftarrow\partial_p\overrightarrow\partial_z
\right)
\right].
\]
于是
\[
\boxed{
(\hat A\hat B)_W
=
A_W
\exp\!\left[
\frac{\ii\hbar}{2}
\left(
\overleftarrow\partial_z\overrightarrow\partial_p
-
\overleftarrow\partial_p\overrightarrow\partial_z
\right)
\right]
B_W
\equiv A_W\star B_W,
}
\]
即正文\eqnrefs{eq:weyl_star_product_derived}。

现在对 Liouville--von Neumann 方程
\[
\ii\hbar\partial_t\hat\rho
=
[\hat H,\hat\rho]
\]
取 Weyl 符号：
\[
\ii\hbar\partial_tW
=
H_W\star W-W\star H_W.
\]
两边除以 \(\ii\hbar\)，得到
\[
\boxed{
\partial_tW
=
\{H_W,W\}_{\rm M}
\equiv
\frac1{\ii\hbar}
(H_W\star W-W\star H_W),
}
\]
即正文\eqnrefs{eq:moyal_general}。

若把星积按 \(\hbar\) 展开，偶数阶在交换子中相消，只保留奇数阶：
\begin{align*}
\{H_W,W\}_{\rm M}
&=
\left(
\partial_zH_W\,\partial_pW
-\partial_pH_W\,\partial_zW
\right)\\
&\quad
-\frac{\hbar^2}{24}
\left(
\partial_z^3H_W\,\partial_p^3W
-3\partial_z^2\partial_pH_W\,
  \partial_p^2\partial_zW
+3\partial_z\partial_p^2H_W\,
  \partial_p\partial_z^2W
-\partial_p^3H_W\,\partial_z^3W
\right)
+\mathcal O(\hbar^4).
\end{align*}
最低阶正是 Poisson 括号，
\[
\{H_W,W\}_{\rm M}
=
\{H_W,W\}_{\rm PB}
+\mathcal O(\hbar^2).
\]
\end{derivation}
经典 Liouville 方程是 Moyal 演化的最低阶梯度极限；对至多二次的 Hamiltonian，三阶及更高导数严格为零，因此量子 Wigner 演化与经典相空间平流精确重合。

\begin{derivation}[从\eqnrefs{eq:relativistic_free_moyal}到\eqnrefs{eq:wigner_shear}]
自由电子 Hamiltonian 只依赖动量，$H_W=E(p)$。从 Moyal 方程的星积定义出发，位置导数只能作用在 $W$ 上，因此
\begin{align*}
 \partial_tW
 &=\frac{1}{\ii\hbar}(E\star W-W\star E)\\
 &=-\sum_{n=0}^{\infty}
 \frac{(-1)^n}{(2n+1)!}
 \left(\frac{\hbar}{2}\right)^{2n}
 E^{(2n+1)}(p)\,\partial_z^{2n+1}W,
\end{align*}
这就是正文\eqnrefs{eq:relativistic_free_moyal}。现在在中心动量 $p_0$ 附近定义二次模型
\begin{equation*}
 E_2(p)=E_0+v_0(p-p_0)+\frac{(p-p_0)^2}{2m_\parallel}.
\end{equation*}
因为 $E_2^{(r)}(p)=0$ 对所有 $r\ge3$ 成立，Moyal 级数不是“近似丢掉量子项”，而是在这个二次 Hamiltonian 内严格终止：
\begin{equation*}
 \partial_tW+
 \left(v_0+\frac{p-p_0}{m_\parallel}\right)\partial_zW=0.
\end{equation*}
这是一个一阶线性偏微分方程。令参数 $s$ 沿特征线变化，要求
\begin{equation*}
 \frac{\dd}{\dd s}W[z(s),p(s),t(s)]
 =\partial_tW\,\dot t+
 \partial_zW\,\dot z+
 \partial_pW\,\dot p=0.
\end{equation*}
与偏微分方程逐项比较，可选
\begin{equation*}
 \dot t=1,
 \qquad
 \dot p=0,
 \qquad
 \dot z=v_0+\frac{p-p_0}{m_\parallel}.
\end{equation*}
于是 $p(s)=p(0)$ 保持不变，而
\begin{equation*}
 z(t)=z(0)+
 \left(v_0+\frac{p-p_0}{m_\parallel}\right)t.
\end{equation*}
沿同一条特征线 $W$ 为常数，因此把初始位置反解为
$z(0)=z-[v_0+(p-p_0)/m_\parallel]t$，得到
\begin{equation*}
 W(z,p,t)=W_0\!\left(
 z-\left[v_0+\frac{p-p_0}{m_\parallel}\right]t,
 p\right),
\end{equation*}
即正文\eqnrefs{eq:wigner_shear}。

“剪切”一词现在有精确几何含义：同一 $p$ 水平线上的分布只做刚性平移，但平移距离依赖 $p$，所以原本竖直的相空间线元会倾斜。若使用完整相对论 $E(p)$，三阶及更高导数重新出现，演化不再是这个线性辛剪切；其误差必须由正文给出的高阶累计相位条件控制。
\end{derivation}

\begin{derivation}[从\eqnrefs{eq:electron_shear_variance}到\eqnrefs{eq:electron_chirp_compression_bound}]
设初态二阶矩有限，不要求波包为 Gaussian。共动二次 Hamiltonian 为
\[
H_{\rm com}=\frac{\delta p^2}{2m_\parallel},
\qquad
[\zeta,\delta p]=\ii\hbar.
\]
先在 Heisenberg 绘景计算算符运动。因为
\begin{align*}
[\zeta,\delta p^2]
&=[\zeta,\delta p]\,\delta p
+\delta p\,[\zeta,\delta p]\\
&=2\ii\hbar\,\delta p,
\end{align*}
所以
\begin{align*}
\dot\zeta
&=\frac{1}{\ii\hbar}[\zeta,H_{\rm com}]
=\frac{\delta p}{m_\parallel},\\
\dot{\delta p}
&=\frac{1}{\ii\hbar}[\delta p,H_{\rm com}]=0.
\end{align*}
积分得到
\[
\zeta(t)=\zeta(0)+\frac{t}{m_\parallel}\delta p(0),
\qquad
\delta p(t)=\delta p(0).
\]
由于 $[\zeta(0),\delta p(0)]=\ii\hbar$，平方时必须保留交叉项的算符次序：
\begin{align*}
\zeta(t)^2
&=\zeta^2(0)
+\frac{t}{m_\parallel}
\bigl[\zeta(0)\delta p(0)+\delta p(0)\zeta(0)\bigr]
+\frac{t^2}{m_\parallel^2}\delta p^2(0).
\end{align*}
取期望值，并使用正文定义
\[
C_{zp}(0)=\frac12
\langle\zeta(0)\delta p(0)+\delta p(0)\zeta(0)\rangle,
\]
便得到
\[
\sigma_z^2(t)
=\sigma_z^2(0)
+\frac{2t}{m_\parallel}C_{zp}(0)
+\frac{t^2}{m_\parallel^2}\sigma_p^2,
\]
即正文\eqnrefs{eq:electron_shear_variance}。把右端视为关于 $t$ 的二次函数，求导得
\[
\frac{\dd\sigma_z^2(t)}{\dd t}
=\frac{2}{m_\parallel}C_{zp}(0)
+\frac{2t}{m_\parallel^2}\sigma_p^2.
\]
若 $\sigma_p^2>0$，驻点唯一且二阶导数
\[
\frac{\dd^2\sigma_z^2(t)}{\dd t^2}
=\frac{2\sigma_p^2}{m_\parallel^2}>0,
\]
因此它确为最小值。解得
\[
t_*=-m_\parallel\frac{C_{zp}(0)}{\sigma_p^2}.
\]
代回二次函数：
\begin{align*}
\sigma_{z,\min}^2
&=\sigma_z^2(0)
-2\frac{C_{zp}^2(0)}{\sigma_p^2}
+\frac{C_{zp}^2(0)}{\sigma_p^2}\\
&=\sigma_z^2(0)-\frac{C_{zp}^2(0)}{\sigma_p^2}.
\end{align*}
最后对正则共轭对 $(\zeta,\delta p)$ 使用 Robertson--Schr\"odinger 不等式
\[
\sigma_z^2(0)\sigma_p^2-C_{zp}^2(0)
\ge\frac14|\langle[\zeta,\delta p]\rangle|^2
=\frac{\hbar^2}{4}.
\]
由于 $\sigma_p^2>0$，两边除以 $\sigma_p^2$ 得
\[
\sigma_z^2(0)-\frac{C_{zp}^2(0)}{\sigma_p^2}
\ge\frac{\hbar^2}{4\sigma_p^2}.
\]
与上式合并即得到正文\eqnrefs{eq:electron_chirp_compression_bound}。
\end{derivation}

\begin{derivation}[从\eqnrefs{eq:dispersion_truncation_control}到\eqnrefs{eq:electron_dispersion_cubic_control}]
取局部工作动量窗 $\mathcal W_{p,1}=\{p:|p-p_0|\le\sigma_p\}$，令 $\delta p=p-p_0$。自由电子的精确正能色散为
\[
E(p)=\sqrt{\me^2c^4+c^2p^2}.
\]
为了把余项界写成可直接计算的量，先从这个一般表达式求高阶导数。一次求导给
\[
E'(p)=\frac{c^2p}{E(p)}.
\]
再次求导时保留商的两项：
\begin{align*}
E''(p)
&=\frac{c^2}{E(p)}
-\frac{c^2p}{E^2(p)}E'(p)\\
&=\frac{c^2}{E(p)}-\frac{c^4p^2}{E^3(p)}\\
&=\frac{\me^2c^6}{E^3(p)}.
\end{align*}
因此
\begin{align*}
E'''(p)
&=-3\me^2c^6E^{-4}(p)E'(p)\\
&=-\frac{3\me^2c^8p}{E^5(p)},
\end{align*}
在 $p=p_0$ 处恢复正文使用的
\[
E'''_0=-\frac{3v_0}{\gamma_0^4\me^2c^2}.
\]
继续求一次导数可得四阶项
\begin{align*}
E^{(4)}(p)
&=-3\me^2c^8
\left[E^{-5}(p)-5pE^{-6}(p)E'(p)\right]\\
&=-\frac{3\me^2c^8}{E^5(p)}
+\frac{15\me^2c^{10}p^2}{E^7(p)}\\
&=\frac{3\me^2c^8}{E^7(p)}
\left(4c^2p^2-\me^2c^4\right).
\end{align*}
这些导数决定二次传播模型被舍弃的谱相位尺度。

对 $E(p_0+\delta p)$ 在 $p_0$ 处使用三阶 Taylor 公式的积分余项，
\begin{align*}
E(p_0+\delta p)-E_2(p_0+\delta p)
&=\frac{E'''_0}{6}\delta p^3+R_4(\delta p),\\
R_4(\delta p)
&=\frac{\delta p^4}{6}
\int_0^1\dd s\,(1-s)^3
E^{(4)}(p_0+s\delta p).
\end{align*}
积分余项保留了从 $p_0$ 到 $p_0+\delta p$ 整条线段上的四阶导数。因为 $|\delta p|\le\sigma_p$ 且 $0\le s\le1$，积分中的动量始终满足 $|p_0+s\delta p-p_0|\le\sigma_p$。于是
\begin{align*}
|R_4(\delta p)|
&\le\frac{|\delta p|^4}{6}
\int_0^1\dd s\,(1-s)^3
\left|E^{(4)}(p_0+s\delta p)\right|\\
&\le\frac{|\delta p|^4}{6}
\sup_{|p-p_0|\le\sigma_p}|E^{(4)}(p)|
\int_0^1\dd s\,(1-s)^3\\
&=\frac{|\delta p|^4}{24}
\sup_{|p-p_0|\le\sigma_p}|E^{(4)}(p)|\\
&\le\frac{\sigma_p^4}{24}
\sup_{|p-p_0|\le\sigma_p}|E^{(4)}(p)|.
\end{align*}
把这个结果代入正文\eqnrefs{eq:dispersion_truncation_control} 的定义，得到
\begin{align*}
\epsilon_{\rm disp}^{(\ge3)}(t;\sigma_p)
&=\frac{t}{\hbar}
\sup_{|\delta p|\le\sigma_p}
\left|E(p_0+\delta p)-E_2(p_0+\delta p)\right|\\
&\le
\frac{t}{\hbar}
\sup_{|\delta p|\le\sigma_p}
\left(\frac{|E'''_0|}{6}|\delta p|^3+|R_4(\delta p)|\right)\\
&\le
\frac{|E'''_0|t\sigma_p^3}{6\hbar}
+\frac{t\sigma_p^4}{24\hbar}
\sup_{|p-p_0|\le\sigma_p}|E^{(4)}(p)|\\
&=\epsilon_3(t)+\epsilon_4(t),
\end{align*}
即正文\eqnrefs{eq:electron_dispersion_cubic_control}。
\end{derivation}



\subsection{时间透镜与 Talbot 传播}

前面的相空间计算表明，若先让波包不同位置具有不同平均速度，自由传播就可能使它压缩。对纵向电子波包，随到达时间变化的光学相位可以产生这种速度相关，作用类似空间光学中的薄透镜，因此称为\emph{时间透镜}。另一方面，若波包含有等间隔的相干能量边带，各边带的二次传播相位还可能在特定距离重新对齐，使周期时间图样复现；这称为\emph{Talbot 复现}。下面分别由传播相位计算聚焦条件和复现条件。

对经典相干光场，PINEM 章节定义的无量纲耦合 $\beta=|\beta|\ee^{\ii\varphi_\beta}$ 决定电子获得的周期相位。用到达延迟 $\tau=t-z/v_0=-\zeta/v_0$ 表示，能量阶数 $\ell$ 的 Fourier 因子固定为 $\ee^{-\ii\ell\omega\tau}$，因此，单次相位调制写成
\begin{equation}
\boxed{
\Psi(\tau)\longrightarrow
\exp\!\left[2\ii|\beta|\sin(\varphi_\beta-\omega\tau)\right]\Psi(\tau).}
\label{eq:electron_periodic_time_phase}
\end{equation}
它不是全局抛物线，因此“时间透镜”只在选定光学相位附近成立。若以相位极值 $\varphi_\beta-\omega\tau_0=\pi/2$ 为中心，令 $\delta\tau=\tau-\tau_0$，则局域展开给
\begin{equation}
\boxed{
\varphi_{\rm lens}(\delta\tau)
=-|\beta|\omega^2\delta\tau^2
+O\!\left(|\beta|\omega^4\delta\tau^4\right).}
\label{eq:electron_local_time_lens_phase}
\end{equation}
常数相位已省略。控制“抛物线透镜”近似的参数是 $|\omega\delta\tau|\ll1$；超出这一时间窗，正弦相位的高阶项会生成周期性焦散/脉冲列，而不是单个理想高斯焦点。

局域二次时间相位先建立到达时间与能量的相关；随后自由色散把这种相关转化为实际的到达时间压缩。粗粒化地看，\eqnrefs{eq:classical_energy_mod_from_beta} 给相位 $\theta=\varphi_\beta-\omega\tau$ 上的能量调制
\[
\delta E(\theta)=2\hbar\omega|\beta|\cos\theta.
\]
经过距离 $L_z$ 后，速度色散把它映射成到达时间偏移。固定传播距离时
\begin{align*}
\tau(E)&=\frac{L_z}{v(E)},\\
\left.\frac{\dd(1/v)}{\dd E}\right|_{E_0}
&=-\frac1{v_0^2}
\left.\frac{\dd v/\dd p}{\dd E/\dd p}\right|_{p_0}
=-\frac{1}{m_\parallel v_0^3},
\end{align*}
其中用了 $\dd E/\dd p=v$ 与 $\dd v/\dd p=1/m_\parallel$。因此对小能量调制
\begin{equation}
\boxed{
\delta\tau_{\rm drift}(\theta)
=
-\frac{2L_z\hbar\omega|\beta|}{m_\parallel v_0^3}\cos\theta
+O\!\left([\delta E(\theta)]^2\right).}
\label{eq:electron_classical_bunching_map}
\end{equation}
这里的一阶式只要求在实际占据能宽内 $1/v(E)$ 的二阶及更高变化可忽略；若能量调制大到越过这一局域线性区，应直接使用精确飞行时间 $L_z/v(E_0+\delta E)$。当映射 $\tau_{\rm out}(\theta)$ 的雅可比行列式变小时，经典束流图像中出现密度聚束；量子图像中，同一现象由\eqnrefs{eq:coherent_energy_comb_time_amplitude} 的边带交叉项干涉形成。两种语言在边带不可分辨、包络跨越许多光学周期的粗粒化极限相接。

另一种特殊相位匹配不是局域聚焦，而是整个离散能量梳的谱相位复现。\secref{subsec:sideband-to-temporal-density-dp46}已经从完整密度谐波得到唯一的 Talbot 时间和长度~\eqnrefs{eq:electron_talbot_time_length_dp46}。这里用能量表象检验同一结果：由\eqnrefs{eq:energy_comb_quadratic_phase}，当 $\vartheta_2(L_T)=2\pi$ 时，任意整数 $\ell$ 都满足 $\exp[-\ii\ell^2\vartheta_2(L_T)]=1$；代入 $\vartheta_2=\hbar\omega^2L_z/(2m_\parallel v_0^3)$ 恰好恢复\eqnrefs{eq:electron_talbot_time_length_dp46} 在 $\kappa_z=\omega/v_0$ 下的 $L_T$。在 $L_z/L_T$ 为有理数时，$\ell^2$ 的模整数结构产生分数 Talbot 图样。初始能散、有限包络、空间电荷以及\eqnrefs{eq:electron_dispersion_cubic_control} 所衡量的高阶色散都会降低复现对比度。

\begin{figure}[H]
  \centering
  \includegraphics[width=0.82\linewidth]{talbot_time_focusing.png}
  \caption{周期相位调制后的自由电子在二次谱相位传播下形成的时域 Talbot 结构。图中从 $c_\ell=J_\ell(2|\beta|)$ 出发，对每个边带乘以 $\exp(-\ii s\ell^2)$ 后重新求和；取 $|\beta|=3$。\cite{Feist2015Nature}}
  \label{fig:talbot_time_focusing}
\end{figure}

经典纵向束流动力学使用 $(z,\delta E)$ 相空间描述啁啾与微聚束。\cite{Pellegrini2016FEL} 当单电子相干、反冲和负 Wigner 结构不可忽略时，Wigner 表示保留同一几何直觉并补上量子相干。因此二者不是互斥理论，而是不同分辨尺度上的连续描述。


精确相对论色散决定电子闭合传播，相空间表示把同一动力学改写成几何剪切，周期相位与二次谱相位则决定聚焦和复现。未追踪环境会把闭合幺正动力学变成约化量子信道，而探测器进一步把量子态映射为经典数据；二者分别由开放系统动力学与量子测量描述。

上述复现只涉及二次色散。若占据带宽内的三阶及更高谱相位在一个 $L_T$ 上积累到 $O(1)$，不同整数边带不能同时回相；这一偏离由正文\eqnrefs{eq:electron-talbot-higher-order-control-dp76} 的 $\epsilon_T^{(\ge3)}$ 控制。

\begin{derivation}[从\eqnrefs{eq:electron_periodic_time_phase}到\eqnrefs{eq:electron_local_time_lens_phase}]
在相位极值附近对正弦调制做 Taylor 展开，推出\eqnrefs{eq:electron_local_time_lens_phase} 并识别四阶项作为有限孔径误差。 透镜近似要求 $|\omega\delta\tau|\ll1$； 超出这一局域时间窗时应使用完整的正弦相位调制。 定义光学相位

\begin{equation*}
\theta(\tau)=\varphi_\beta-\omega\tau .
\end{equation*}
电子的周期相位调制为 $2|\beta|\sin\theta$。选择相位极值
\begin{equation*}
\theta_0=\frac{\pi}{2},
\qquad
\varphi_\beta-\omega\tau_0=\frac{\pi}{2},
\end{equation*}
并写
\begin{equation*}
\delta\tau=\tau-\tau_0,
\qquad
\delta\theta=-\omega\delta\tau.
\end{equation*}
利用
\begin{equation*}
\sin\!\left(\frac\pi2+\delta\theta\right)
=\cos\delta\theta
=1-\frac{\delta\theta^2}{2}
+\frac{\delta\theta^4}{24}
+O(\delta\theta^6),
\end{equation*}
相位指数的实函数变成
\begin{align*}
2|\beta|\sin\theta
&=2|\beta|
-|\beta|\delta\theta^2
+\frac{|\beta|}{12}\delta\theta^4
+O(|\beta|\delta\theta^6)\\
&=2|\beta|
-|\beta|\omega^2\delta\tau^2
+\frac{|\beta|\omega^4\delta\tau^4}{12}
+O(|\beta|\omega^6\delta\tau^6).
\end{align*}
第一项只产生与 $\delta\tau$ 无关的整体相位，去掉它以后得到正文的局域二次相位
\begin{equation*}
\varphi_{\rm lens}(\delta\tau)
=-|\beta|\omega^2\delta\tau^2
+O(|\beta|\omega^4\delta\tau^4).
\end{equation*}
四次项相对于二次项的绝对比值为
\begin{equation*}
\left|
\frac{(|\beta|/12)\omega^4\delta\tau^4}
{|\beta|\omega^2\delta\tau^2}
\right|
=\frac{(\omega\delta\tau)^2}{12}.
\end{equation*}
\end{derivation}
抛物线时间透镜的局域控制量是 $|\omega\delta\tau|$；它不再远小于一时，应使用完整周期正弦相位。
\begin{derivation}[从\eqnrefs{eq:energy_comb_quadratic_phase}到\eqnrefs{eq:electron_talbot_time_length_dp46}]
设占据的能量边带由整数 $\ell\in\mathbb Z$ 标记，并先只保留二次色散。第 $\ell$ 个边带在传播长度 $L_z$ 后获得
\[
U_\ell^{(2)}(L_z)
=\exp[-\ii\ell^2\vartheta_2(L_z)],
\qquad
\vartheta_2(L_z)=\frac{\hbar\omega^2L_z}{2m_\parallel v_0^3}.
\]
若传播前纯态写成
\[
|\psi(0)\rangle=\sum_\ell c_\ell|\ell\rangle,
\]
传播后为
\[
|\psi(L_z)\rangle
=\sum_\ell c_\ell
\ee^{-\ii\ell^2\vartheta_2(L_z)}|\ell\rangle.
\]
密度矩阵元因此满足
\begin{align*}
\rho_{\ell m}(L_z)
&=c_\ell c_m^*
\ee^{-\ii\ell^2\vartheta_2}
\ee^{+\ii m^2\vartheta_2}\\
&=\rho_{\ell m}(0)
\ee^{-\ii(\ell^2-m^2)\vartheta_2}.
\end{align*}
对混态同样成立，因为自由传播在线性算符层面逐个矩阵元乘上同一相对相位。

完全原位复现要求所有已占据矩阵元同时恢复。若希望得到与具体占据集合无关的整数格复现条件，就要求
\[
\ee^{-\ii(\ell^2-m^2)\vartheta_2(L_T)}=1
\qquad
\forall\,\ell,m\in\mathbb Z.
\]
取 $m=0$，必要条件为
\[
\ell^2\vartheta_2(L_T)=2\pi n_\ell,
\qquad n_\ell\in\mathbb Z.
\]
其中 $\ell=1$ 已经要求
\[
\vartheta_2(L_T)=2\pi n_1.
\]
最小正复现取 $n_1=1$，即
\[
\vartheta_2(L_T)=2\pi.
\]
反过来，把这个值代回任意整数 $\ell,m$：
\begin{align*}
(\ell^2-m^2)\vartheta_2(L_T)
&=2\pi(\ell^2-m^2),
\end{align*}
右端总是 $2\pi$ 的整数倍，所以该条件不仅必要，而且对完整整数能量格充分。于是
\begin{align*}
2\pi
&=\frac{\hbar\omega^2L_T}{2m_\parallel v_0^3},\\
L_T
&=\frac{4\pi m_\parallel v_0^3}{\hbar\omega^2},\\
T_T=\frac{L_T}{v_0}
&=\frac{4\pi m_\parallel v_0^2}{\hbar\omega^2},
\end{align*}
得到正文\eqnrefs{eq:electron_talbot_time_length_dp46}。

半个 Talbot 长度对应
\[
\vartheta_2(L_T/2)=\pi.
\]
此时单边带相位为
\begin{align*}
U_\ell^{(2)}(L_T/2)
&=\ee^{-\ii\pi\ell^2}\\
&=(-1)^{\ell^2}\\
&=(-1)^\ell,
\end{align*}
最后一步使用整数 $\ell$ 与 $\ell^2$ 具有相同奇偶性。另一方面，共动坐标平移
\[
\Delta\zeta=\frac{\pi}{\kappa_z}
\]
会让第 $\ell$ 个空间 Fourier 分量获得
\[
\ee^{\ii\ell\kappa_z\Delta\zeta}
=\ee^{\ii\pi\ell}=(-1)^\ell.
\]
因此 $L_T/2$ 的传播与半周期平移具有相同的格点相位作用，而 $L_T$ 才对应原位完全复现。
\end{derivation}

\subsection{有限高斯波包：解析解 TDSE 与时间聚焦}
\label{sec:electron-gaussian-lens-analytic}

时间透镜的聚焦也可以直接从一个有限波包的波函数计算。取已建立的纵向二阶色散，在随中心电子运动的坐标 $\zeta=z-v_0t$ 中，包络满足含时 Schrödinger 方程（time-dependent Schrödinger equation, TDSE）
\begin{equation}
 \ii\hbar\partial_t\psi(\zeta,t)
 =-\frac{\hbar^2}{2m_\parallel}\partial_\zeta^2\psi(\zeta,t),
 \qquad \int|\psi|^2\dd\zeta=1.
 \label{eq:electron-gaussian-envelope-tdse}
\end{equation}
这个方程与自由非相对论粒子的 Schrödinger 方程形式相同，但系数 $m_\parallel=\gamma_0^3m_{\mathrm e}$ 来自相对论能量的二次展开。中心能量和匀速运动已被移入载波相位，$\psi$ 只描述相对于中心电子的包络变化。它的适用条件仍是实际动量范围内高阶色散积累的相位可以忽略。

用复参数 $a_0$ 表示一般纯高斯输入，$\Re a_0>0$：
\begin{equation*}
 \psi(\zeta,0)=\mathcal N_0\ee^{-a_0\zeta^2},\qquad
 \mathcal N_0=\left(\frac{2\Re a_0}{\pi}\right)^{1/4}.
\end{equation*}
$a_0$ 的实部决定密度宽度，虚部决定二次空间相位；$[a_0]=\mathrm{m^{-2}}$。连续 TDSE 的精确高斯解为
\begin{equation}
 \psi(\zeta,t)=\frac{\mathcal N_0}{\sqrt{d(t)}}
 \exp\!\left[-\frac{a_0\zeta^2}{d(t)}\right],\qquad
 d(t)=1+\frac{2\ii\hbar a_0t}{m_\parallel}.
 \label{eq:electron-gaussian-tdse-solution}
\end{equation}
平方根支路从 $d(0)=1$ 连续选取。$d$ 无量纲；前因子同时包含保持归一化所需的幅度变化与共同传播相位，不能只保留指数中的宽度变化。

\begin{exbox}[带二次相位的波包何时最窄]
令初始密度标准差为 $\sigma_{\zeta0}$，空间啁啾参数为实数 $C_\zeta$，即
\begin{equation*}
 a_0=\frac{1-\ii C_\zeta}{4\sigma_{\zeta0}^2},\qquad
 t_{\rm dif}=\frac{2m_\parallel\sigma_{\zeta0}^2}{\hbar},\qquad
 u_t=\frac{t}{t_{\rm dif}}.
\end{equation*}
这里 $C_\zeta$ 是无量纲实数，控制位置表象中的二次相位；前面用 $C_p$ 表示的是动量表象中的二次相位。$t_{\rm dif}$ 具有时间单位，它是无啁啾高斯包络因自由传播而明显展宽的时间尺度，$u_t$ 则以这个时间为单位。代入一般解，有 $d=1+(C_\zeta+\ii)u_t$。利用
\begin{equation*}
 \Re\frac{a_0}{d}=\frac{\Re a_0}{|d|^2},\qquad
 |d|^2=(1+C_\zeta u_t)^2+u_t^2,
\end{equation*}
便从概率密度的高斯指数读出
\begin{equation*}
 \sigma_\zeta^2(t)=\sigma_{\zeta0}^2
 \left[1+2C_\zeta u_t+(1+C_\zeta^2)u_t^2\right].
\end{equation*}
这是关于时间的二次式，配方即可求极小值：
\begin{equation*}
 \frac{\sigma_\zeta^2(t)}{\sigma_{\zeta0}^2}
 =(1+C_\zeta^2)
 \left(u_t+\frac{C_\zeta}{1+C_\zeta^2}\right)^2
 +\frac1{1+C_\zeta^2}.
\end{equation*}
若 $C_\zeta<0$，未来存在焦点，
\begin{equation}
 t_f=-\frac{C_\zeta}{1+C_\zeta^2}t_{\rm dif},\qquad
 L_f=v_0t_f,\qquad
 \sigma_{\zeta,\min}=\frac{\sigma_{\zeta0}}{\sqrt{1+C_\zeta^2}}.
 \label{eq:electron-gaussian-focus}
\end{equation}
若 $C_\zeta=0$，波包从初始时刻起展宽；若 $C_\zeta>0$，二次式的焦点位于过去，未来同样单调展宽。因此聚焦需要正确符号的相位相关，而不只是足够大的能宽。

若输入本来无啁啾，在时间透镜的相位极值附近乘入 $\exp[-\ii|\beta|\omega^2\tau^2]$，并使用 $\zeta=-v_0\tau$，便得到
\begin{equation*}
 C_\zeta=-4|\beta|\omega^2\sigma_{\tau0}^2,
 \qquad \sigma_{\tau0}=\sigma_{\zeta0}/v_0.
\end{equation*}
其符号恰好允许未来聚焦；将它代入\eqnrefs{eq:electron-gaussian-focus} 就得到对任意允许参数成立的焦距与压缩比。到达时间宽度取 $\sigma_\tau\simeq\sigma_\zeta/v_0$，还要求速度展宽相对中心速度很小。
\end{exbox}

为了理解压缩过程中动量分布怎样变化，再计算位置与动量的协方差。记 $\Delta\zeta=\hat\zeta-\langle\hat\zeta\rangle$、$\Delta p=\hat p-\langle\hat p\rangle$；由于两个算符不对易，用对称乘积定义实协方差。对上述初态，在位置表象使用 $\hat p=-\ii\hbar\partial_\zeta$，得
\begin{equation*}
 \sigma_p^2=\frac{\hbar^2(1+C_\zeta^2)}{4\sigma_{\zeta0}^2},\qquad
 C_{\zeta p}\equiv\tfrac12\langle\Delta\zeta\Delta p+\Delta p\Delta\zeta\rangle
 =\frac{\hbar C_\zeta}{2}.
\end{equation*}
所以 $\sigma_{\zeta0}^2\sigma_p^2-C_{\zeta p}^2=\hbar^2/4$。自由传播保持 $\sigma_p$ 不变，只把协方差变为 $C_{\zeta p}(t)=C_{\zeta p}(0)+t\sigma_p^2/m_\parallel$；焦点恰在该协方差过零时，且 $\sigma_{\zeta,\min}\sigma_p=\hbar/2$。因此焦点处位置与动量不再相关，宽度乘积恰好达到最小不确定值。

时间透镜的适用条件也应作用于\emph{相位本身}。若检查包络的 $|\tau|\le K\sigma_{\tau0}$ 区域，其中 $K$ 是选定的覆盖倍数，则正弦相位余项的首项为
\begin{equation*}
 |\delta\varphi_4|\lesssim\frac{|\beta|}{12}
 (\omega K\sigma_{\tau0})^4.
\end{equation*}
要可靠地用二次相位替代完整作用，应同时检查 $\omega K\sigma_{\tau0}\ll1$ 和 $|\delta\varphi_4|\ll1$，并检查高阶色散在 $t_f$ 内的累积误差。仅有局部几何孔径很小，在极强驱动下仍不足以保证波函数误差很小。

\begin{derivation}[从自由传播核到高斯解，并直接代回 TDSE]
在二阶包络模型中，连续动量谱给传播核
\begin{equation*}
 K(\zeta,\zeta';t)=\int\frac{\dd p}{2\pi\hbar}
 \exp\!\left[\frac{\ii p(\zeta-\zeta')}{\hbar}
             -\frac{\ii p^2t}{2m_\parallel\hbar}\right]
 =\sqrt{\frac{m_\parallel}{2\pi\ii\hbar t}}
 \exp\!\left[\frac{\ii m_\parallel(\zeta-\zeta')^2}{2\hbar t}\right],
 \quad t>0.
\end{equation*}
振荡高斯积分可先加正的收敛因子，再取其趋零极限。将初态代入 $\psi(\zeta,t)=\int K\psi(\zeta',0)\dd\zeta'$，整理为
\begin{equation*}
 -A\zeta'^2+B\zeta'+C,
 \quad A=a_0-\frac{\ii m_\parallel}{2\hbar t},\quad
 B=-\frac{\ii m_\parallel\zeta}{\hbar t},\quad
 C=\frac{\ii m_\parallel\zeta^2}{2\hbar t}.
\end{equation*}
$\Re A>0$ 保证积分收敛；使用 $\int\ee^{-Ax^2+Bx}\dd x=\sqrt{\pi/A}\ee^{B^2/(4A)}$，合并前因子与指数即得\eqnrefs{eq:electron-gaussian-tdse-solution}。

还可独立代回检验。写 $a(t)=a_0/d(t)$、$\mathcal N(t)=\mathcal N_0/\sqrt{d(t)}$，则
\begin{equation*}
 \dot a=-\frac{2\ii\hbar}{m_\parallel}a^2,\qquad
 \frac{\dot{\mathcal N}}{\mathcal N}=-\frac{\ii\hbar}{m_\parallel}a,
 \qquad
 \frac{\partial_\zeta^2\psi}{\psi}=-2a+4a^2\zeta^2.
\end{equation*}
TDSE 左右两侧的常数项与 $\zeta^2$ 项于是分别相等。再用
$\Re(a_0/d)=\Re a_0/|d|^2$，对 $|\psi|^2$ 作高斯积分，归一化保持为一。所以传播后的函数同时满足初始条件、运动方程和归一化条件。
\end{derivation}


\section{开放系统动力学}

自由传播对应电子与其它自由度已经解耦的阶段。若辐射、材料涨落或未探测模式继续与电子交换能量或相位，就把实验仍要预测和控制的自由度称为\emph{系统} $S$，把不直接读出的其余自由度称为\emph{环境} $B$。系统与环境的总态仍按量子力学幺正演化；实验只读取 $S$ 时，对 $B$ 作偏迹得到约化态。环境在偏迹前保存的相关性可以在稍后反馈，因此约化动力学自然产生记忆。

只知道当前电子态一般不足以预测下一时刻，因为环境可能还保存先前相互作用留下的信息。我们先像消去联立方程中的未知量一样，消去环境相关，得到含历史积分的系统方程；随后再检查什么时候历史积分可以简化为只依赖当前态的演化。

\subsection{约化动力学}

设系统、环境与相互作用 Hamiltonian 分别为 $\hat H_S(t)$、$\hat H_B$ 与 $\hat H_{SB}(t)$，
\begin{equation}
\boxed{
\hat H(t)=\hat H_S(t)+\hat H_B+\hat H_{SB}(t),
\qquad
\frac{\dd}{\dd t}\hat\rho_{SB}(t)=\mathcal L(t)\hat\rho_{SB}(t).}
\label{eq:open_system_joint_liouville_equation}
\end{equation}
其中
\[
\mathcal L(t)X=-\frac{\ii}{\hbar}[\hat H(t),X]
\]
是联合态的 Liouville 生成元。这里 $\mathcal L$ 不是作用在 Hilbert 空间态矢上的普通算符，而是把一个算符 $X$ 映射成另一个算符的\emph{超算符}（superoperator）；Liouville--von Neumann 方程正是把 Schr\"odinger 绘景的幺正演化改写成密度算符空间中的一阶线性方程。以后出现的 $\mathcal P$、$\mathcal Q$、$\mathcal K$ 与 $\Phi$ 都属于同一类对象，因此它们之间的乘法表示超算符复合。

环境偏迹只保留
$\rho_S(t)=\operatorname{Tr}_B\rho_{SB}(t)$，即所有系统可观测量所需的约化态；但被迹掉的环境仍可能保存先前交换出去的振幅和相位，并在稍后重新作用于系统。约化动力学的核心问题因此是：怎样在不继续显式传播全部环境自由度的情况下，把这部分反馈严格编码进系统方程。

为把保留的系统信息与被消去的相关子空间分开，选定参考环境态 $\rho_B$，并定义投影超算符
\[
\mathcal PX=\operatorname{Tr}_B(X)\otimes\rho_B,
\qquad
\mathcal Q=1-\mathcal P.
\]
$\mathcal P\rho$ 保留“当前系统态 $\otimes$ 参考环境态”，而 $\mathcal Q\rho$ 包含相对于这一参考分解的全部系统--环境相关性。这里“投影”发生在\emph{算符空间}而不是电子的空间波函数上：$\mathcal P$ 选出我们希望显式保留的约化信息，$\mathcal Q$ 收集被暂时消去的相关自由度。先让 $\mathcal Q$ 中的信息从过去时刻 $s$ 演化到当前时刻 $t$，再把它反馈到 $\mathcal P$ 子空间，就得到 Nakajima--Zwanzig 方程
\begin{equation}
\boxed{
\frac{\dd}{\dd t}\mathcal P\rho(t)
=\mathcal P\mathcal L(t)\mathcal P\rho(t)
+\int_{t_0}^{t}\dd s\,\mathcal K(t,s)\mathcal P\rho(s)
+\mathcal I(t,t_0),}
\label{eq:exact_nakajima_zwanzig}
\end{equation}
其中记忆核与 $\mathcal Q$ 子空间传播子为
\[
\mathcal K(t,s)
=\mathcal P\mathcal L(t)\mathcal U_Q(t,s)\mathcal Q\mathcal L(s)\mathcal P,
\qquad
\mathcal U_Q(t,s)=\mathcal T\exp\!\left[\int_s^t\dd\tau\,\mathcal Q\mathcal L(\tau)\mathcal Q\right].
\]
这条乘积可以按时间顺序阅读：$\mathcal Q\mathcal L(s)\mathcal P$ 在时刻 $s$ 把系统信息送入被消去的相关子空间，$\mathcal U_Q(t,s)$ 保存并传播这部分信息，$\mathcal P\mathcal L(t)$ 再把它送回系统。因此 $\mathcal K(t,s)$ 就是“时刻 $s$ 的系统状态对时刻 $t$ 的变化率有多大影响”的算符权重，称为\emph{记忆核}。除了这种传播过程中形成的记忆，初态本身若已经含有系统--环境相关，也会产生独立的非齐次源项
\begin{equation}
\boxed{
\mathcal I(t,t_0)
=\mathcal P\mathcal L(t)\mathcal U_Q(t,t_0)\mathcal Q\rho(t_0).}
\label{eq:exact_nz_initial_correlation}
\end{equation}
因此非 Markov 性有两个不同来源：记忆核使当前导数依赖历史状态，非因子化初态则通过 $\mathcal I$ 产生额外边界源。只有同时控制记忆核与初始相关源项，时间局域主方程才是受控近似。


时间平移不变时，$\mathcal K(t,s)=\mathcal K(t-s)$，Laplace 变换把历史卷积转化为频域乘法。与普通量子力学中 $(z-H)^{-1}$ 的预解算符相同，这里在 Liouville 算符空间定义约化预解式
\begin{equation}
\boxed{
\mathcal R_P(z)
=\left[z-\mathcal P\mathcal L\mathcal P-\widetilde K(z)\right]^{-1},
\qquad
\widetilde K(z)
=\mathcal P\mathcal L\mathcal Q
(z-\mathcal Q\mathcal L\mathcal Q)^{-1}
\mathcal Q\mathcal L\mathcal P,}
\label{eq:nz_reduced_resolvent_definition}
\end{equation}
则完整联合预解式投影到 $\mathcal P$ 子空间以后严格满足
\begin{equation}
\boxed{
\mathcal P(z-\mathcal L)^{-1}\mathcal P
=\mathcal R_P(z).}
\label{eq:nz_projected_full_resolvent_schur_identity}
\end{equation}
这就是 Liouville 空间中的 Feshbach--Schur 消元。$\widetilde K(z)$ 与\chapref{chap:general-interaction}的 Hamiltonian 自能承担同一数学角色：它记录被消去子空间离开再返回以后对保留子空间产生的频率依赖修正，因此称为约化自能。$\mathcal R_P(z)$ 的孤立极点对应指数衰减或振荡模；连续谱造成的支割与其它非解析结构则产生长时间非指数尾。若系统实际采样的频带内 $\widetilde K(z)$ 的频率依赖足够弱，历史反馈可以被一个近似常数的局域生成元表示，这才进入 Markov 极限。

精确记忆方程同时保留系统--环境相关、环境记忆和不同系统跃迁频率之间的相干耦合。若 $H_S|n\rangle=E_n|n\rangle$，任意能级差定义一个\emph{Bohr 频率} $\omega_{mn}=(E_m-E_n)/\hbar$；它就是系统算符在相互作用绘景中出现的基本振荡频率。久期近似所比较的正是不同 Bohr 频率之差与约化动力学弛豫时间的倒数。

下面进行常用的弱耦合推导。先按耦合强度保留二阶项，这一步称为 Born 近似，所得方程仍含历史态 $\rho_S(t-\tau)$。若环境相关衰减得足够快，再在记忆积分中用当前态替代缓慢变化的历史态，作 Markov 近似。最后比较不同 Bohr 频率的拍频与弛豫速度，舍去长期快速振荡的交叉项，作\emph{久期近似}。在相应条件下，将得到 Gorini--Kossakowski--Sudarshan--Lindblad（GKSL）方程，通常简称 Lindblad 方程。为逐步看出这些近似，进入相互作用绘景，写
\[
H_I(t)=\lambda\sum_{\alpha}A_\alpha(t)\otimes B_\alpha(t),
\qquad
\operatorname{Tr}_B[B_\alpha\rho_B]=0,
\]
其中 $\lambda$ 只用于追踪耦合阶数。对精确 von Neumann 方程积分一次再代回自身，得到二次积分恒等式
\begin{equation}
\boxed{
\dot\rho_I(t)
=-\frac{\ii}{\hbar}[H_I(t),\rho_I(0)]
-\frac1{\hbar^2}\int_0^t\dd\tau\,
[H_I(t),[H_I(t-\tau),\rho_I(t-\tau)]].}
\label{eq:exact_second_order_integral_identity_open_system}
\end{equation}
Born 近似只在积分核内部把联合态近似为 $\rho_{S,I}(t-\tau)\otimes\rho_B$。对平稳环境定义相关矩阵
\[
C_{\alpha\beta}(\tau)
=\operatorname{Tr}_B[B_\alpha(\tau)B_\beta(0)\rho_B],
\]
得到的是\emph{二阶 Born 时间非局域方程}：耦合阶数已经截断，但历史态 $\rho_{S,I}(t-\tau)$ 和有限积分上限仍完整保留。它可以看成精确 Nakajima--Zwanzig 方程的二阶记忆核近似：
\begin{equation}
\boxed{
\begin{aligned}
\dot\rho_{S,I}(t)
={}&-\frac{\lambda^2}{\hbar^2}
\sum_{\alpha\beta}\int_0^t\dd\tau\,
\Big\{
C_{\alpha\beta}(\tau)
[A_\alpha(t),A_\beta(t-\tau)\rho_{S,I}(t-\tau)]\\
&\hspace{25mm}+C_{\beta\alpha}(-\tau)
[\rho_{S,I}(t-\tau)A_\beta(t-\tau),A_\alpha(t)]
\Big\}+O(\lambda^3).
\end{aligned}}
\label{eq:born_nonmarkov_redfield_kernel}
\end{equation}
因此 Born 截断只回答“保留到耦合的哪一阶”，还没有把有记忆动力学变成时间局域动力学。

在实际有限工作空间内，取一个非负可积包络 $\mathcal C_*(\tau)$，使约化记忆超算符满足
$\|\mathcal K(\tau)\|\le\mathcal C_*(\tau)$，并定义它的前两个时间矩
\[
M_0\equiv\int_0^\infty\dd\tau\,\mathcal C_*(\tau),
\qquad
M_1\equiv\int_0^\infty\dd\tau\,\tau\mathcal C_*(\tau),
\qquad
\tau_B\equiv\frac{M_1}{M_0}.
\]
对目标时间窗再定义
$\rho_*(t)\equiv\sup_{0\le s\le t}\|\rho_S(s)\|$ 与
$\tau_R^{-1}(t)\equiv\sup_{0\le s\le t}\|\dot\rho_S(s)\|/\rho_*(t)$。这样 $\tau_B$ 与 $\tau_R$ 都由实际核和实际约化轨迹确定。用 $g_{SB}$ 表示相互作用产生的典型频率尺度，Born 截断至少要求
\begin{equation}
\boxed{
\epsilon_{\rm B,op}=g_{SB}\tau_B\ll1,}
\label{eq:open_system_born_control_parameter}
\end{equation}
即环境仍保留记忆的时间内，相互作用不足以积累 $O(1)$ 的系统--环境相关性。Markov 近似还要求系统弛豫时间 $\tau_R$ 远大于 $\tau_B$，并且被截去的长时间相关尾足够小：
\begin{equation}
\boxed{
\epsilon_{\rm M}(t)
=\max\!\left[
\frac{\tau_B}{\tau_R(t)},
\frac{\int_t^\infty\dd\tau\,\mathcal C_*(\tau)}
{\int_0^\infty\dd\tau\,\mathcal C_*(\tau)}
\right]\ll1.}
\label{eq:markov_timescale_control_parameter}
\end{equation}
第一项给出把 $\rho_S(t-\tau)$ 局域为 $\rho_S(t)$ 时、相对于自然核尺度 $\rho_*(t)M_0$ 的严格上界；第二项给出把积分上限从 $t$ 延伸到 $\infty$ 的同一归一化尾部上界。二者分别检验历史态变化与相关尾，必须独立满足。

当 $\epsilon_{\rm M}(t)\ll1$ 时，在\eqnrefs{eq:born_nonmarkov_redfield_kernel} 中作两步 Markov 局域化：先以当前态 $\rho_{S,I}(t)$ 替换核内的历史态，再把已经衰减完毕的环境相关积分上限延伸到 $\infty$。得到
\begin{equation}
\boxed{
\begin{aligned}
\dot\rho_{S,I}(t)
={}&-\frac{\lambda^2}{\hbar^2}
\sum_{\alpha\beta}\int_0^\infty\dd\tau\,
\Big\{
C_{\alpha\beta}(\tau)
[A_\alpha(t),A_\beta(t-\tau)\rho_{S,I}(t)]\\
&\hspace{25mm}+C_{\beta\alpha}(-\tau)
[\rho_{S,I}(t)A_\beta(t-\tau),A_\alpha(t)]
\Big\}+O(\lambda^3).
\end{aligned}}
\label{eq:redfield_markov_nonsecular}
\end{equation}
这条二阶、Markov、但尚未做 Bohr 频率久期化的时间局域方程称为\emph{Redfield 方程}。它已经没有历史态，却仍保留不同 Bohr 频率块之间的相干交叉项，因此一般尚不能保证任意初态下的完全正性。下一步的久期近似正是把这些长期快速振荡的交叉块分离开。

完成 Markov 局域化以后，环境通过相关谱进入生成元。定义双边量子噪声谱
\begin{equation}
\boxed{
\Gamma_{\alpha\beta}(\omega)
=\frac{\lambda^2}{\hbar^2}\int_{-\infty}^{\infty}\dd\tau\,
\ee^{\ii\omega\tau}C_{\alpha\beta}(\tau),
\qquad
\boldsymbol\Gamma(\omega)\succeq0.}
\label{eq:quantum_bath_noise_spectrum}
\end{equation}
其半正定性来自环境密度算符的正性。若 $H_S$ 在目标时间窗内可视为静态，系统耦合算符可按 Bohr 频率分解
\[
A_\alpha(t)=\sum_\omega\ee^{-\ii\omega t}A_\alpha(\omega),
\qquad
[H_S,A_\alpha(\omega)]=-\hbar\omega A_\alpha(\omega).
\]
严格周期驱动时，应把这一分解改写成 Floquet 准能块；一般非周期驱动则不能强行套用固定 Bohr 频率的久期化。

把\eqnrefs{eq:redfield_markov_nonsecular} 中的系统算符按 Bohr 频率展开后，不同 $\omega,\omega'$ 频率块仍通过 $\ee^{\ii(\omega'-\omega)t}$ 相互耦合，因此 Redfield 生成元一般尚不具有 GKSL 形式。若实际参与动力学的最小非零 Bohr 频率间隔为 $\Delta\omega_B$，则需要存在粗粒化时间 $\Delta t$ 满足
\[
\tau_B\ll\Delta t\ll\tau_R,
\qquad
\Delta\omega_B\Delta t\gg1.
\]
等价的必要尺度控制量可写成
\begin{equation}
\boxed{
\epsilon_{\rm sec}=\frac1{\Delta\omega_B\tau_R}\ll1.}
\label{eq:secular_control_parameter}
\end{equation}
只有这一条件独立成立，才可丢弃 $\omega\neq\omega'$ 的交叉频率块。最终得到 Born--Markov--久期近似下的 GKSL 主方程
\begin{equation}
\boxed{
\dot\rho_S
=-\frac{\ii}{\hbar}[H_S+H_{\rm LS},\rho_S]
+\sum_\omega\sum_{\alpha\beta}\Gamma_{\alpha\beta}(\omega)
\left[
A_\beta(\omega)\rho_SA_\alpha^\dagger(\omega)
-\frac12\{A_\alpha^\dagger(\omega)A_\beta(\omega),\rho_S\}
\right].}
\label{eq:born_markov_secular_gksl}
\end{equation}
$H_{\rm LS}$ 称为 \emph{Lamb 位移哈密顿量}：环境涨落除了造成不可逆跃迁，还会通过虚交换改变系统能级或共振频率；这一相干频移来自单边环境相关积分的反厄米部分。对每个半正定矩阵 $\boldsymbol\Gamma(\omega)$ 作谱分解，所得非负本征值成为 Lindblad 跃迁率，本征向量给出相应跳跃算符。环境相关谱的正性与久期近似后的频率块分离因而共同保证生成元的完全正结构。

若环境处于温度 $T$ 的 Gibbs 态，$\beta_T=(k_BT)^{-1}$，热平衡还把正频率与负频率过程锁在一起。这个约束称为 Kubo--Martin--Schwinger（KMS）关系；在能量本征态语言中，它就是“上跃迁与下跃迁的热权重必须符合 Boltzmann 因子”的算符形式：
\begin{equation}
\boxed{
\Gamma_{\alpha\beta}(-\omega)
=\ee^{-\beta_T\hbar\omega}\Gamma_{\beta\alpha}(+\omega).}
\label{eq:kms_detailed_balance_electron}
\end{equation}
因此热平衡中的吸收率与发射率由同一谱密度通过详细平衡关系联结。$\omega=0$ 频率块描述不交换净能量的纯退相干，非零 Bohr 频率块描述系统与环境之间的能量交换；二者统一由矩阵谱 $\boldsymbol\Gamma(\omega)$ 控制。

对单通道环境，\eqnrefs{eq:kms_detailed_balance_electron} 直接给出相反能量交换方向的速率比；$\omega=0$ 时热因子为一，而非零 Bohr 频率携带明确的能量交换 $\hbar\omega$。

\begin{derivation}[从\eqnrefs{eq:open_system_joint_liouville_equation}到\eqnrefs{eq:exact_nz_initial_correlation}]
联合态满足
\begin{equation*}
 \dot\rho(t)=\mathcal L(t)\rho(t).
\end{equation*}
插入恒等分解 $1=\mathcal P+\mathcal Q$：
\begin{align*}
 \dot\rho
 &=\mathcal L(\mathcal P+\mathcal Q)\rho\\
 &=\mathcal L\mathcal P\rho+\mathcal L\mathcal Q\rho.
\end{align*}
分别左乘 $\mathcal P$ 与 $\mathcal Q$，利用 $\mathcal P^2=\mathcal P$、$\mathcal Q^2=\mathcal Q$、$\mathcal P\mathcal Q=\mathcal Q\mathcal P=0$，得到两条与原 Liouville 方程完全等价的耦合方程：
\begin{align*}
 \frac{\dd}{\dd t}\mathcal P\rho(t)
 &=\mathcal P\mathcal L(t)\mathcal P\rho(t)
 +\mathcal P\mathcal L(t)\mathcal Q\rho(t),\\
 \frac{\dd}{\dd t}\mathcal Q\rho(t)
 &=\mathcal Q\mathcal L(t)\mathcal P\rho(t)
 +\mathcal Q\mathcal L(t)\mathcal Q\rho(t).
\end{align*}
关键是精确解出第二条 $\mathcal Q$ 方程，而不是把 $\mathcal Q\rho$ 直接置零。

定义限制在 $\mathcal Q$ 子空间内的传播子
\begin{equation*}
 \mathcal U_Q(t,s)
 =\mathcal T\exp\!\left[
 \int_s^t\dd\tau\,
 \mathcal Q\mathcal L(\tau)\mathcal Q
 \right].
\end{equation*}
由时间有序指数的定义，它满足
\begin{align*}
 \partial_t\mathcal U_Q(t,s)
 &=\mathcal Q\mathcal L(t)\mathcal Q\,\mathcal U_Q(t,s),\\
 \mathcal U_Q(s,s)&=\mathcal Q,\\
 \mathcal U_Q(t,s)\mathcal Q
 &=\mathcal Q\mathcal U_Q(t,s)=\mathcal U_Q(t,s).
\end{align*}
为严格推出常数变易公式，定义
\begin{equation*}
 Y(t)\equiv\mathcal U_Q(t_0,t)\mathcal Q\rho(t),
\end{equation*}
其中 $\mathcal U_Q(t_0,t)=\mathcal U_Q(t,t_0)^{-1}$。逆传播子的导数由
\begin{equation*}
 \partial_t\mathcal U_Q(t_0,t)
 =-\mathcal U_Q(t_0,t)\mathcal Q\mathcal L(t)\mathcal Q
\end{equation*}
给出，因此
\begin{align*}
 \dot Y(t)
 ={}&-\mathcal U_Q(t_0,t)
 \mathcal Q\mathcal L(t)\mathcal Q\rho(t)\\
 &+\mathcal U_Q(t_0,t)
 \frac{\dd}{\dd t}\mathcal Q\rho(t).
\end{align*}
把第二条耦合方程代入：
\begin{align*}
 \dot Y(t)
 ={}&-\mathcal U_Q(t_0,t)
 \mathcal Q\mathcal L(t)\mathcal Q\rho(t)\\
 &+\mathcal U_Q(t_0,t)
 \left[
 \mathcal Q\mathcal L(t)\mathcal P\rho(t)
 +\mathcal Q\mathcal L(t)\mathcal Q\rho(t)
 \right]\\
 &=\mathcal U_Q(t_0,t)
 \mathcal Q\mathcal L(t)\mathcal P\rho(t).
\end{align*}
从 $t_0$ 积分到 $t$：
\begin{equation*}
 Y(t)-Y(t_0)
 =\int_{t_0}^{t}\dd s\,
 \mathcal U_Q(t_0,s)
 \mathcal Q\mathcal L(s)\mathcal P\rho(s).
\end{equation*}
又因为
\begin{equation*}
 Y(t_0)=\mathcal Q\rho(t_0),
\end{equation*}
左乘 $\mathcal U_Q(t,t_0)$，并利用组合律
\begin{equation*}
 \mathcal U_Q(t,t_0)\mathcal U_Q(t_0,s)=\mathcal U_Q(t,s),
\end{equation*}
得到精确的 $\mathcal Q$ 分量：
\begin{equation*}
 \mathcal Q\rho(t)
 =\mathcal U_Q(t,t_0)\mathcal Q\rho(t_0)
 +\int_{t_0}^{t}\dd s\,
 \mathcal U_Q(t,s)
 \mathcal Q\mathcal L(s)\mathcal P\rho(s).
\end{equation*}

将得到的 $\mathcal Q\rho(t)$ 精确解代回第一条 $\mathcal P$ 方程：
\begin{align*}
 \frac{\dd}{\dd t}\mathcal P\rho(t)
 ={}&\mathcal P\mathcal L(t)\mathcal P\rho(t)\\
 &+\mathcal P\mathcal L(t)
 \mathcal U_Q(t,t_0)\mathcal Q\rho(t_0)\\
 &+\int_{t_0}^{t}\dd s\,
 \mathcal P\mathcal L(t)
 \mathcal U_Q(t,s)
 \mathcal Q\mathcal L(s)\mathcal P\rho(s).
\end{align*}
正文已经把第二行记作
\begin{equation*}
 \mathcal I(t,t_0)
 \equiv
 \mathcal P\mathcal L(t)
 \mathcal U_Q(t,t_0)\mathcal Q\rho(t_0),
\end{equation*}
就是正文\eqnrefs{eq:exact_nz_initial_correlation}。正文的积分核记作
\begin{equation*}
 \mathcal K(t,s)
 \equiv
 \mathcal P\mathcal L(t)
 \mathcal U_Q(t,s)
 \mathcal Q\mathcal L(s)\mathcal P,
\end{equation*}
便得到正文\eqnrefs{eq:exact_nakajima_zwanzig}：
\begin{equation*}
 \frac{\dd}{\dd t}\mathcal P\rho(t)
 =\mathcal P\mathcal L(t)\mathcal P\rho(t)
 +\int_{t_0}^{t}\dd s\,
 \mathcal K(t,s)\mathcal P\rho(s)
 +\mathcal I(t,t_0).
\end{equation*}

直接检查 $t=t_0$：积分项消失，且
\begin{align*}
 \left.\frac{\dd}{\dd t}\mathcal P\rho(t)\right|_{t_0}
 &=\mathcal P\mathcal L(t_0)
 \left[\mathcal P+\mathcal Q\right]\rho(t_0)\\
 &=\mathcal P\mathcal L(t_0)\rho(t_0),
\end{align*}
与原联合 Liouville 方程一致，所以消元没有丢失初始信息。

若 $\mathcal Q\rho(t_0)=0$，仅有 $\mathcal I=0$；记忆积分仍然保留。因而“因子化初态”与“Markov 近似”是两件不同的事：前者只消去非齐次源，后者还必须证明 $\mathcal K(t,s)$ 在 $t-s$ 上足够短程，才能把历史态 $\mathcal P\rho(s)$ 近似成当前态。
\end{derivation}

\begin{derivation}[从\eqnrefs{eq:nz_reduced_resolvent_definition}到\eqnrefs{eq:nz_projected_full_resolvent_schur_identity}]
假设 $\mathcal L$ 时间无关，并取 $z$ 位于相应预解集，使随后出现的逆算符存在。首先对 $\mathcal Q$ 子空间传播子作单边 Laplace 变换：
\begin{align*}
\int_0^\infty\dd t\,\ee^{-zt}\ee^{t\mathcal Q\mathcal L\mathcal Q}
&=\int_0^\infty\dd t\,
\exp[-t(z-\mathcal Q\mathcal L\mathcal Q)]\\
&=(z-\mathcal Q\mathcal L\mathcal Q)^{-1}.
\end{align*}
因此时间域记忆核
\[
\mathcal K(t)=\mathcal P\mathcal L\mathcal Q
\ee^{t\mathcal Q\mathcal L\mathcal Q}
\mathcal Q\mathcal L\mathcal P
\]
的 Laplace 变换为
\[
\widetilde K(z)
=\mathcal P\mathcal L\mathcal Q
(z-\mathcal Q\mathcal L\mathcal Q)^{-1}
\mathcal Q\mathcal L\mathcal P.
\]
这给出正文\eqnrefs{eq:nz_reduced_resolvent_definition} 中的约化自能。

现在直接从完整预解式验证同一结构。按 $\mathcal P\oplus\mathcal Q$ 分解任意 Liouville 空间向量
\[
|X)=|X_P)+|X_Q),
\qquad
|X_P)=\mathcal P|X),\quad |X_Q)=\mathcal Q|X),
\]
并考虑线性方程
\[
(z-\mathcal L)|X)=|Y).
\]
左右分别投影到两个子空间：
\begin{align*}
(z-\mathcal P\mathcal L\mathcal P)|X_P)
-\mathcal P\mathcal L\mathcal Q|X_Q)
&=|Y_P),\\
-\mathcal Q\mathcal L\mathcal P|X_P)
+(z-\mathcal Q\mathcal L\mathcal Q)|X_Q)
&=|Y_Q).
\end{align*}
从第二式解出被消去分量：
\[
|X_Q)
=(z-\mathcal Q\mathcal L\mathcal Q)^{-1}
\left[|Y_Q)+\mathcal Q\mathcal L\mathcal P|X_P)\right].
\]
代回第一式得到
\begin{align*}
\Bigl[&z-\mathcal P\mathcal L\mathcal P
-\mathcal P\mathcal L\mathcal Q
(z-\mathcal Q\mathcal L\mathcal Q)^{-1}
\mathcal Q\mathcal L\mathcal P\Bigr]|X_P)\\
&=|Y_P)
+\mathcal P\mathcal L\mathcal Q
(z-\mathcal Q\mathcal L\mathcal Q)^{-1}|Y_Q).
\end{align*}
若右端只位于保留子空间，即 $|Y_Q)=0$，则
\[
|X_P)=
\left[z-\mathcal P\mathcal L\mathcal P-\widetilde K(z)\right]^{-1}|Y_P).
\]
另一方面，由 $|X)=(z-\mathcal L)^{-1}|Y)$ 且 $|Y)=|Y_P)$，同一个 $|X_P)$ 也等于
\[
|X_P)=\mathcal P(z-\mathcal L)^{-1}\mathcal P|Y_P).
\]
由于该等式对任意 $|Y_P)$ 成立，得到算符恒等式
\[
\mathcal P(z-\mathcal L)^{-1}\mathcal P
=\left[z-\mathcal P\mathcal L\mathcal P-\widetilde K(z)\right]^{-1},
\]
即正文\eqnrefs{eq:nz_projected_full_resolvent_schur_identity}。
\end{derivation}

\begin{derivation}[从\eqnrefs{eq:exact_second_order_integral_identity_open_system}到\eqnrefs{eq:born_nonmarkov_redfield_kernel}]
相互作用绘景中的联合态满足
\[
\dot\rho_I(t)=-\frac{\ii}{\hbar}[H_I(t),\rho_I(t)],
\qquad
H_I(t)=\lambda\sum_\alpha A_\alpha(t)\otimes B_\alpha(t).
\]
从 $0$ 积分到 $t$：
\[
\rho_I(t)=\rho_I(0)
-\frac{\ii}{\hbar}\int_0^t\dd s\,[H_I(s),\rho_I(s)].
\]
把这个积分式代回原微分方程，并令 $s=t-\tau$，得到正文\eqnrefs{eq:exact_second_order_integral_identity_open_system}：
\[
\dot\rho_I(t)
=-\frac{\ii}{\hbar}[H_I(t),\rho_I(0)]
-\frac1{\hbar^2}\int_0^t\dd\tau\,
[H_I(t),[H_I(t-\tau),\rho_I(t-\tau)]].
\]
取初态因子化 $\rho_I(0)=\rho_S(0)\otimes\rho_B$，并令
$\Tr_B(B_\alpha\rho_B)=0$。对环境取迹以后，一阶项为
\begin{align*}
\Tr_B[H_I(t),\rho_S(0)\otimes\rho_B]
&=\lambda\sum_\alpha
[A_\alpha(t),\rho_S(0)]\Tr_B[B_\alpha(t)\rho_B]\\
&=0.
\end{align*}
Born 截断只用于二阶积分核：
\[
\rho_I(t-\tau)
=\rho_{S,I}(t-\tau)\otimes\rho_B+O(\lambda).
\]
由于核前已有两个 $H_I\propto\lambda$，代入后被丢弃部分从 $O(\lambda^3)$ 开始。

把双重对易子完整展开，记
$\rho_S^-=\rho_{S,I}(t-\tau)$：
\begin{align*}
&[H_I(t),[H_I(t-\tau),\rho_S^-\otimes\rho_B]]\\
={}&H_I(t)H_I(t-\tau)(\rho_S^-\otimes\rho_B)
-H_I(t)(\rho_S^-\otimes\rho_B)H_I(t-\tau)\\
&-H_I(t-\tau)(\rho_S^-\otimes\rho_B)H_I(t)
+(\rho_S^-\otimes\rho_B)H_I(t-\tau)H_I(t).
\end{align*}
代入 $H_I=\lambda\sum A_\alpha\otimes B_\alpha$。第一项的环境偏迹为
\begin{align*}
&\lambda^2\sum_{\alpha\beta}
A_\alpha(t)A_\beta(t-\tau)\rho_S^-\,
\Tr_B[B_\alpha(t)B_\beta(t-\tau)\rho_B]\\
&=\lambda^2\sum_{\alpha\beta}
C_{\alpha\beta}(\tau)
A_\alpha(t)A_\beta(t-\tau)\rho_S^-.
\end{align*}
平稳性使相关函数只依赖时间差。第二项利用迹的循环性：
\begin{align*}
&\Tr_B[B_\alpha(t)\rho_BB_\beta(t-\tau)]\\
&=\Tr_B[B_\beta(t-\tau)B_\alpha(t)\rho_B]
=C_{\beta\alpha}(-\tau),
\end{align*}
故给
\[
-\lambda^2\sum_{\alpha\beta}
C_{\beta\alpha}(-\tau)
A_\alpha(t)\rho_S^-A_\beta(t-\tau).
\]
第三、第四项同理分别为
\begin{align*}
&-\lambda^2\sum_{\alpha\beta}
C_{\alpha\beta}(\tau)
A_\beta(t-\tau)\rho_S^-A_\alpha(t),\\
&+\lambda^2\sum_{\alpha\beta}
C_{\beta\alpha}(-\tau)
\rho_S^-A_\beta(t-\tau)A_\alpha(t).
\end{align*}
把第一项与第三项组合成一个对易子，把第二项与第四项组合成另一个对易子：
\begin{align*}
\Tr_B[H_I(t),[H_I(t-\tau),\rho_S^-\otimes\rho_B]]
=\lambda^2\sum_{\alpha\beta}\Big\{&
C_{\alpha\beta}(\tau)
[A_\alpha(t),A_\beta(t-\tau)\rho_S^-]\\
+&C_{\beta\alpha}(-\tau)
[\rho_S^-A_\beta(t-\tau),A_\alpha(t)]\Big\}.
\end{align*}
代回二次积分恒等式后得到
\begin{align*}
\dot\rho_{S,I}(t)
={}&-\frac{\lambda^2}{\hbar^2}
\sum_{\alpha\beta}\int_0^t\dd\tau\,\Big\{
C_{\alpha\beta}(\tau)
[A_\alpha(t),A_\beta(t-\tau)\rho_{S,I}(t-\tau)]\\
&\hspace{30mm}+C_{\beta\alpha}(-\tau)
[\rho_{S,I}(t-\tau)A_\beta(t-\tau),A_\alpha(t)]\Big\}
+O(\lambda^3),
\end{align*}
即正文\eqnrefs{eq:born_nonmarkov_redfield_kernel}。
\end{derivation}

\begin{derivation}[从\eqnrefs{eq:born_nonmarkov_redfield_kernel}到\eqnrefs{eq:markov_timescale_control_parameter}]
把正文\eqnrefs{eq:born_nonmarkov_redfield_kernel} 抽象写成
\[
\dot\rho_S(t)=\int_0^t\dd\tau\,\mathcal K(\tau)\rho_S(t-\tau),
\qquad
\|\mathcal K(\tau)\|\le\mathcal C_*(\tau).
\]
Markov 局域化把历史态换成当前态。历史替换造成的差严格为
\begin{align*}
\delta\dot\rho_{S,{\rm hist}}(t)
&\equiv
\int_0^t\dd\tau\,\mathcal K(\tau)
[\rho_S(t-\tau)-\rho_S(t)],\\
\rho_S(t-\tau)-\rho_S(t)
&=-\int_{t-\tau}^{t}\dd s\,\dot\rho_S(s).
\end{align*}
由正文定义
$\sup_{0\le s\le t}\|\dot\rho_S(s)\|=\rho_*(t)/\tau_R(t)$，逐次使用三角不等式与超算符范数的次乘性，得到
\begin{align*}
\|\delta\dot\rho_{S,{\rm hist}}(t)\|
&\le
\int_0^t\dd\tau\,\|\mathcal K(\tau)\|
\int_{t-\tau}^{t}\dd s\,\|\dot\rho_S(s)\|\\
&\le
\frac{\rho_*(t)}{\tau_R(t)}
\int_0^t\dd\tau\,\tau\mathcal C_*(\tau)\\
&\le
\frac{\rho_*(t)M_1}{\tau_R(t)}.
\end{align*}
用自然核尺度 $\rho_*(t)M_0$ 归一化，便有严格上界
\begin{equation*}
\frac{\|\delta\dot\rho_{S,{\rm hist}}(t)\|}{\rho_*(t)M_0}
\le
\frac{M_1}{M_0\tau_R(t)}
=\frac{\tau_B}{\tau_R(t)}.
\end{equation*}
第二个误差来自把已经局域化的有限积分上限延伸到 $\infty$。定义
\[
\delta\dot\rho_{S,{\rm tail}}(t)
\equiv
\int_t^\infty\dd\tau\,\mathcal K(\tau)\rho_S(t),
\]
则
\begin{align*}
\|\delta\dot\rho_{S,{\rm tail}}(t)\|
&\le
\rho_*(t)\int_t^\infty\dd\tau\,\mathcal C_*(\tau),\\
\frac{\|\delta\dot\rho_{S,{\rm tail}}(t)\|}{\rho_*(t)M_0}
&\le
\frac{\int_t^\infty\dd\tau\,\mathcal C_*(\tau)}
{\int_0^\infty\dd\tau\,\mathcal C_*(\tau)}.
\end{align*}
取两个严格上界的最大值，就得到正文\eqnrefs{eq:markov_timescale_control_parameter}。
\end{derivation}

\begin{derivation}[从\eqnrefs{eq:quantum_bath_noise_spectrum}到\eqnrefs{eq:born_markov_secular_gksl}]
Markov 局域化后的 Redfield 方程仍含不同 Bohr 频率之间的交叉项。先分解系统算符：
\[
A_\alpha(t)=\sum_\omega\ee^{-\ii\omega t}A_\alpha(\omega),
\qquad
[H_S,A_\alpha(\omega)]=-\hbar\omega A_\alpha(\omega).
\]
于是
\[
A_\beta(t-\tau)
=\sum_\omega\ee^{-\ii\omega t}\ee^{+\ii\omega\tau}A_\beta(\omega).
\]
定义单边相关积分
\[
G_{\alpha\beta}(\omega)
\equiv\frac{\lambda^2}{\hbar^2}
\int_0^\infty\dd\tau\,
\ee^{\ii\omega\tau}C_{\alpha\beta}(\tau).
\]
平稳环境满足
$C_{\beta\alpha}(-\tau)=C_{\alpha\beta}^*(\tau)$，故正文的双边谱可以写成
\begin{align*}
\Gamma_{\alpha\beta}(\omega)
&=\frac{\lambda^2}{\hbar^2}
\int_{-\infty}^{\infty}\dd\tau\,
\ee^{\ii\omega\tau}C_{\alpha\beta}(\tau)\\
&=G_{\alpha\beta}(\omega)+G_{\beta\alpha}^*(\omega).
\end{align*}
因此单边积分唯一分解为
\[
G_{\alpha\beta}(\omega)
=\frac12\Gamma_{\alpha\beta}(\omega)
+\ii S_{\alpha\beta}(\omega),
\qquad
S_{\alpha\beta}=S_{\beta\alpha}^*.
\]

将 Bohr 分解代入 Markov--Redfield 生成元后，任意 $\omega,\omega'$ 交叉块都乘有
$\ee^{\ii(\omega'-\omega)t}$。在粗粒化时间 $\Delta t$ 上，
\begin{align*}
\frac1{\Delta t}\int_t^{t+\Delta t}\dd s\,
\ee^{\ii(\omega'-\omega)s}
&=\ee^{\ii(\omega'-\omega)(t+\Delta t/2)}
\frac{2\sin[(\omega'-\omega)\Delta t/2]}
{(\omega'-\omega)\Delta t}.
\end{align*}
当 $\tau_B\ll\Delta t\ll\tau_R$ 且
$|\omega'-\omega|\Delta t\gg1$ 时，$\omega'\ne\omega$ 的块被 sinc 因子压低，只保留相同 Bohr 频率的块。

对固定 $\omega$，久期化后的生成元为
\begin{align*}
\mathcal L_\omega\rho
={}&\sum_{\alpha\beta}
G_{\alpha\beta}(\omega)
[A_\beta(\omega)\rho,A_\alpha^\dagger(\omega)]\\
&+\sum_{\alpha\beta}
G_{\beta\alpha}^*(\omega)
[A_\alpha(\omega),\rho A_\beta^\dagger(\omega)].
\end{align*}
代入 $G=\Gamma/2+\ii S$。厄米的 $\Gamma$ 部分逐项展开为
\begin{align*}
\mathcal D_\omega\rho
={}&\sum_{\alpha\beta}\Gamma_{\alpha\beta}(\omega)
A_\beta(\omega)\rho A_\alpha^\dagger(\omega)\\
&-\frac12\sum_{\alpha\beta}\Gamma_{\alpha\beta}(\omega)
A_\alpha^\dagger(\omega)A_\beta(\omega)\rho\\
&-\frac12\sum_{\alpha\beta}\Gamma_{\alpha\beta}(\omega)
\rho A_\alpha^\dagger(\omega)A_\beta(\omega)\\
={}&\sum_{\alpha\beta}\Gamma_{\alpha\beta}(\omega)
\left[
A_\beta(\omega)\rho A_\alpha^\dagger(\omega)
-\frac12\{A_\alpha^\dagger(\omega)A_\beta(\omega),\rho\}
\right].
\end{align*}
反厄米部分定义
\[
H_{\rm LS}
\equiv\hbar\sum_\omega\sum_{\alpha\beta}
S_{\alpha\beta}(\omega)
A_\alpha^\dagger(\omega)A_\beta(\omega).
\]
由 $S_{\alpha\beta}=S_{\beta\alpha}^*$ 得
$H_{\rm LS}=H_{\rm LS}^\dagger$，相应项组合为
\[
-\frac{\ii}{\hbar}[H_{\rm LS},\rho].
\]
恢复 Schr\"odinger 绘景并把所有 Bohr 频率块求和，得到
\[
\dot\rho_S
=-\frac{\ii}{\hbar}[H_S+H_{\rm LS},\rho_S]
+\sum_\omega\sum_{\alpha\beta}\Gamma_{\alpha\beta}(\omega)
\left[
A_\beta(\omega)\rho_SA_\alpha^\dagger(\omega)
-\frac12\{A_\alpha^\dagger(\omega)A_\beta(\omega),\rho_S\}
\right],
\]
即正文\eqnrefs{eq:born_markov_secular_gksl}。
\end{derivation}

\begin{derivation}[从\eqnrefs{eq:quantum_bath_noise_spectrum}到\eqnrefs{eq:kms_detailed_balance_electron}]
设环境处于 Gibbs 态
\[
\rho_B=\frac{\ee^{-\beta_T H_B}}{Z},
\qquad
H_B|n\rangle=E_n|n\rangle,
\qquad
Z=\sum_n\ee^{-\beta_TE_n}.
\]
相互作用绘景环境算符满足
\[
B_\alpha(t)=\ee^{+\ii H_Bt/\hbar}B_\alpha\ee^{-\ii H_Bt/\hbar}.
\]
在能量本征态间插入完备关系，矩阵元为
\[
\langle n|B_\alpha(t)|m\rangle
=\ee^{+\ii(E_n-E_m)t/\hbar}(B_\alpha)_{nm}.
\]
因此平稳二点函数可以写成
\begin{align*}
C_{\alpha\beta}(t)
&=\Tr_B[B_\alpha(t)B_\beta(0)\rho_B]\\
&=\frac1Z\sum_{n,m}
\ee^{-\beta_TE_n}
\ee^{-\ii(E_m-E_n)t/\hbar}
(B_\alpha)_{nm}(B_\beta)_{mn}.
\end{align*}
按正文的双边 Fourier 定义
\[
\Gamma_{\alpha\beta}(\omega)
=\frac{\lambda^2}{\hbar^2}
\int_{-\infty}^{+\infty}\dd t\,
\ee^{+\ii\omega t}C_{\alpha\beta}(t),
\]
使用
\[
\int_{-\infty}^{+\infty}\dd t\,
\ee^{\ii(\omega-\Omega)t}=2\pi\delta(\omega-\Omega),
\]
得到
\[
\Gamma_{\alpha\beta}(\omega)
=\frac{2\pi\lambda^2}{\hbar^2Z}
\sum_{n,m}\ee^{-\beta_TE_n}
(B_\alpha)_{nm}(B_\beta)_{mn}
\delta\!\left(\omega-\frac{E_m-E_n}{\hbar}\right).
\]
现在将 $\omega$ 换成 $-\omega$：
\[
\Gamma_{\alpha\beta}(-\omega)
=\frac{2\pi\lambda^2}{\hbar^2Z}
\sum_{n,m}\ee^{-\beta_TE_n}
(B_\alpha)_{nm}(B_\beta)_{mn}
\delta\!\left(\omega+\frac{E_m-E_n}{\hbar}\right).
\]
交换哑指标 $n\leftrightarrow m$ 后，
\[
\Gamma_{\alpha\beta}(-\omega)
=\frac{2\pi\lambda^2}{\hbar^2Z}
\sum_{n,m}\ee^{-\beta_TE_m}
(B_\alpha)_{mn}(B_\beta)_{nm}
\delta\!\left(\omega-\frac{E_m-E_n}{\hbar}\right).
\]
在该 $\delta$ 函数的支撑上
\[
E_m-E_n=\hbar\omega,
\]
所以 Boltzmann 权重严格满足
\[
\ee^{-\beta_TE_m}
=\ee^{-\beta_T(E_m-E_n)}\ee^{-\beta_TE_n}
=\ee^{-\beta_T\hbar\omega}\ee^{-\beta_TE_n}.
\]
将与求和指标无关的热因子提出，并按照 $\Gamma_{\beta\alpha}(+\omega)$ 的定义识别剩余和式，得到
\[
\Gamma_{\alpha\beta}(-\omega)
=\ee^{-\beta_T\hbar\omega}
\Gamma_{\beta\alpha}(+\omega),
\]
即正文\eqnrefs{eq:kms_detailed_balance_electron}。
\end{derivation}


\subsection{高斯环境与涨落--耗散定理（FDT）}

上一节已经把环境的作用写成精确的记忆核，但对任意多体环境而言，这个记忆核本身仍需要知道环境的高阶相关函数，因此通常不能直接求出。这里选择一类既常见又可以完整计算的环境：环境 Hamiltonian 对玻色坐标至多是二次的，系统与环境之间只保留线性耦合，并且环境初态是高斯态。所谓高斯态，是指任意高阶相关函数都可以由二点相关函数组合得到。正因为这一性质，环境自由度可以在这个模型内精确消去；所得二次核来自高斯统计本身，而不是弱耦合 Born 截断。

取单通道耦合
\[
\hat H_{SB}=\hat A\otimes\hat B,
\qquad
\hat B=\sum_\lambda
(g_\lambda\hat b_\lambda+g_\lambda^*\hat b_\lambda^\dagger).
\]
若环境有非零平均值，先定义
\[
\bar B(t)\equiv\langle\hat B(t)\rangle,
\qquad
\delta\hat B(t)\equiv\hat B(t)-\bar B(t),
\]
并把 $\hat A\bar B(t)$ 吸收到重整化后的系统哈密顿量中。这样，真正需要描述的环境作用只剩下零均值算符 $\delta\hat B$ 的涨落。取可分离初态 $\rho_S(t_i)\otimes\rho_B$ 时可以得到最常用的闭式结果；若系统与环境在初始时刻已经相关，则还必须保留\eqnrefs{eq:exact_nz_initial_correlation} 中的边界源项。

为区分“涨落大小”与“因果响应”，定义较大/较小二点关联函数（greater/lesser correlators）
\[
C^>(t,s)=\langle\delta\hat B(t)\delta\hat B(s)\rangle,
\qquad
C^<(t,s)=\langle\delta\hat B(s)\delta\hat B(t)\rangle,
\]
以及它们的对称/反对称组合
\begin{equation}
\boxed{
C_{\rm sym}(t,s)
\equiv\frac12\langle\{\delta\hat B(t),\delta\hat B(s)\}\rangle,
\qquad
C_{\rm asym}(t,s)
\equiv\frac{1}{2\ii}\langle[\delta\hat B(t),\delta\hat B(s)]\rangle.}
\label{eq:noise_dissipation_kernels_general}
\end{equation}
$C_{\rm sym}$ 衡量环境涨落的对称部分，$C_{\rm asym}$ 来自算符对易子并决定因果反作用。密度矩阵的演化天然包含两条历史：ket 沿时间向前演化，bra 对应另一条历史。分别把系统变量沿这两条历史的取值记为 $A_+$ 与 $A_-$，并定义
\[
A_\Delta=A_+-A_-,
\qquad
A_\Sigma=\frac{A_++A_-}{2}.
\]
把环境自由度迹掉以后，环境对这两条系统历史的全部影响可以压缩为一个乘法因子，这个对象称为 Feynman--Vernon 影响泛函，全文记为 $\mathscr I_{\rm FV}$。由于环境是高斯的，它由二点函数完全决定：
\begin{equation}
\boxed{
\mathscr I_{\rm FV}[A_+,A_-]
=\exp\!\left\{-\frac{1}{\hbar^2}
\int_{t_i}^{t_f}\dd t\int_{t_i}^{t}\dd s\,
A_\Delta(t)
\left[C_{\rm sym}(t,s)A_\Delta(s)+2\ii C_{\rm asym}(t,s)A_\Sigma(s)\right]
\right\}.}
\label{eq:gaussian_feynman_vernon_influence}
\end{equation}
这个影响泛函保留 $t$ 与 $s$ 两个独立时间，因此一般描述有色噪声和非局域记忆。若 $C_{\rm sym}(t,s)$、$C_{\rm asym}(t,s)$ 在相关时间远短于系统时间尺度时收缩成近局域核，才分别产生白噪声、指数型退相干或局域摩擦的 Markov 极限。$C_{\rm sym}$ 抑制两条历史之间的相干，$C_{\rm asym}$ 改变相位并产生因果反作用；二者来自同一个环境二点函数。

只有进一步要求环境平稳且处于 Gibbs 态，噪声与响应才由热平衡约束相互锁定。按全文傅里叶约定定义
\begin{equation}
\boxed{
S_B^{>/<}(\omega)
\equiv\int_{-\infty}^{\infty}\dd\tau\,
\ee^{\ii\omega\tau}C^{>/<}(\tau).}
\label{eq:bath_greater_lesser_spectra}
\end{equation}
若 $\rho_B=\ee^{-\beta_T H_B}/Z_B$，其中 $\beta_T=(k_BT)^{-1}$，谱展开直接给出 KMS 细致平衡关系
\[
S_B^<(\omega)=\ee^{-\beta_T\hbar\omega}S_B^>(\omega).
\]
再定义
\[
\widetilde C_{\rm sym}(\omega)
\equiv\frac12\left[S_B^>(\omega)+S_B^<(\omega)\right],
\qquad
\mathscr S_B^{(-)}(\omega)
\equiv\frac{1}{\hbar}\left[S_B^>(\omega)-S_B^<(\omega)\right],
\]
以及推迟响应函数
\[
\chi_B^R(\tau)
\equiv-\frac{\ii}{\hbar}\Theta(\tau)
\langle[\delta\hat B(\tau),\delta\hat B(0)]\rangle.
\]
在上述符号约定下 $\mathscr S_B^{(-)}=-2\operatorname{Im}\chi_B^R$，于是涨落--耗散定理（fluctuation--dissipation 定理）为
\begin{equation}
\boxed{
\widetilde C_{\rm sym}(\omega)
=-\hbar\coth\!\left(\frac{\beta_T\hbar\omega}{2}\right)
\operatorname{Im}\chi_B^R(\omega).}
\label{eq:fdt_from_defined_correlators}
\end{equation}
热平衡结构沿“热态 $\to$ KMS 条件 $\to$ 涨落--耗散关系”单向确定；非平衡高斯环境仍满足\eqnrefs{eq:gaussian_feynman_vernon_influence}，却一般不满足\eqnrefs{eq:fdt_from_defined_correlators}。

同一组核也可以写成随机力语言。取系统坐标 $x_S$ 线性耦合环境，$A=x_S$。由于 $C_{\rm sym}$ 是半正定协方差核，Hubbard--Stratonovich 变换允许引入零均值高斯随机力 $\xi$，满足
\[
\langle\xi(t)\xi(s)\rangle=C_{\rm sym}(t,s).
\]
对粗粒化有效作用量在路径差 $x_{S,\Delta}$ 的线性阶取驻值条件，得到广义 Langevin 方程
\begin{equation}
\boxed{
M\ddot x_S(t)+V'[x_S(t)]
+\int_{t_i}^{t}\dd s\,\chi_B^R(t,s)x_S(s)
=\xi(t),
\qquad
\langle\xi(t)\xi(s)\rangle=C_{\rm sym}(t,s).}
\label{eq:generalized_langevin_from_influence}
\end{equation}
因此记忆力与随机力不是两个独立的经验输入；在线性高斯环境中，它们分别读取同一量子关联函数的对易子与反对易子部分。只有环境记忆时间远短于系统时间尺度时，记忆积分才可进一步局域化；再取 $k_BT\gg\hbar\omega$，量子涨落--耗散关系才退化成经典 Einstein 关系。

在同一个厄米耦合 $H_I=\lambda x_S\otimes B$ 下，Born 二阶但仍保留非 Markov 记忆的方程可以写成
\begin{equation}
\boxed{
\begin{aligned}
\dot\rho_{S,I}(t)
={}&-\frac{\lambda^2}{\hbar^2}\int_0^t\dd\tau\,
C_{\rm sym}(\tau)[x_S(t),[x_S(t-\tau),\rho_{S,I}(t-\tau)]]\\
&-\frac{\ii\lambda^2}{\hbar^2}\int_0^t\dd\tau\,
C_{\rm asym}(\tau)[x_S(t),\{x_S(t-\tau),\rho_{S,I}(t-\tau)\}]
+O(\lambda^3).
\end{aligned}}
\label{eq:born_nonmarkov_noise_dissipation_form}
\end{equation}
Born 主方程与影响泛函使用同一组 $C_{\rm sym},C_{\rm asym}$；前者把约化动力学按耦合展开，后者则在高斯环境中对环境收缩精确重求和。

对谐振子系统 $V(x_S)=M\Omega_0^2x_S^2/2$，Langevin 方程在频域定义经过环境修饰的因果响应
\begin{equation}
\boxed{
\mathcal R_{x_S}(\omega)
\equiv\left[
M\bigl(\Omega_0^2-(\omega+\ii0^+)^2\bigr)
+\chi_B^R(\omega)
\right]^{-1},
\qquad
x_S(\omega)=\mathcal R_{x_S}(\omega)\xi(\omega).}
\label{eq:dressed_oscillator_retarded_response}
\end{equation}
于是平稳对称坐标谱为
\begin{equation}
\boxed{
S_{x_Sx_S}^{\rm sym}(\omega)
=|\mathcal R_{x_S}(\omega)|^2\widetilde C_{\rm sym}(\omega).}
\label{eq:dressed_coordinate_spectrum_dp68}
\end{equation}
若环境处于热平衡，把\eqnrefs{eq:fdt_from_defined_correlators} 代入便得到系统可观测量自身的涨落--耗散关系，
\begin{equation}
\boxed{
S_{x_Sx_S}^{\rm sym}(\omega)
=\hbar\coth\!\left(\frac{\beta_T\hbar\omega}{2}\right)
[-\operatorname{Im}\chi_{x_S}^R(\omega)],}
\label{eq:dressed_system_fdt_from_bath_fdt}
\end{equation}
按全文符号约定，$\chi_{x_S}^R=-\mathcal R_{x_S}$。因此环境的涨落--耗散关系经过推迟响应修饰后自动传递给系统可观测量。


高斯环境之所以特殊，是因为所有三阶及更高连通相关函数都严格为零；这不是一般线性环境的必然性质。若材料自由度本身非谐、环境被制备在 Fock/cat/squeezed-non-Gaussian 态，或电子耦合到具有明显跳跃统计的离散缺陷，二点函数不足以确定电子动力学。一般情况下仍可以沿闭合时间路径组织环境迹，但影响泛函的对数必须保留完整的连通累积量层级。把两条 Keldysh 支路统一记作轮廓变量 $z\in\mathcal C$，并把支路方向符号吸收到轮廓源 $A_{\mathcal C}(z)$ 中，则
\begin{equation}
\boxed{
\ln\mathscr I[A_+,A_-]
=\sum_{n=1}^{\infty}
\frac1{n!}
\left(-\frac{\ii}{\hbar}\right)^n
\int_{\mathcal C}\dd z_1\cdots\dd z_n\,
A_{\mathcal C}(z_1)\cdots A_{\mathcal C}(z_n)
C_c^{(n)}(z_1,\ldots,z_n),}
\label{eq:non_gaussian_influence_cumulant_hierarchy}
\end{equation}
其中
\[
C_c^{(n)}(z_1,\ldots,z_n)
\equiv
\left\langle\!\left\langle
\mathcal T_{\mathcal C}
\delta B(z_1)\cdots\delta B(z_n)
\right\rangle\!\right\rangle
\]
是环境的轮廓有序 $n$ 点连通累积量。高斯环境正是
$C_c^{(n>2)}=0$ 的特例，所以\eqnrefs{eq:gaussian_feynman_vernon_influence} 才在二阶终止；非高斯环境不能靠“把 Green 函数换成一个更复杂的 Green 函数”吸收掉，因为 Green 函数只编码二点响应。

一个最直观的可计算极限是纯退相干型耦合：系统路径只通过一个与环境算符对易的累积相位
\begin{equation*}
\hat\Phi
\equiv\frac1\hbar\int_{t_i}^{t_f}\dd t\,f(t)\delta B(t)
\end{equation*}
读取环境。此时路径相干因子就是环境相位的特征函数
\begin{equation}
\boxed{
\mathcal V\ee^{\ii\varphi}
=\left\langle\ee^{-\ii\hat\Phi}\right\rangle
=\exp\!\left[
\sum_{n=1}^{\infty}
\frac{(-\ii)^n}{n!}\mathcal C_n^{(\Phi)}
\right],}
\label{eq:non_gaussian_phase_characteristic_cumulants}
\end{equation}
其中
$\mathcal C_n^{(\Phi)}=\langle\!\langle\hat\Phi^n\rangle\!\rangle$。
若 $\langle\hat\Phi\rangle=0$，前四阶给
\begin{equation*}
\ln(\mathcal V\ee^{\ii\varphi})
=-\frac{\mathcal C_2^{(\Phi)}}{2}
+\frac{\ii\mathcal C_3^{(\Phi)}}{6}
+\frac{\mathcal C_4^{(\Phi)}}{24}+O(\mathcal C_5^{(\Phi)}).
\end{equation*}
所以二阶 cumulant 控制 Gaussian 可见度衰减，三阶 cumulant 首先产生由偏斜统计引起的附加相位，四阶 cumulant 则修正 Gaussian 指数包络。为了判断“近似 Gaussian”是否受控，可在实际滤波函数 $f(t)$ 已经固定后定义
\begin{equation}
\boxed{
\epsilon_{\rm NG}^{(3)}
\equiv
\frac{|\mathcal C_3^{(\Phi)}|}{[\mathcal C_2^{(\Phi)}]^{3/2}},
\qquad
\epsilon_{\rm NG}^{(4)}
\equiv
\frac{|\mathcal C_4^{(\Phi)}|}{[\mathcal C_2^{(\Phi)}]^{2}}.}
\label{eq:non_gaussian_standardized_cumulant_control}
\end{equation}
这些量必须连同实际电子时间窗一起计算；环境本身“弱非线性”并不自动保证电子所采样的相位统计接近 Gaussian。

这条层级对自由电子尤其有用。EELS 或弱 PINEM 的低阶谱通常主要读取环境的二点关联；经过整形的电子相干、重复条件测量和 heralding 则可以把更高连通 cumulant 转换成可观测电子统计。

因此，两个环境即使具有完全相同的二点 Green/FDT 谱，也仍可能因为三点、四点连通函数不同而给出不同的电子输出。只用二点 Green 张量预测全部量子统计时，所依赖的额外条件是“环境动力学线性且相关初态为 Gaussian”，而不是 Green 张量本身自动包含所有高阶统计。



% Detailed bridge between Born memory kernels and the Gaussian/Langevin form.

\begin{derivation}[从\eqnrefs{eq:born_nonmarkov_redfield_kernel}到\eqnrefs{eq:born_nonmarkov_noise_dissipation_form}]
记

\begin{equation*}
A=q(t),\qquad B_\tau=q(t-\tau),\qquad \rho_\tau=\rho_{S,I}(t-\tau),
\end{equation*}
并把平稳环境相关函数写成
\begin{equation*}
C(\tau)=C_{\rm sym}(\tau)+\ii C_{\rm asym}(\tau),\qquad
C(-\tau)=C_{\rm sym}(\tau)-\ii C_{\rm asym}(\tau).
\end{equation*}
二阶 Born 核中的算符组合为
\begin{align*}
\mathcal K_\tau
={}&C(\tau)[A,B_\tau\rho_\tau]
+C(-\tau)[\rho_\tau B_\tau,A]\\
={}&(C_{\rm sym}+\ii C_{\rm asym})(AB_\tau\rho_\tau-B_\tau\rho_\tau A)
+(C_{\rm sym}-\ii C_{\rm asym})(\rho_\tau B_\tau A-A\rho_\tau B_\tau).
\end{align*}
先收集 $C_{\rm sym}$ 项：
\begin{align*}
\mathcal K_\tau^{({\rm sym})}
&=C_{\rm sym}(AB_\tau\rho_\tau-A\rho_\tau B_\tau-B_\tau\rho_\tau A+\rho_\tau B_\tau A)\\
&=C_{\rm sym}[A,[B_\tau,\rho_\tau]].
\end{align*}
再收集 $C_{\rm asym}$ 项：
\begin{align*}
\mathcal K_\tau^{({\rm asym})}
&=\ii C_{\rm asym}(AB_\tau\rho_\tau+A\rho_\tau B_\tau-B_\tau\rho_\tau A-\rho_\tau B_\tau A)\\
&=\ii C_{\rm asym}[A,\{B_\tau,\rho_\tau\}].
\end{align*}
因此
\begin{equation*}
\mathcal K_\tau
=C_{\rm sym}(\tau)[q(t),[q(t-\tau),\rho_\tau]]
+\ii C_{\rm asym}(\tau)[q(t),\{q(t-\tau),\rho_\tau\}].
\end{equation*}
把它代回二阶 Born 方程前面的 $-\lambda^2/\hbar^2$ 系数和 $\tau$ 积分，就得到\eqnrefs{eq:born_nonmarkov_noise_dissipation_form}。
\end{derivation}

\begin{derivation}[从\eqnrefs{eq:generalized_langevin_from_influence}到\eqnrefs{eq:dressed_oscillator_retarded_response}]
对谐振子势 $V(q)=M\Omega_0^2q^2/2$，正文\eqnrefs{eq:generalized_langevin_from_influence} 可写成
\[
M\ddot q(t)+M\Omega_0^2q(t)
+\int_{-\infty}^{t}\dd s\,\chi_B^R(t-s)q(s)
=\xi(t).
\]
在推迟核中令 $\tau=t-s$，则
\[
\int_{-\infty}^{t}\dd s\,\chi_B^R(t-s)q(s)
=\int_0^\infty\dd\tau\,\chi_B^R(\tau)q(t-\tau).
\]
采用正文的 Fourier 约定
$q(t)=\int \dd\omega\,q(\omega)\ee^{-\ii\omega t}/(2\pi)$，卷积项变成
\begin{align*}
&\int_0^\infty\dd\tau\,\chi_B^R(\tau)
\int\frac{\dd\omega}{2\pi}q(\omega)
\ee^{-\ii\omega(t-\tau)}\\
&=\int\frac{\dd\omega}{2\pi}
\left[
\int_0^\infty\dd\tau\,
\chi_B^R(\tau)\ee^{+\ii\omega\tau}
\right]q(\omega)\ee^{-\ii\omega t}\\
&=\int\frac{\dd\omega}{2\pi}
\chi_B^R(\omega)q(\omega)\ee^{-\ii\omega t}.
\end{align*}
同时 $\ddot q(t)\mapsto-\omega^2q(\omega)$。因此每个频率分量满足
\[
\left[
M\bigl(\Omega_0^2-(\omega+\ii0^+)^2\bigr)
+\chi_B^R(\omega)
\right]q(\omega)=\xi(\omega).
\]
对乘在 $q(\omega)$ 前的标量核求逆，得到
\[
q(\omega)=\mathcal R_q(\omega)\xi(\omega),
\qquad
\mathcal R_q(\omega)
=\frac{1}{M[\Omega_0^2-(\omega+\ii0^+)^2]+\chi_B^R(\omega)},
\]
即正文\eqnrefs{eq:dressed_oscillator_retarded_response}。
\end{derivation}

\begin{derivation}[从\eqnrefs{eq:dressed_coordinate_spectrum_dp68}到\eqnrefs{eq:dressed_system_fdt_from_bath_fdt}]

正文\eqnrefs{eq:dressed_coordinate_spectrum_dp68} 表明系统坐标谱等于响应模方乘环境噪声谱。热平衡时，环境噪声谱又由\eqnrefs{eq:fdt_from_defined_correlators} 与环境响应的虚部关联。利用复数逆函数的虚部恒等式，把 $|\mathcal R_{x_S}|^2[-\operatorname{Im}\chi_B^R]$ 改写成系统自身响应的耗散部分，就得到\eqnrefs{eq:dressed_system_fdt_from_bath_fdt}。

令
\[
D(\omega)=M[\Omega_0^2-(\omega+\ii0^+)^2]+\chi_B^R(\omega),
\qquad
\mathcal R_{x_S}=D^{-1}.
\]
于是
\[
\operatorname{Im}\mathcal R_{x_S}
=-|\mathcal R_{x_S}|^2\operatorname{Im}\chi_B^R.
\]
按正文约定 $\chi_{x_S}^R=-\mathcal R_{x_S}$，因此
\[
-\operatorname{Im}\chi_{x_S}^R
=|\mathcal R_{x_S}|^2[-\operatorname{Im}\chi_B^R].
\]
将这一恒等式和环境 FDT~\eqnrefs{eq:fdt_from_defined_correlators} 代入\eqnrefs{eq:dressed_coordinate_spectrum_dp68}，得到\eqnrefs{eq:dressed_system_fdt_from_bath_fdt}。
\end{derivation}



\begin{derivation}[从\eqnrefs{eq:non_gaussian_influence_cumulant_hierarchy}到\eqnrefs{eq:non_gaussian_phase_characteristic_cumulants}]
从任意环境平均
\[
\mathscr I
=\left\langle
\mathcal T_{\mathcal C}
\exp\!\left[-\frac{\ii}{\hbar}
\int_{\mathcal C}\dd z\,
A_{\mathcal C}(z)\delta B(z)
\right]
\right\rangle
\]
出发。定义轮廓源
\[
X\equiv-\frac{\ii}{\hbar}
\int_{\mathcal C}\dd z\,
A_{\mathcal C}(z)\delta B(z).
\]
普通矩生成函数与连通生成函数的关系是
\[
\ln\langle\ee^X\rangle
=\sum_{n=1}^{\infty}\frac1{n!}
\langle\!\langle X^n\rangle\!\rangle.
\]
由于 $X$ 对环境算符是线性的，第 $n$ 阶连通矩直接成为
\begin{align*}
\langle\!\langle X^n\rangle\!\rangle
={}&\left(-\frac{\ii}{\hbar}\right)^n
\int_{\mathcal C}\dd z_1\cdots\dd z_n\,
A_{\mathcal C}(z_1)\cdots A_{\mathcal C}(z_n)\\
&\times
\left\langle\!\left\langle
\mathcal T_{\mathcal C}
\delta B(z_1)\cdots\delta B(z_n)
\right\rangle\!\right\rangle.
\end{align*}
代回即得到正文\eqnrefs{eq:non_gaussian_influence_cumulant_hierarchy}。若环境是 Gaussian，Wick 定理使所有 $n>2$ 连通 cumulant 消失，因此二阶结果不是“弱耦合截断”，而是在 Gaussian 条件下的精确终止。

纯退相干极限中，系统两条路径只给环境施加一个 c-number 滤波函数 $f(t)$，且相关耦合在不同时刻彼此对易。定义
\[
\hat\Phi=\frac1\hbar\int\dd t\,f(t)\delta B(t).
\]
此时轮廓有序问题退化为普通特征函数
\[
\chi_\Phi(1)=\langle\ee^{-\ii\hat\Phi}\rangle.
\]
对随机变量/可交换算符的 cumulant 生成函数
\[
K_\Phi(s)
\equiv\ln\langle\ee^{s\hat\Phi}\rangle
=\sum_{n=1}^{\infty}\frac{s^n}{n!}\mathcal C_n^{(\Phi)},
\]
取 $s=-\ii$，得到
\[
\ln\chi_\Phi(1)
=\sum_{n=1}^{\infty}\frac{(-\ii)^n}{n!}\mathcal C_n^{(\Phi)},
\]
即正文\eqnrefs{eq:non_gaussian_phase_characteristic_cumulants}。

若均值已经吸收到确定性相位中，$\mathcal C_1^{(\Phi)}=0$。逐阶展开：
\[
(-\ii)^2=-1,
\qquad
(-\ii)^3=+\ii,
\qquad
(-\ii)^4=+1,
\]
因此
\[
\ln\chi_\Phi
=-\frac{\mathcal C_2^{(\Phi)}}{2}
+\frac{\ii\mathcal C_3^{(\Phi)}}{6}
+\frac{\mathcal C_4^{(\Phi)}}{24}+\cdots.
\]
模与辐角分别由实部和虚部给出。若只保留 Gaussian 二阶项，误差是否小必须比较被忽略项与二阶项，而不能比较裸环境 Hamiltonian 中某个非线性系数。把 cumulant 按方差标准化便得到正文\eqnrefs{eq:non_gaussian_standardized_cumulant_control}。当 $\mathcal C_2^{(\Phi)}$ 很小时，这些标准化比值可能变大，所以“噪声总体很弱”与“噪声近似 Gaussian”是两个不同命题。
\end{derivation}
\begin{derivation}[从\eqnrefs{eq:noise_dissipation_kernels_general}到\eqnrefs{eq:gaussian_feynman_vernon_influence}]
环境为 Gaussian 态且系统--环境耦合对环境变量线性时，连通累积量在二阶精确终止。闭时路径的四个支路必须显式重组，才能看出噪声结构 $A_\Delta A_\Delta$ 与响应结构 $A_\Delta A_\Sigma$ 分别从何而来。

从零均值环境算符 $\delta B$ 的轮廓平均出发：
\begin{equation*}
\mathscr I_{\rm FV}
=\left\langle\mathcal T_{\mathcal C}
\exp\!\left[-\frac{\ii}{\hbar}
\int_{\mathcal C}\dd z\,A_{\mathcal C}(z)\delta B(z)\right]
\right\rangle_B.
\end{equation*}
高斯累积量定理给
\begin{equation*}
\ln\mathscr I_{\rm FV}
=-\frac{1}{2\hbar^2}
\int_{\mathcal C}\dd z\int_{\mathcal C}\dd z'\,
A_{\mathcal C}(z)
C_{\mathcal C}(z,z')
A_{\mathcal C}(z'),
\end{equation*}
其中
\begin{equation*}
C_{\mathcal C}(z,z')
\equiv\langle\mathcal T_{\mathcal C}\delta B(z)\delta B(z')\rangle.
\end{equation*}

闭时路径有向前的 $+$ 支和向后的 $-$ 支。把轮廓积分改写成普通实时间积分时，支路方向带来
\begin{equation*}
s_+=+1,
\qquad
s_-=-1,
\end{equation*}
所以
\begin{equation*}
\ln\mathscr I_{\rm FV}
=-\frac{1}{2\hbar^2}
\sum_{a,b=\pm}s_a s_b
\int_{t_i}^{t_f}\dd t\int_{t_i}^{t_f}\dd s\,
A_a(t)C^{ab}(t,s)A_b(s).
\end{equation*}
四个实时间相关函数不是四个独立对象，而都由 greater/lesser 两个函数组成：
\begin{align*}
C^{++}(t,s)
&=\Theta(t-s)C^>(t,s)+\Theta(s-t)C^<(t,s),\\
C^{--}(t,s)
&=\Theta(s-t)C^>(t,s)+\Theta(t-s)C^<(t,s),\\
C^{+-}(t,s)&=C^<(t,s),\\
C^{-+}(t,s)&=C^>(t,s).
\end{align*}
这里 $++$ 是普通时间有序，$--$ 是反时间有序，而 $+-$ 与 $-+$ 的次序由闭时路径本身固定。

先只看三角域 $t>s$。四支路之和的被积式为
\begin{align*}
\mathcal Q(t,s)
={}&A_+(t)C^>(t,s)A_+(s)
-A_+(t)C^<(t,s)A_-(s)\\
&-A_-(t)C^>(t,s)A_+(s)
+A_-(t)C^<(t,s)A_-(s).
\end{align*}
另一半三角域 $s>t$ 交换积分哑变量 $t\leftrightarrow s$ 后给出完全相同的 $\mathcal Q(t,s)$。因此原来正方形积分前的 $1/2$ 与两个三角域的因子 2 恰好抵消：
\begin{equation*}
\ln\mathscr I_{\rm FV}
=-\frac{1}{\hbar^2}
\int_{t_i}^{t_f}\dd t\int_{t_i}^{t}\dd s\,
\mathcal Q(t,s).
\end{equation*}

现在才引入正文的和、差变量
\begin{equation*}
A_+=A_\Sigma+\frac{A_\Delta}{2},
\qquad
A_-=A_\Sigma-\frac{A_\Delta}{2},
\end{equation*}
以及
\begin{equation*}
C^>=C_{\rm sym}+\ii C_{\rm asym},
\qquad
C^<=C_{\rm sym}-\ii C_{\rm asym}.
\end{equation*}
把这四个等式逐项代入 $\mathcal Q$。先收集 $C_{\rm sym}$：
\begin{align*}
\mathcal Q_{\rm sym}
={}&C_{\rm sym}
\bigl[A_+(t)-A_-(t)\bigr]
\bigl[A_+(s)-A_-(s)\bigr]\\
={}&A_\Delta(t)C_{\rm sym}(t,s)A_\Delta(s).
\end{align*}
再收集 $\ii C_{\rm asym}$：
\begin{align*}
\mathcal Q_{\rm asym}
={}&\ii C_{\rm asym}
\bigl[A_+(t)-A_-(t)\bigr]
\bigl[A_+(s)+A_-(s)\bigr]\\
={}&2\ii A_\Delta(t)C_{\rm asym}(t,s)A_\Sigma(s).
\end{align*}
因此
\begin{equation*}
\mathcal Q(t,s)
=A_\Delta(t)
\left[
C_{\rm sym}(t,s)A_\Delta(s)
+2\ii C_{\rm asym}(t,s)A_\Sigma(s)
\right].
\end{equation*}
代回三角域积分，得到
\begin{equation*}
\boxed{
\mathscr I_{\rm FV}[A_+,A_-]
=\exp\!\left\{-\frac{1}{\hbar^2}
\int_{t_i}^{t_f}\dd t\int_{t_i}^{t}\dd s\,
A_\Delta(t)
\left[
C_{\rm sym}(t,s)A_\Delta(s)
+2\ii C_{\rm asym}(t,s)A_\Sigma(s)
\right]
\right\},}
\end{equation*}
即正文\eqnrefs{eq:gaussian_feynman_vernon_influence}。
\end{derivation}

这个代数同时给出两种核的物理角色。若两条系统历史完全相同，$A_\Delta=0$，影响泛函等于 1；对称核总是乘历史之差，因此控制历史可区分性与退相干，而反对称核只在“差变量 $\times$ 平均变量”中出现，因而编码因果反作用而不是独立的随机噪声强度。
\begin{derivation}[从\eqnrefs{eq:bath_greater_lesser_spectra}到\eqnrefs{eq:fdt_from_defined_correlators}]

正文\eqnrefs{eq:bath_greater_lesser_spectra} 已把环境的正、负频率信息分别编码在 $S_B^>$ 与 $S_B^<$ 中。热平衡额外提供的内容是 Gibbs 权重，因此先在能量本征基中比较两种谱的权重，再把它们的和、差分别连接到 $\widetilde C_{\rm sym}$ 与推迟响应。令 $H_B|n\rangle=E_n|n\rangle$。greater 谱为
\begin{equation*}
S_B^>(\omega)
=\frac{2\pi}{Z_B}\sum_{n,m}\ee^{-\beta_TE_n}|B_{nm}|^2
\delta\!\left(\omega-\frac{E_m-E_n}{\hbar}\right).
\end{equation*}
lesser 谱写成
\begin{equation*}
S_B^<(\omega)
=\frac{2\pi}{Z_B}\sum_{n,m}\ee^{-\beta_TE_n}|B_{nm}|^2
\delta\!\left(\omega+\frac{E_m-E_n}{\hbar}\right).
\end{equation*}
交换哑指标 $n\leftrightarrow m$ 后，
\begin{align*}
S_B^<(\omega)
&=\frac{2\pi}{Z_B}\sum_{n,m}\ee^{-\beta_TE_m}|B_{nm}|^2
\delta\!\left(\omega-\frac{E_m-E_n}{\hbar}\right)\\
&=\ee^{-\beta_T\hbar\omega}S_B^>(\omega),
\end{align*}
因为 $\delta$ 分布只在 $E_m-E_n=\hbar\omega$ 上有支撑。这就是频域 KMS/详细平衡关系。于是
\begin{align*}
\widetilde C_{\rm sym}(\omega)
&=\frac12\big[S_B^>(\omega)+S_B^<(\omega)\big]\\
&=\frac12S_B^>(\omega)\big(1+\ee^{-\beta_T\hbar\omega}\big),\\
\mathscr S_B^{(-)}(\omega)
&=\frac1\hbar\big[S_B^>(\omega)-S_B^<(\omega)\big]\\
&=\frac1\hbar S_B^>(\omega)\big(1-\ee^{-\beta_T\hbar\omega}\big).
\end{align*}
消去 $S_B^>$ 得
\begin{align*}
\frac{\widetilde C_{\rm sym}(\omega)}{\mathscr S_B^{(-)}(\omega)}
&=\frac{\hbar}{2}
\frac{1+\ee^{-\beta_T\hbar\omega}}{1-\ee^{-\beta_T\hbar\omega}}\\
&=\frac{\hbar}{2}\coth\!\left(\frac{\beta_T\hbar\omega}{2}\right).
\end{align*}
另一方面，由
\begin{equation*}
\chi_B^R(t)=-\frac{\ii}{\hbar}\Theta(t)\langle[B(t),B(0)]\rangle
\end{equation*}
的 Fourier 变换可得谱函数与推迟响应的关系
\begin{equation*}
\mathscr S_B^{(-)}(\omega)=-2\operatorname{Im}\chi_B^R(\omega).
\end{equation*}
代入上一比值，得到
\begin{equation*}
\widetilde C_{\rm sym}(\omega)
=-\hbar\coth\!\left(\frac{\beta_T\hbar\omega}{2}\right)
\operatorname{Im}\chi_B^R(\omega),
\end{equation*}
即\eqnrefs{eq:fdt_from_defined_correlators}。只有再取 $\beta_T\hbar\omega\ll1$，才会进一步得到经典高温形式。
\end{derivation}
\begin{derivation}[从\eqnrefs{eq:gaussian_feynman_vernon_influence}到\eqnrefs{eq:generalized_langevin_from_influence}]
取线性耦合 $A=x_S$。要从影响泛函得到运动方程，必须先把它的两个核分成“噪声权重”和“因果反作用”两部分，而不能直接把指数读成 Langevin 方程。

首先利用 $C_{\rm sym}(t,s)=C_{\rm sym}(s,t)$。正文影响泛函中的对称核积分只覆盖 $t>s$ 的三角域，因此
\begin{align*}
 &\int_{t_i}^{t_f}\dd t\int_{t_i}^{t}\dd s\,
 x_\Delta(t)C_{\rm sym}(t,s)x_\Delta(s)\\
 &\qquad=\frac12
 \int_{t_i}^{t_f}\dd t\int_{t_i}^{t_f}\dd s\,
 x_\Delta(t)C_{\rm sym}(t,s)x_\Delta(s).
\end{align*}
于是这一部分恰好是高斯特征泛函
\begin{align*}
 &\exp\!\left[-\frac{1}{2\hbar^2}
 \iint\dd t\dd s\,
 x_\Delta(t)C_{\rm sym}(t,s)x_\Delta(s)\right]\\
 &\qquad=
 \int\mathcal D\xi\,\mathcal P[\xi]
 \exp\!\left[\frac{\ii}{\hbar}
 \int\dd t\,\xi(t)x_\Delta(t)\right],
\end{align*}
其中 $\mathcal P[\xi]$ 是零均值高斯测度，其协方差由
\begin{equation*}
 \langle\xi(t)\xi(s)\rangle_\xi=C_{\rm sym}(t,s)
\end{equation*}
唯一固定。这里的 $\xi$ 首先只是 Hubbard--Stratonovich 辅助场；只有在相应半经典极限中，才把它解释成一条经典随机力轨迹。

再处理反对称核。由正文定义
\begin{equation*}
 \chi_B^R(t,s)
 =\frac{2}{\hbar}\Theta(t-s)C_{\rm asym}(t,s),
\end{equation*}
所以影响泛函中的相位项精确改写为
\begin{align*}
 &-\frac{2\ii}{\hbar^2}
 \int_{t>s}\dd t\dd s\,
 x_\Delta(t)C_{\rm asym}(t,s)x_\Sigma(s)\\
 &\qquad=-\frac{\ii}{\hbar}
 \int_{t_i}^{t_f}\dd t\,x_\Delta(t)
 \int_{t_i}^{t}\dd s\,
 \chi_B^R(t,s)x_\Sigma(s).
\end{align*}
阶跃函数使积分只依赖 $s\le t$ 的过去历史，这就是“记忆力是因果的”的数学来源。

系统自身的作用量取
\begin{equation*}
 S_S[x]=\int\dd t\left[\frac{M}{2}\dot x^2-V(x)\right],
\end{equation*}
并写 $x_\pm=x_\Sigma\pm x_\Delta/2$，则动能差为
\begin{align*}
 \frac{M}{2}(\dot x_+^2-\dot x_-^2)
 &=M\dot x_\Sigma\dot x_\Delta\\
 &=-M\ddot x_\Sigma x_\Delta
 +\frac{\dd}{\dd t}(M\dot x_\Sigma x_\Delta),
\end{align*}
而势能差的 Taylor 展开给
\begin{equation*}
 V(x_+)-V(x_-)
 =V'(x_\Sigma)x_\Delta+O(x_\Delta^3).
\end{equation*}
在两端固定路径差 $x_\Delta(t_i)=x_\Delta(t_f)=0$ 时，边界项消失，因此
\begin{equation*}
 S_S[x_+]-S_S[x_-]
 =\int\dd t\,x_\Delta(t)
 [-M\ddot x_\Sigma(t)-V'(x_\Sigma(t))]
 +O(x_\Delta^3).
\end{equation*}

把系统作用量、耗散相位和辅助高斯场合并，总有效指数中关于 $x_\Delta$ 的线性项为
\begin{align*}
 \frac{\ii}{\hbar}\int\dd t\,x_\Delta(t)
 \Bigg[&-M\ddot x_\Sigma(t)-V'(x_\Sigma(t))\\
 &-\int_{t_i}^{t}\dd s\,\chi_B^R(t,s)x_\Sigma(s)
 +\xi(t)\Bigg].
\end{align*}
粗粒化的半经典驻值方程就是令这个线性系数对每个 $t$ 都为零。最后在物理路径上令 $x_+=x_-=x_S$，即 $x_\Delta=0$、$x_\Sigma=x_S$，得到
\begin{equation*}
 M\ddot x_S(t)+V'[x_S(t)]
 +\int_{t_i}^{t}\dd s\,\chi_B^R(t,s)x_S(s)
 =\xi(t),
\end{equation*}
并同时保留
\begin{equation*}
 \langle\xi(t)\xi(s)\rangle=C_{\rm sym}(t,s).
\end{equation*}
这才是正文\eqnrefs{eq:generalized_langevin_from_influence}。推导清楚地区分了三个层次：$C_{\rm asym}$ 通过推迟核生成确定性的历史反作用；$C_{\rm sym}$ 通过高斯特征泛函生成涨落；把辅助场 $\xi$ 当作经典随机力还需要额外的半经典解释，而不是影响泛函本身自动给出的本体论结论。
\end{derivation}
\subsection{退相干与扩散}

环境对电子的影响可以从密度矩阵的两类分量看出。若主要改变不同通道的相对相位，非对角元会减小，干涉对比度随之降低；若还交换能量，对角元也会随时间变化，能谱可能逐渐展宽。实验中的到达时间波动也能使多次记录平均后的干涉减弱，但单次电子态未必因此变混。本节分别计算这些效应，并比较它们对相干、布居和能宽的影响。

先看实验参考相位的不确定性。若一发电子相对参考光场有随机延迟 $\delta t$，能量格相干 $\rho_{\ell m}$ 获得相位 $\exp[-\ii(\ell-m)\omega\delta t]$。当 $\delta t$ 是方差 $\sigma_t^2$ 的零均值高斯变量时，对未记录延迟作系综平均得到
\begin{equation}
\boxed{
\rho_{\ell m}\longrightarrow
\rho_{\ell m}
\exp\!\left[-\frac12(\ell-m)^2\omega^2\sigma_t^2\right].}
\label{eq:arrival_jitter_dephasing}
\end{equation}
单次电子并未在传播中必然与量子环境纠缠；损失的是实验者对光学相位参考的知识。因此改变定时诊断或后选择原则上可以恢复这类对比度，而真正不可逆的环境偏迹一般不能这样撤销。

为了描述能量格上的动态随机相位，引入无量纲格点指标算符
\[
\hat\Lambda_e\equiv\sum_{\ell\in\mathbb Z}\ell|\ell\rangle\langle\ell|.
\]
令频率噪声 $\mathfrak n_\varphi(t)$ 以
\[
H(t)=H_e+\hbar\mathfrak n_\varphi(t)\hat\Lambda_e
\]
耦合，其中 $\langle\mathfrak n_\varphi\rangle=0$，$\mathfrak n_\varphi$ 的 SI 单位为 s$^{-1}$，平稳相关函数为
\[
C_\varphi(\tau)
\equiv\langle\mathfrak n_\varphi(t+\tau)\mathfrak n_\varphi(t)\rangle.
\]
若 $H_e$ 在所选能量格基底中对角，则 $[H_e,\hat\Lambda_e]=0$，高斯相位平均可精确完成；这一对易性本身不要求相邻能级等间隔。定义累积相位方差的一半
\[
\Xi_\varphi(t)
\equiv\int_0^t\dd\tau\,(t-\tau)C_\varphi(\tau),
\]
则
\begin{equation}
\boxed{
\rho_{\ell m}(t)
=\rho_{\ell m}(0)\ee^{-\ii\omega_{\ell m}t}
\exp[-(\ell-m)^2\Xi_\varphi(t)].}
\label{eq:exact_gaussian_dephasing_electron}
\end{equation}
格点距离 $|\ell-m|$ 越大，相干携带的相对光学相位阶数越高，因此同一相位噪声对远副对角元抑制更强。

\eqnrefs{eq:exact_gaussian_dephasing_electron} 同时区分准静态与短记忆两个极限。若实验窗内可取 $C_\varphi(\tau)\simeq C_\varphi(0)$，则
\[
\Xi_\varphi(t)\simeq\frac12C_\varphi(0)t^2,
\]
相干呈高斯时间衰减；若环境关联时间 $\tau_B$ 远短于系统时间，定义
\[
\Gamma_\varphi\equiv2\int_0^\infty\dd\tau\,C_\varphi(\tau),
\]
则 $\Xi_\varphi(t)\simeq\Gamma_\varphi t/2$，并得到
\begin{equation}
\boxed{
\dot\rho_e
=-\frac{\ii}{\hbar}[H_e,\rho_e]
+\Gamma_\varphi\mathcal D[\Lambda_e]\rho_e,
\qquad
\mathcal D[X]\rho=X\rho X^\dagger-\frac12\{X^\dagger X,\rho\}.}
\label{eq:energy_lattice_pure_dephasing_gksl}
\end{equation}
因此 $\mathcal D[\Lambda_e]$ 由能量格相位噪声的短相关时间极限固定。若环境必须量子化，同一结构应从 $\Lambda_e\otimes B$ 的零频噪声谱经 Born--Markov--久期极限得到；经典随机噪声与量子环境偏迹对应不同的物理假设。

\begin{figure}
\centering
\includegraphics[width=0.72\textwidth]{dephasing_coherence.png}
\caption{能量格相位噪声对不同副对角线的选择性抑制。相隔 $|\ell-m|$ 越远的相干项，对同一相位噪声越敏感。}
\label{fig:dephasing_coherence}
\end{figure}

纯退相干只消除相位，不改变理想能量格布居。若环境能够真实吸收或发射一个交换量子，应使用跃迁通道。远离正能谱下界时，可在双边有效格上定义幺正平移算符 $\hat b|\ell\rangle=|\ell-1\rangle$，并取
\[
J_\uparrow=\sqrt{\Gamma_\uparrow}\,\hat b^\dagger,
\qquad
J_\downarrow=\sqrt{\Gamma_\downarrow}\,\hat b.
\]
对 $P_\ell=\langle\ell|\rho|\ell\rangle$ 取 GKSL 方程的对角元，得到生灭主方程
\begin{equation}
\boxed{
\dot P_\ell
=\Gamma_\uparrow P_{\ell-1}
+\Gamma_\downarrow P_{\ell+1}
-(\Gamma_\uparrow+\Gamma_\downarrow)P_\ell.}
\label{eq:bilateral_energy_random_walk_master}
\end{equation}
其漂移和方差增长率为
\begin{equation}
\boxed{
\frac{\dd\langle\ell\rangle}{\dd t}
=\Gamma_\uparrow-\Gamma_\downarrow,
\qquad
\frac{\dd\operatorname{Var}(\ell)}{\dd t}
=\Gamma_\uparrow+\Gamma_\downarrow.}
\label{eq:energy_random_walk_moments}
\end{equation}
因此相干哈密顿量跃迁与非相干量子跳跃的实验标志不同：前者的早期展宽由振幅干涉决定，后者在 Markov 区域产生线性方差增长。

双边平移只是在远离真实谱底时的有效描述。若使用带最低格点 $\ell_{\min}$ 的单边平移算符 $s$，则向下跳跃在边界被阻塞，平均漂移变为
\[
\frac{\dd\langle\ell\rangle}{\dd t}
=\Gamma_\uparrow
-\Gamma_\downarrow[1-P_{\ell_{\min}}(t)].
\]
只有整个观测窗内 $P_{\ell_{\min}}\ll1$，才可以把\eqnrefs{eq:bilateral_energy_random_walk_master} 当成平移不变的一般动力学方程。

若反冲曲率已先把一条电子跃迁与其余格点谱隔离，才可以把上述开放动力学继续投影到有限维共振子空间。以反冲选出的 Jaynes--Cummings 子空间为例，在 Born、Markov、久期以及谱投影条件都满足后，可写
\begin{equation}
\boxed{
\dot\rho
=-\frac{\ii}{\hbar}[H_{\rm JC}+H_{\rm LS},\rho]
+\Gamma_{\rm cav}\mathcal D[a]\rho
+\Gamma_1\mathcal D[\sigma_-]\rho
+\frac{\Gamma_\varphi}{2}\mathcal D[\sigma_z]\rho.}
\label{eq:open_jc_effective_master_equation}
\end{equation}
电子相干的总衰减率是 $\Gamma_2=\Gamma_1/2+\Gamma_\varphi$，于是该\emph{低维有效模型内部}可定义协同性
\begin{equation}
\boxed{
\mathscr C_{\rm coop}\equiv\frac{2|g_c|^2}{\Gamma_{\rm cav}\Gamma_2}
=\frac{4|g_c|^2}{\Gamma_{\rm cav}(\Gamma_1+2\Gamma_\varphi)}.}
\label{eq:open_jc_cooperativity}
\end{equation}
它是耦合强度与腔线宽、电子相干线宽乘积构成的无量纲协同性。在坏腔绝热消去极限，它等于腔介导电子振幅衰减率与 $\Gamma_2$ 的比值；是否出现可振荡的正常模分裂还要独立检查 $g_c$ 与两条振幅线宽之差，因而协同性本身不能替代强耦合、单模谱隔离、反冲/Feshbach 泄漏或非 Markov 记忆判据。

协同性与正常模振荡判据描述不同问题：前者比较坏腔消元后的有效通道与电子退相干，后者由完整 $2\times2$ 非厄米本征值决定。

\begin{derivation}[从\eqnrefs{eq:arrival_jitter_dephasing}到\eqnrefs{eq:exact_gaussian_dephasing_electron}]
先处理 shot-to-shot 到达时间抖动。若第 $\ell$ 个能量格相对于第 $m$ 个能量格多积累相位
\[
\Delta\phi_{\ell m}=(\ell-m)\omega\,\delta t,
\]
则一次给定 $\delta t$ 的实验把密度矩阵元乘以
\[
\ee^{-\ii\Delta\phi_{\ell m}}
=\ee^{-\ii(\ell-m)\omega\delta t}.
\]
对零均值 Gaussian 分布
\[
p(\delta t)=\frac1{\sqrt{2\pi}\sigma_t}
\exp\!\left(-\frac{\delta t^2}{2\sigma_t^2}\right)
\]
求实验平均。令 $k=(\ell-m)\omega$，配方得到
\begin{align*}
&-\frac{\delta t^2}{2\sigma_t^2}-\ii k\delta t\\
&\qquad=-\frac{(\delta t+\ii k\sigma_t^2)^2}{2\sigma_t^2}
-\frac12k^2\sigma_t^2.
\end{align*}
Gaussian 积分归一化因子不变，因此
\begin{align*}
\left\langle\ee^{-\ii k\delta t}\right\rangle
&=\int\dd\delta t\,p(\delta t)\ee^{-\ii k\delta t}\\
&=\exp\!\left(-\frac12k^2\sigma_t^2\right)\\
&=\exp\!\left[-\frac12(\ell-m)^2\omega^2\sigma_t^2\right],
\end{align*}
即正文\eqnrefs{eq:arrival_jitter_dephasing}。

再处理随时间变化的零均值 Gaussian 频率噪声 $\mathfrak n_\varphi(t)$。定义累积随机相位
\[
X_t=\int_0^t\dd s\,\mathfrak n_\varphi(s).
\]
给定一条噪声轨迹时，能量格密度矩阵元为
\[
\rho_{\ell m}(t|\mathfrak n_\varphi)
=\rho_{\ell m}(0)\ee^{-\ii\omega_{\ell m}t}
\ee^{-\ii(\ell-m)X_t}.
\]
因为 $X_t$ 是 Gaussian 过程的线性泛函，它自身仍是零均值 Gaussian 变量，所以其特征函数满足
\[
\left\langle\ee^{-\ii qX_t}\right\rangle
=\exp\!\left[-\frac{q^2}{2}\langle X_t^2\rangle\right].
\]
取 $q=\ell-m$。由平稳相关函数
\[
C_\varphi(s-s')
=\langle\mathfrak n_\varphi(s)\mathfrak n_\varphi(s')\rangle
\]
有
\begin{align*}
\langle X_t^2\rangle
&=\int_0^t\dd s\int_0^t\dd s'\,C_\varphi(s-s')\\
&=2\int_0^t\dd s\int_0^s\dd s'\,C_\varphi(s-s')\\
&=2\int_0^t\dd s\int_0^s\dd\tau\,C_\varphi(\tau)\\
&=2\int_0^t\dd\tau\,(t-\tau)C_\varphi(\tau).
\end{align*}
其中第二步利用实平稳噪声相关函数的偶对称性，第三步取 $\tau=s-s'$，最后交换积分次序。因而
\[
\frac12\langle X_t^2\rangle
=\int_0^t\dd\tau\,(t-\tau)C_\varphi(\tau)
\equiv\Xi_\varphi(t).
\]
对噪声平均后得到
\[
\rho_{\ell m}(t)
=\rho_{\ell m}(0)\ee^{-\ii\omega_{\ell m}t}
\exp[-(\ell-m)^2\Xi_\varphi(t)],
\]
即正文\eqnrefs{eq:exact_gaussian_dephasing_electron}。
\end{derivation}
\begin{derivation}[从\eqnrefs{eq:exact_gaussian_dephasing_electron}到\eqnrefs{eq:energy_lattice_pure_dephasing_gksl}]
设 $C_\varphi(\tau)$ 绝对可积且一阶绝对矩有限，并定义
\begin{equation*}
I_0\equiv\int_0^\infty\dd\tau\,C_\varphi(\tau),
\qquad
M_1^{(\varphi)}\equiv\int_0^\infty\dd\tau\,\tau|C_\varphi(\tau)|.
\end{equation*}
相位耦合算符 $\Lambda_e=\sum_\ell\ell|\ell\rangle\langle\ell|$ 在能量格上对角。对精确退相干指数，直接加减 $tI_0$ 得到恒等式
\begin{align*}
\Xi_\varphi(t)
&=\int_0^t\dd\tau\,(t-\tau)C_\varphi(\tau)\\
&=tI_0-r_\varphi(t),\\
r_\varphi(t)
&\equiv
t\int_t^\infty\dd\tau\,C_\varphi(\tau)
+\int_0^t\dd\tau\,\tau C_\varphi(\tau).
\end{align*}
余项满足严格界
\begin{equation*}
|r_\varphi(t)|
\le
t\int_t^\infty\dd\tau\,|C_\varphi(\tau)|
+M_1^{(\varphi)}.
\end{equation*}
当右端相对于 $t|I_0|$ 很小时，才可用 $\Xi_\varphi(t)=tI_0+r_{\rm M}(t)$ 且 $|r_{\rm M}(t)|\ll t|I_0|$ 的时间局域极限。
代入正文关系
\begin{equation*}
\frac{\Gamma_\varphi}{2}=\int_0^\infty\dd\tau\,C_\varphi(\tau),
\end{equation*}
于是主导项为 $\Xi_\varphi(t)=\Gamma_\varphi t/2$。再计算 Lindblad 耗散子。由于
\begin{equation*}
\Lambda_e|\ell\rangle=\ell|\ell\rangle,
\end{equation*}
有
\begin{align*}
\langle\ell|\Lambda_e\rho\Lambda_e|m\rangle&=\ell m\rho_{\ell m},\\
\langle\ell|\Lambda_e^2\rho|m\rangle&=\ell^2\rho_{\ell m},\\
\langle\ell|\rho\Lambda_e^2|m\rangle&=m^2\rho_{\ell m}.
\end{align*}
因此
\begin{align*}
\langle\ell|\mathcal D[\Lambda_e]\rho|m\rangle
&=\ell m\rho_{\ell m}-\frac12(\ell^2+m^2)\rho_{\ell m}\\
&=-\frac12(\ell-m)^2\rho_{\ell m}.
\end{align*}
对角元 $\ell=m$ 的右端为零，所以人口保持不变；非对角元满足
\begin{equation*}
\dot\rho_{\ell m}=-\frac{\Gamma_\varphi}{2}(\ell-m)^2\rho_{\ell m},
\end{equation*}
解为
\begin{equation*}
\rho_{\ell m}(t)=\rho_{\ell m}(0)\exp[-\Gamma_\varphi t(\ell-m)^2/2],
\end{equation*}
与精确高斯退相干律的短记忆极限一致，即\eqnrefs{eq:energy_lattice_pure_dephasing_gksl}。
\end{derivation}
\begin{derivation}[从\eqnrefs{eq:bilateral_energy_random_walk_master}到\eqnrefs{eq:energy_random_walk_moments}]
由最近邻 Lindblad 跳跃算符推出双向随机游走主方程~\eqnrefs{eq:bilateral_energy_random_walk_master}，再对 $\ell$ 与 $\ell^2$ 求和得到矩方程~\eqnrefs{eq:energy_random_walk_moments}。 这里跃迁率时间局域且在所取格点范围内平移不变； 边界效应可忽略或单独处理。 双边位移满足 $\hat b^\dagger\hat b=\hat b\hat b^\dagger=I$。 因此

\[
\langle\ell|\mathcal D[\hat b^\dagger]\rho|\ell\rangle
=P_{\ell-1}-P_\ell,
\qquad
\langle\ell|\mathcal D[\hat b]\rho|\ell\rangle
=P_{\ell+1}-P_\ell,
\]
线性组合即得布居主方程。乘以 $\ell$ 后求和，并分别令 $j=\ell-1$、$j=\ell+1$，有
\[
\sum_\ell\ell P_{\ell-1}=\langle\ell\rangle+1,
\qquad
\sum_\ell\ell P_{\ell+1}=\langle\ell\rangle-1,
\]
从而得到平均漂移。二阶矩需要分别计算两条平移支路。对第一条支路令 $j=\ell-1$，于是 $\ell=j+1$：
\begin{align*}
\sum_\ell\ell^2P_{\ell-1}
&=\sum_j(j+1)^2P_j\\
&=\sum_j(j^2+2j+1)P_j\\
&=\langle\ell^2\rangle+2\langle\ell\rangle+1.
\end{align*}
对第二条支路令 $j=\ell+1$，于是 $\ell=j-1$：
\begin{align*}
\sum_\ell\ell^2P_{\ell+1}
&=\sum_j(j-1)^2P_j\\
&=\langle\ell^2\rangle-2\langle\ell\rangle+1.
\end{align*}
设向上、向下随机跃迁率分别为 $\Gamma_+$ 与 $\Gamma_-$。代入布居主方程后
\begin{align*}
\frac{\dd}{\dd t}\langle\ell^2\rangle
&=\Gamma_+\left(2\langle\ell\rangle+1\right)
+\Gamma_-\left(-2\langle\ell\rangle+1\right)\\
&=2(\Gamma_+-\Gamma_-)\langle\ell\rangle
+(\Gamma_++\Gamma_-).
\end{align*}
另一方面
\[
\frac{\dd}{\dd t}\langle\ell\rangle
=\Gamma_+-\Gamma_-.
\]
因此方差 $\operatorname{Var}(\ell)=\langle\ell^2\rangle-\langle\ell\rangle^2$ 满足
\begin{align*}
\frac{\dd}{\dd t}\operatorname{Var}(\ell)
&=\frac{\dd}{\dd t}\langle\ell^2\rangle
-2\langle\ell\rangle
\frac{\dd}{\dd t}\langle\ell\rangle\\
&=\Gamma_++\Gamma_-.
\end{align*}
\end{derivation}
平均能量格漂移由两方向跃迁率之差控制，而方差增长率由两者之和控制。
\begin{derivation}[从\eqnrefs{eq:open_jc_effective_master_equation}到\eqnrefs{eq:open_jc_cooperativity}]
先从主方程读取电子相干线宽。对 $\rho_{ge}=\langle g|\rho|e\rangle$，布居弛豫耗散子给
\begin{align*}
\langle g|\mathcal D[\sigma_-]\rho|e\rangle
&=\langle g|\sigma_-\rho\sigma_+|e\rangle
-\frac12\langle g|\{\sigma_+\sigma_-,\rho\}|e\rangle\\
&=-\frac12\rho_{ge},
\end{align*}
而
\begin{equation*}
\langle g|\mathcal D[\sigma_z]\rho|e\rangle
=\langle g|\sigma_z\rho\sigma_z-\rho|e\rangle
=-2\rho_{ge}.
\end{equation*}
所以
\begin{equation*}
\dot\rho_{ge}\big|_{\rm diss}
=-\left(\frac{\Gamma_1}{2}+\Gamma_\varphi\right)\rho_{ge}
\equiv-\Gamma_2\rho_{ge}.
\end{equation*}

在线性单激发子空间 $\{|e,0\rangle,|g,1\rangle\}$ 中，零失谐时的均值方程为
\begin{equation*}
\frac{\dd}{\dd t}
\begin{pmatrix}c_e\\c_c\end{pmatrix}
=
\underbrace{\begin{pmatrix}
-\Gamma_2&-\ii g_c\\
-\ii g_c^*&-\Gamma_{\rm cav}/2
\end{pmatrix}}_{\mathsf M}
\begin{pmatrix}c_e\\c_c\end{pmatrix}.
\end{equation*}
矩阵 $\mathsf M$ 的两个本征值由二次方程精确给出
\begin{equation*}
\lambda_\pm
=-\frac{\Gamma_2+\Gamma_{\rm cav}/2}{2}
\pm
\sqrt{
\frac{(\Gamma_2-\Gamma_{\rm cav}/2)^2}{4}
-|g_c|^2}.
\end{equation*}
因此时间域出现振荡型正常模分裂的条件是根号变成纯虚数，即
\begin{equation*}
|g_c|>
\frac12\left|\Gamma_2-\frac{\Gamma_{\rm cav}}{2}\right|.
\end{equation*}
这是一条关于本征值结构的判据；实验上要把两条谱峰清楚分辨，还需分裂大于相应总线宽。

另一方面，在坏腔极限中可从第二行写出完全等价的恒等式
\begin{equation*}
c_c
=-\frac{2\ii g_c^*}{\Gamma_{\rm cav}}c_e
-\frac{2}{\Gamma_{\rm cav}}\dot c_c.
\end{equation*}
把它代回第一行，不丢项时得到
\begin{equation*}
\dot c_e
=-\left(\Gamma_2+\frac{2|g_c|^2}{\Gamma_{\rm cav}}\right)c_e
+\frac{2\ii g_c}{\Gamma_{\rm cav}}\dot c_c.
\end{equation*}
当快变量条件使最后一项相对于前两项可忽略时，腔介导的电子振幅衰减率为
\begin{equation*}
\Gamma_{\rm cav\to e}^{({\rm amp})}
\equiv\frac{2|g_c|^2}{\Gamma_{\rm cav}}.
\end{equation*}
将它与电子自身相干衰减率比较，严格得到正文采用的协同性定义
\begin{equation*}
\frac{\Gamma_{\rm cav\to e}^{({\rm amp})}}{\Gamma_2}
=\frac{2|g_c|^2}{\Gamma_{\rm cav}\Gamma_2}
\equiv\mathscr C_{\rm coop}.
\end{equation*}
\end{derivation}
\subsection{有限相关时间噪声：可见度、回波与近似误差}
\label{sec:electron-colored-noise-applications}

如果环境的随机频率偏移在一段时间内仍彼此相关，电子先后积累的相位也会相关。这时用一个常数衰减率描述全部演化，可能低估或高估相干。下面从一般相位噪声的关联函数出发，计算自由积累与理想回波，并用一个可解析的相关函数比较它们。

沿用前文无量纲能量格算符 $\hat\Lambda_e$ 和频率噪声 $\mathfrak n_\varphi(t)$。假定后者是均值为零的平稳经典高斯过程，并且所研究的相位演化阶段满足 $[H_e,\hat\Lambda_e]=0$。外加控制用已知实函数 $y_\varphi(t)$ 表示，噪声耦合为 $\hbar y_\varphi(t)\mathfrak n_\varphi(t)\hat\Lambda_e$。$y_\varphi=1$ 对应自由相位积累；理想反号控制使 $y_\varphi$ 在指定时刻改变符号。这个反号是对选定相干子空间的控制假设，在连续电子能谱上实现它需要另行指定器件与脉冲，不能由符号变换直接推断实验可行性。

对能量格相干 $\rho_{\ell m}$，定义阶差 $d_{\ell m}=\ell-m$。去掉已知自由传播相位后，剩余相干比为
\begin{equation}
 \mathcal V_{\ell m}(t)
 =\exp[-d_{\ell m}^2\chi_\varphi[y_\varphi;t]],\qquad
 \chi_\varphi[y_\varphi;t]
 =\frac12\int_0^t\dd s\int_0^t\dd s'\,
 y_\varphi(s)y_\varphi(s')C_\varphi(s-s').
 \label{eq:electron-noise-filter-time}
\end{equation}
$\mathcal V_{\ell m}$ 是单个密度矩阵元的幅值衰减因子；实际干涉可见度若混合多个阶差，需先将各项按测量算符组合。$C_\varphi$ 的单位为 $\mathrm{s^{-2}}$，故 $\chi_\varphi$ 无量纲。高斯性与对易性使该式在这个噪声模型内精确，无需弱噪声假设。非高斯过程需要高阶累积量，不对易耦合则需要时间排序的算符核。

按照全书 Fourier 约定，定义双边频率噪声谱与控制滤波振幅
\begin{equation*}
 S_\varphi(\omega)=\int_{-\infty}^{\infty}\dd\tau\,
 C_\varphi(\tau)\ee^{\ii\omega\tau},\qquad
 F_{y}(\omega,t)=\int_0^t\dd s\,y_\varphi(s)\ee^{\ii\omega s}.
\end{equation*}
$S_\varphi$ 的单位为 $\mathrm{s^{-1}}$，$F_y$ 为时间，频域形式为
\begin{equation}
 \chi_\varphi[y_\varphi;t]
 =\frac12\int_{-\infty}^{\infty}\frac{\dd\omega}{2\pi}
 S_\varphi(\omega)|F_y(\omega,t)|^2.
 \label{eq:electron-noise-filter-frequency}
\end{equation}
平稳经典噪声的谱非负，所以 $\chi_\varphi\geq0$，单条相干不会超出其初值。此式的因子 $1/2$ 来自高斯特征函数，$1/(2\pi)$ 来自双边 Fourier 测度；若改用单边谱，必须同时修改测度与谱归一化。

现在代入具有指数相关函数的高斯噪声，令 $\sigma_\varphi$ 为频率噪声的均方根、$\tau_\varphi$ 为相关时间：
\begin{equation}
 C_\varphi(\tau)=\sigma_\varphi^2\ee^{-|\tau|/\tau_\varphi},\qquad
 S_\varphi(\omega)=\frac{2\sigma_\varphi^2\tau_\varphi}{1+\omega^2\tau_\varphi^2}.
 \label{eq:electron-exponential-noise}
\end{equation}
这就是 Ornstein--Uhlenbeck 平稳噪声的相关函数。对自由积累 $y_\varphi=1$，时间积分给
\begin{equation}
 \chi_{\rm free}(t)
 =\sigma_\varphi^2\tau_\varphi^2
 \left[\frac{t}{\tau_\varphi}-1+\ee^{-t/\tau_\varphi}\right].
 \label{eq:electron-exponential-dephasing}
\end{equation}
它在短时呈二次增长，长时才近似线性。将线性部分写成前文 $\Gamma_\varphi t/2$，得到 $\Gamma_\varphi=2\sigma_\varphi^2\tau_\varphi$。常数率近似相对于精确相干的比值为
\begin{equation}
 \frac{\mathcal V_{\ell m}^{\rm M}}{\mathcal V_{\ell m}^{\rm exact}}
 =\exp\!\left[-d_{\ell m}^2\sigma_\varphi^2\tau_\varphi^2
 (1-\ee^{-t/\tau_\varphi})\right].
 \label{eq:electron-markov-visibility-error}
\end{equation}
因此 $t\gg\tau_\varphi$ 保证瞬时斜率接近常数，却不自动保证从 $t=0$ 初始化的相干幅值准确。若要求该比值与一之差小于给定误差，还需控制 $d_{\ell m}^2\sigma_\varphi^2\tau_\varphi^2$。高阶边带相干对同一记忆时间更敏感；这把一般 Markov 时间尺度判据落实到具体可观测量。

若在 $t/2$ 作理想反号，使前半段 $y_\varphi=1$、后半段 $y_\varphi=-1$，则
\begin{equation}
 \chi_{\rm echo}(t)=\sigma_\varphi^2\tau_\varphi^2
 \left[\frac{t}{\tau_\varphi}-3
 +4\ee^{-t/(2\tau_\varphi)}-\ee^{-t/\tau_\varphi}\right].
 \label{eq:electron-exponential-echo}
\end{equation}
准静态相位在两半程互相抵消，短时领先项从 $t^2$ 变成 $t^3$。有限相关时间内新产生的涨落仍然保留，所以回波不等于关闭环境。有限脉冲宽度、反号误差或反号期间发生的能量交换需回到真实 $y_\varphi$ 与联合动力学。

\begin{exbox}[不同阶相干对记忆误差的敏感性]
设相邻边带的相干衰减因子为 $v(t)=\exp[-\chi_{\rm free}(t)]$。求相隔 $d$ 阶的相干，以及常数率近似误差不超过 $0<\varepsilon<1$ 时允许的阶差。

由\eqnrefs{eq:electron-noise-filter-time}，相隔 $d$ 阶的因子为 $v(t)^{d^2}$。因此，即使相邻边带的相干只略有减小，更远的非对角元也可能明显衰减。记无量纲时间 $x=t/\tau_\varphi>0$、噪声强度 $a_\varphi=\sigma_\varphi\tau_\varphi>0$，则常数率近似相对于精确结果的比值为
\[
R_d=\exp[-d^2a_\varphi^2(1-\ee^{-x})].
\]
要求 $1-R_d\le\varepsilon$，对两边取对数，得到
\[
|d|\le\frac{\sqrt{-\ln(1-\varepsilon)}}
{a_\varphi\sqrt{1-\ee^{-x}}}.
\]
长时间 $x\gg1$ 时，右边趋于有限值 $\sqrt{-\ln(1-\varepsilon)}/a_\varphi$。因此延长观测时间虽使瞬时衰减率接近常数，却不会消除最初记忆阶段积累的相对幅值误差。

对理想回波，直接相减两个指数得
\[
\chi_{\rm free}-\chi_{\rm echo}
=2a_\varphi^2(1-\ee^{-x/2})^2,
\qquad
\frac{\mathcal V_d^{\rm echo}}{\mathcal V_d^{\rm free}}
=\exp[2d^2a_\varphi^2(1-\ee^{-x/2})^2].
\]
回波提高保留下来的相干，但两者在长时间仍共同衰减。若另有独立的到达时间抖动，还需乘上\eqnrefs{eq:arrival_jitter_dephasing}的衰减因子；若两种噪声相关，则必须保留交叉协方差。
\end{exbox}

\begin{figure}[htbp]
 \centering
 \includegraphics[width=0.75\textwidth]{electron_colored_noise.pdf}
 \caption{指数相关噪声中的自由相干、理想回波和常数率近似，取 $\sigma_\varphi\tau_\varphi=0.20$。相隔三个边带的相干比相邻边带更早区分不同近似；回波主要抑制低频相位噪声。}
 \label{fig:electron-colored-noise}
\end{figure}
\figref{fig:electron-colored-noise} 将近似误差和实际信号放在同一尺度上。判定某个近似足够好时，应使用所测阶差和曝光时间，而非只依据环境相关时间的绝对长短。

\begin{derivation}[高斯平均与自由积累闭式]
在相互作用绘景中，各时刻的噪声哈密顿量彼此对易，单次实现的传播算符为
\begin{equation*}
 U_I(t)=\exp\!\left[-\ii\hat\Lambda_e\int_0^t\dd s\,
 y_\varphi(s)\mathfrak n_\varphi(s)\right].
\end{equation*}
令随机相位 $\Theta_y=\int_0^t\dd s\,y_\varphi(s)\mathfrak n_\varphi(s)$，则
$\rho_{\ell m,I}(t)=\rho_{\ell m,I}(0)\langle\ee^{-\ii d_{\ell m}\Theta_y}\rangle$。
由于 $\Theta_y$ 是零均值高斯变量，其特征函数为 $\exp[-d_{\ell m}^2\langle\Theta_y^2\rangle/2]$，直接得到\eqnrefs{eq:electron-noise-filter-time}。将 $C_\varphi(\tau)=\int\dd\omega\,S_\varphi(\omega)\ee^{-\ii\omega\tau}/(2\pi)$ 代入双积分，两段时间积分互为复共轭，得到频域模平方。

自由积累时，相关函数为偶函数，所以
\begin{align*}
 \chi_{\rm free}(t)
 &=\int_0^t\dd\tau\,(t-\tau)C_\varphi(\tau)\\
 &=\sigma_\varphi^2\left[t\tau_\varphi(1-\ee^{-t/\tau_\varphi})
 -\tau_\varphi^2\{1-(1+t/\tau_\varphi)\ee^{-t/\tau_\varphi}\}\right]\\
 &=\sigma_\varphi^2\tau_\varphi^2(t/\tau_\varphi-1+\ee^{-t/\tau_\varphi}).
\end{align*}
短时展开给 $\chi_{\rm free}=\sigma_\varphi^2t^2/2-\sigma_\varphi^2t^3/(6\tau_\varphi)+\order(t^4)$；这同时检验初始值和初始斜率都为零。对白噪声极限，必须保持 $\Gamma_\varphi=2\sigma_\varphi^2\tau_\varphi$ 不变并令 $\tau_\varphi\to0$；若反而固定 $\sigma_\varphi$，噪声谱的总低频强度趋零，不会得到有限退相干率。
\end{derivation}

\begin{derivation}[回波的跨区间相关项]
令半程时间 $a_t=t/2$，并令 $\Theta_1,\Theta_2$ 为两半程积累相位。反号后总相位为 $\Theta_1-\Theta_2$，故
\begin{equation*}
 \chi_{\rm echo}=\frac12\operatorname{Var}(\Theta_1-\Theta_2)
 =2\chi_{\rm free}(a_t)-\langle\Theta_1\Theta_2\rangle.
\end{equation*}
跨区间协方差不能丢弃，因为回波正是利用它消去缓慢变化的噪声。指数相关函数给
\begin{align*}
 \langle\Theta_1\Theta_2\rangle
 &=\sigma_\varphi^2\int_0^{a_t}\dd s\int_{a_t}^{2a_t}\dd s'\,
 \ee^{-(s'-s)/\tau_\varphi}\\
 &=\sigma_\varphi^2\tau_\varphi^2(1-\ee^{-a_t/\tau_\varphi})^2.
\end{align*}
代回即得\eqnrefs{eq:electron-exponential-echo}。其短时展开为
$\chi_{\rm echo}=\sigma_\varphi^2t^3/(12\tau_\varphi)+\order(t^4)$，长时自由积累和回波具有相同斜率 $\sigma_\varphi^2\tau_\varphi$。频域中 $F_y(0,t)=0$ 表明回波消除严格静态的噪声分量，而有限频率的权重仍非零。
\end{derivation}


\subsection{时间局域生成元与非 Markov 动力学}

时间卷积和时间局域是同一约化动力学的两种表示方式。Nakajima--Zwanzig 形式把被环境保存的信息显式写入历史卷积；time-convolutionless（TCL）形式则把这些历史效应重新编码进随时间变化的局域生成元。若从初态到时刻 $t$ 的约化动力学映射 $\Phi_{t,t_0}$ 在所研究算符空间上可逆，就可以直接定义这一精确的 TCL 生成元。令
\begin{equation}
\boxed{
\rho_S(t)=\Phi_{t,t_0}[\rho_S(t_0)],
\qquad
\mathcal K_{\rm TCL}(t)
\equiv
\dot\Phi_{t,t_0}\,\Phi_{t,t_0}^{-1},}
\label{eq:exact_tcl_generator_definition}
\end{equation}
则恒等地有
\begin{equation}
\boxed{
\dot\rho_S(t)
=\mathcal K_{\rm TCL}(t)[\rho_S(t)].}
\label{eq:exact_tcl_master_equation}
\end{equation}
TCL 是 time-convolutionless 的缩写。上式表明，生成元在当前时刻只作用于 $\rho_S(t)$，但它的系数已经包含此前系统--环境相关积累下来的全部历史信息。若 $\Phi_{t,t_0}$ 在某时刻失去可逆性，$\mathcal K_{\rm TCL}$ 可以出现奇点，而动力学映射 $\Phi_{t,t_0}$ 仍可保持有限；只有再满足短环境记忆等附加条件，时间依赖生成元才可进一步约化为通常的 Markov 生成元。

对常用的因子化投影
$\mathcal P X=\operatorname{Tr}_B(X)\otimes\rho_B$，并在相互作用绘景中满足
$\mathcal P\mathcal L_I(t)\mathcal P=0$ 时，TCL 的二阶生成元为
\begin{equation}
\boxed{
\mathcal K_{\rm TCL}^{(2)}(t)
=\int_{t_0}^{t}\dd s\,
\mathcal P\mathcal L_I(t)\mathcal L_I(s)\mathcal P,}
\label{eq:tcl2_generator_projected_liouvillian}
\end{equation}
其中 $\mathcal L_I X\equiv-\ii[H_I/\hbar,X]$。它与 Born--Redfield 历史核具有相同的二阶环境相关函数，但作用在当前态 $\rho_S(t)$ 上，而不是 $\rho_S(s)$。只有进一步把积分上限延长到无穷并用快速衰减相关函数替代有限历史，才进入 Markov 极限；再做 Bohr 频率久期化后，生成元才具有标准 GKSL 结构。

因此需要区分三个逻辑层级。第一，时间局域只是一种表示方式，不等于无记忆；第二，某一时刻的生成元能写成 Lindblad 形状，也不意味着所有瞬时速率都非负；第三，在有限维工作窗中，若时间局域标准形的所有耗散率始终非负，则动力学具有 CP-divisibility，这是比“方程没有记忆积分”更强的 Markov 性条件。自由电子能量格若与窄带腔模或结构共振发生可逆信息回流，TCL 速率可以暂时变负，此时仍应把它理解为非 Markov 约化动力学，而不是负概率。

对于实验拟合，这个区分很重要。若只需要传播一个已知初态，NZ 记忆核和 TCL 时间局域生成元是同一约化动力学的两种表示；若要把主方程参数解释成真正的常数弛豫率，则必须额外验证环境相关时间、系统 Bohr 频率分离以及 CP-divisibility，而不能仅因数值上成功拟合一个时间局域方程就称其为 Markov。


单时刻动力学映射还有一个更深的限制。$\Phi_{t,t_0}$ 只回答“给定初态以后，在时刻 $t$ 得到什么状态”；如果实验在 $t_1,t_2,\ldots$ 之间主动插入电子碰撞、脉冲、测量或后选择，而环境在这些时刻之间保留记忆，那么末态一般不能由同一组两时刻映射简单复合得到。自由电子量子光学尤其容易进入这一情形：一列电子可以先后穿过同一个长寿命腔模、声子模或材料极化，前一个电子在环境中留下的量子信息会改变后一个电子看到的输入状态。独立碰撞模型只有在这部分记忆已经被重置或不可分辨时才成立\cite{Ciccarello2022CollisionModels,Pollock2018OperationalMarkov,Keeling2026ProcessTensor}。

为描述真正的多时刻实验，把显式保留的系统记为 $S$，把所有跨时刻保存信息的自由度合并成\emph{记忆空间} $M$。取 $t_0<t_1<\cdots<t_N$；在每个 $t_k$ 允许实验者对 $S$ 施加任意完全正操作 $\mathcal A_k$；它可以表示一次已知电子碰撞、一个控制脉冲、一个带结果的量子仪器分支，或恒等操作。相邻时刻之间，$S+M$ 由 CPTP 映射 $\mathcal U_{k+1:k}$ 演化。于是末态由多线性映射
\begin{equation}
\boxed{
\begin{aligned}
\rho_S(t_N)
={}&\mathfrak T_{N:0}
[\mathcal A_{N-1},\ldots,\mathcal A_0]\\
\equiv{}&\operatorname{Tr}_M\!\Big[
\mathcal U_{N:N-1}(\mathcal A_{N-1}\otimes\mathcal I_M)
\cdots
\mathcal U_{1:0}(\mathcal A_0\otimes\mathcal I_M)
\rho_{SM}(t_0)
\Big]
\end{aligned}}
\label{eq:process_tensor_multitime_definition}
\end{equation}
给出。这个把一整串中间操作映射到末态及其联合概率的多线性对象称为\emph{过程张量}（process tensor）。普通密度算符回答“在一个时刻测什么会得到什么概率”，普通量子信道回答“给定一次输入操作后状态怎样改变”；过程张量则把问题提升为“在多个时刻依次插入不同操作时，整条记录的概率怎样改变”。因此它具有类似多时刻状态的作用，但并不是一个没有额外约束的普通密度算符：时间槽之间还必须满足因果归一化条件。

为了把这个多线性对象变成可直接做矩阵运算的对象，使用 Choi 表示。线性量子操作虽然以算符为输入，但任意输入都可在线性基 $|i\rangle\langle j|$ 上展开，所以只要保存操作对这些基矩阵的全部输出，就已经保存了整个操作。把输入基指标与输出算符一起排成一个更大的矩阵，得到 Choi 算符。对
$\mathcal A:L(\mathcal H^{\rm in})\to L(\mathcal H^{\rm out})$，定义
\[
A
=\sum_{ij}|i\rangle\langle j|_{\rm in}
\otimes
\mathcal A(|i\rangle\langle j|)_{\rm out}.
\]
它作用在输入算符 $X$ 上的恢复公式为
\[
\mathcal A(X)
=\operatorname{Tr}_{\rm in}
\left[A\,(X^{\mathsf T}\otimes\mathbb I_{\rm out})\right].
\]
反复使用这条恒等式，就存在一个正半定多时刻 Choi 算符
$\Upsilon_{N:0}\succeq0$，使
\begin{equation}
\boxed{
\rho_S(t_N)
=\Upsilon_{N:0}
\star A_{N-1}\star\cdots\star A_0,}
\label{eq:process_tensor_choi_link}
\end{equation}
其中 $\star$ 称为 \emph{link product}（连接积）：它不是新的动力学，而是把一个操作的输出指标与下一个操作的输入指标收缩掉的记号；具体运算就是按上一式在共享输入空间取转置并偏迹。这样，多时刻实验网络的“把线接起来”变成普通矩阵指标收缩。过程张量的因果归一化条件保证：若所有 $\mathcal A_k$ 都是保迹信道，则最终态仍归一；若某个 $\mathcal A_k$ 是量子仪器中的一个非保迹分支，则最终迹直接给该整条实验记录的联合概率。

这一定义也精确说明“Markov 碰撞模型”究竟额外假设了什么。若每个时间步都使用一个与过去无关的新鲜辅助系统，而且旧辅助系统在下一步以前永久丢弃，则记忆空间在每一步都被重置，过程张量的 Choi 算符沿时间方向因子化为独立的一步信道；末态退化为
\[
\rho_S(t_N)
=\mathscr E_{N:N-1}\circ\mathcal A_{N-1}\circ\cdots
\circ\mathscr E_{1:0}\circ\mathcal A_0[\rho_S(t_0)].
\]
若进一步每一周期完全相同，才可以写成前文频闪映射的幂。反之，只要同一个腔模、材料极化或声子在相邻电子之间没有重置，$S$ 与 $M$ 会在第一次碰撞后相关；第二次碰撞的结果不仅取决于当前 $\rho_S$，还取决于过去操作怎样改变了 $M$。这时即使每个固定初态的两时刻动力学都能拟合得很好，也不能据此推出真正的多时刻 Markov 性。

对自由电子实验，这个层级给出一个直接判据。若环境记忆时间 $\tau_M$ 远短于相邻电子间隔 $T_e$，且单次碰撞后残余系统--环境关联在 $T_e$ 内衰减到实验分辨率以下，可以使用独立碰撞/Poisson 信道；若 $\tau_M\gtrsim T_e$，则应把相关环境模式并入显式系统，或直接保留过程张量。因而“电子逐个到达”并不等于“电子作用彼此独立”：真正需要比较的是电子到达间隔与环境保存量子信息的时间尺度。


过程张量在形式上随时间槽数指数增长，因此还需要回答一个计算问题：如果环境记忆只有有限复杂度，能否不用存储整个 $\Upsilon_{N:0}$？这里使用的 Schmidt 分解与双粒子纯态的 Schmidt 分解是同一个线性代数思想，只是现在把\emph{算符本身}视为 Hilbert--Schmidt 空间中的向量。选定一个时间切割以后，左侧指标代表过去，右侧指标代表未来；需要多少个独立乘积项才能连接两侧，就量化了跨过这个切割必须保存多少记忆。把过程 Choi 算符在第 $k$ 个时间槽处划分为“过去”与“未来”两部分，
\[
\mathcal H_{\rm proc}
=
\mathcal H_{\le k}\otimes\mathcal H_{>k},
\]
则总可以写成
\begin{equation}
\boxed{
\Upsilon_{N:0}
=
\sum_{r=1}^{R_k}
s_r^{(k)}
A_r^{(\le k)}
\otimes
B_r^{(>k)},}
\label{eq:process_tensor_temporal_schmidt}
\end{equation}
其中 $s_r^{(k)}\ge0$ 是该时间切割的算符 Schmidt 奇异值，$R_k$ 是非零奇异值个数，即 Schmidt rank。若所有切割的 rank 都受到同一个常数 $D_M$ 控制，就可以把这些逐切割分解串成 matrix-product operator（MPO，矩阵乘积算符）：每个时间槽只保留一个局域张量，相邻张量用内部 \emph{bond} 指标连接。最大 bond dimension $D_M$ 就是跨任意时间切割需要携带的内部算符记忆维数，因此它比单一“相关时间”更细致地刻画记忆复杂度。

这个结构对显式系统--记忆模型有直接上界。若所有历史信息都能压缩进一个有限维记忆空间
$\mathcal H_M$，$\dim\mathcal H_M=d_M$，那么在 Liouville 空间中跨任意时间切割传递的记忆指标至多有 $d_M^2$ 个，因此
\begin{equation}
\boxed{
R_k\le d_M^2,
\qquad
D_M\le d_M^2.}
\label{eq:process_tensor_memory_bond_bound}
\end{equation}
例如只需显式保留一个截断到 $n=0,\ldots,N_M$ 的长寿命腔模时，$d_M=N_M+1$；过程张量虽然跨越很多电子时刻，其时间方向的精确 bond dimension 仍由这个记忆 Hilbert 空间而不是电子总数决定。

若 $R_k$ 很大但奇异值快速衰减，可以保留前 $D$ 个 Schmidt 分量。对 Hilbert--Schmidt 范数，单个时间切割的最优截断误差严格为
\begin{equation}
\boxed{
\epsilon_{k,D}^{(2)}
=
\left[
\sum_{r>D}
\bigl(s_r^{(k)}\bigr)^2
\right]^{1/2}.}
\label{eq:process_tensor_schmidt_truncation_error}
\end{equation}
这只是算符表示误差，不应自动解释成任意实验概率的统一 diamond-norm 上界；实际可观测误差还要与允许操作集的范数相乘。不过它提供了一个可计算判据：若所有实验相关时间切割都存在小的 $\epsilon_{k,D}^{(2)}$，则有限记忆 MPO 是受控近似，而不是经验性的“只记前几步”。

这个低维结构也改变层析的可辨识条件。若因果和低秩参数化把允许过程限制在 $d_{\rm eff}$ 维局部参数流形，并用矩阵 $B$ 把有效参数 $\boldsymbol\vartheta$ 映射到前文线性系数 $\boldsymbol\theta$，
\[
\delta\boldsymbol\theta=B\,\delta\boldsymbol\vartheta,
\]
则实验只需要在这一切空间上满秩：
\begin{equation}
\boxed{
\operatorname{rank}(XB)=d_{\rm eff},
\qquad
\mathcal F_{\rm eff}
=
B^{\mathsf T}X^{\mathsf T}WXB.}
\label{eq:process_tensor_lowrank_identifiability}
\end{equation}
因此有限记忆不是“可以少测一些”的模糊说法，而是把待识别方向从完整过程空间压缩到一个明确的低维切空间。若某个低秩模型拟合很好但 $XB$ 仍有零模，实验仍不能唯一判断对应的环境记忆参数。

低秩表示还必须回答一个更直接的实验问题：MPO 截断误差究竟会把电子计数概率改多少？固定实验设置 $r$ 和完整记录 $y$ 后，所有“已知的东西”——初始制备、各时刻控制操作以及最终 POVM——可以按与过程张量互补的方式连接成一个正算符。这个算符称为 \emph{过程测试算符}；它在多时刻理论中承担的角色与单时刻 Born 规则中的 POVM 元相同。记为 $T_{r,y}$，则
\begin{equation}
\boxed{p_{r,y}=\Tr[T_{r,y}\Upsilon_{N:0}].}
\label{eq:process_tester_probability_linear_functional}
\end{equation}
若用近似过程 $\widetilde\Upsilon$ 代替真过程，Hilbert--Schmidt Cauchy--Schwarz 给
\begin{equation}
\boxed{
|p_{r,y}-\widetilde p_{r,y}|
\le\|T_{r,y}\|_2\,
\|\Upsilon-\widetilde\Upsilon\|_2.}
\label{eq:process_probability_hs_error_bound}
\end{equation}
若连续作若干次 SVD 截断，并把第 $k$ 次对当前 MPO 的实际改变量记为 $\epsilon_k$，三角不等式给
\begin{equation}
\boxed{
\|\Upsilon-\widetilde\Upsilon\|_2\le\sum_k\epsilon_k,
\qquad
\frac12\sum_y|p_{r,y}-\widetilde p_{r,y}|
\le
\frac12\left(\sum_y\|T_{r,y}\|_2\right)\sum_k\epsilon_k.}
\label{eq:process_mpo_total_variation_error_bound}
\end{equation}
同样的 MPO 表示误差，对不同测量设计不一定产生同样大的实验误差；真正相关的是过程测试算符对被舍弃时间相关方向的灵敏度。

还要注意，逐切割 SVD 只优化算符距离，并不自动保持 $\widetilde\Upsilon\succeq0$ 和过程张量的递归因果约束。因此实际数值压缩应把低 bond-dimension MPO 当作参数化流形，在压缩后重新施加正性和因果归一化，或直接在这些约束内变分拟合。否则即使算符误差很小，也可能出现负概率或未来影响过去的非物理解。


该界控制的是指定过程测试算符族上的可观测误差；若需要对任意辅助系统和控制操作给最坏情形界，则必须改用更强的 operational process norm。

\begin{derivation}[从\eqnrefs{eq:exact_tcl_generator_definition}到\eqnrefs{eq:tcl2_generator_projected_liouvillian}]
先证明精确时间局域式。对
\[
\rho_S(t)=\Phi_t\rho_S(t_0)
\]
求导：
\[
\dot\rho_S(t)=\dot\Phi_t\rho_S(t_0).
\]
若 $\Phi_t$ 在所考察算符工作窗上可逆，则
\[
\rho_S(t_0)=\Phi_t^{-1}\rho_S(t),
\]
因此
\[
\dot\rho_S(t)
=\dot\Phi_t\Phi_t^{-1}\rho_S(t)
\equiv\mathcal K_{\rm TCL}(t)\rho_S(t),
\]
得到正文\eqnrefs{eq:exact_tcl_master_equation}。因此 $\mathcal K_{\rm TCL}(t)=\dot\Phi_t\Phi_t^{-1}$ 的定义域只由 $\Phi_t$ 的可逆工作窗决定；环境记忆随后全部进入这个时间依赖生成元。

二阶约化动力学可直接从联合 Liouville 方程迭代一次得到。相互作用绘景中
\[
\dot\rho_{SB}(t)=\mathcal L_I(t)\rho_{SB}(t),
\qquad
\mathcal L_I(t)X=-\frac{\ii}{\hbar}[H_I(t),X].
\]
形式积分一次：
\[
\rho_{SB}(t)
=\rho_{SB}(t_0)
+\int_{t_0}^{t}\dd s\,\mathcal L_I(s)\rho_{SB}(s).
\]
再代回微分方程并投影，得到到二阶为止
\begin{align*}
\frac{\dd}{\dd t}\mathcal P\rho_{SB}(t)
&=\mathcal P\mathcal L_I(t)\mathcal P\rho_{SB}(t_0)\\
&\quad+\int_{t_0}^{t}\dd s\,
\mathcal P\mathcal L_I(t)\mathcal L_I(s)
\mathcal P\rho_{SB}(t_0)
+O(H_I^3).
\end{align*}
采用因子化投影并令环境平均相互作用已经吸收到 $H_S$ 中，使
$\mathcal P\mathcal L_I\mathcal P=0$，一阶项消失。为了把右端改写成当前态而不是初态，注意
\[
\mathcal P\rho_{SB}(t)
=\mathcal P\rho_{SB}(t_0)+O(H_I),
\]
而二阶积分前面已经含两个 $\mathcal L_I$；用当前投影态替换初态只会产生 $O(H_I^3)$ 误差。因此
\[
\frac{\dd}{\dd t}\mathcal P\rho_{SB}(t)
=\int_{t_0}^{t}\dd s\,
\mathcal P\mathcal L_I(t)\mathcal L_I(s)\mathcal P
\rho_{SB}(t)+O(H_I^3),
\]
即正文\eqnrefs{eq:tcl2_generator_projected_liouvillian}。

若 $H_I(t)=\sum_\alpha A_\alpha(t)\otimes B_\alpha(t)$，对环境取迹后，TCL2 系数显式由有限时间相关函数
\[
C_{\alpha\beta}(t,s)
=\operatorname{Tr}_B[B_\alpha(t)B_\beta(s)\rho_B]
\]
组成。此时积分仍只到当前时刻 $t$，所以环境记忆仍保留在时间依赖系数中。只有当正文\eqnrefs{eq:markov_timescale_control_parameter} 的历史替换上界与相关尾上界都很小，才可把有限历史核替换为时间局域、无限上限的系数；再经过 Bohr 频率分块与久期化，才得到正文前面使用的 Born--Markov--GKSL 链。
\end{derivation}

\begin{derivation}[从\eqnrefs{eq:process_tensor_multitime_definition}到\eqnrefs{eq:process_tensor_choi_link}]
先证明过程张量对每一个实验操作都是线性的。以两个中间时刻为例，
\[
\rho_S(t_2)
=\operatorname{Tr}_M\!
\left[
\mathcal U_{2:1}(\mathcal A_1\otimes\mathcal I_M)
\mathcal U_{1:0}(\mathcal A_0\otimes\mathcal I_M)
\rho_{SM}(t_0)
\right].
\]
若
$\mathcal A_0=\lambda\mathcal A_0^{(1)}+\mu\mathcal A_0^{(2)}$，则逐项代入定义有
\begin{align*}
\mathfrak T_{2:0}[\mathcal A_1,\mathcal A_0]
={}&\operatorname{Tr}_M\!\left[
\mathcal U_{2:1}(\mathcal A_1\otimes\mathcal I_M)
\mathcal U_{1:0}((\lambda\mathcal A_0^{(1)}+\mu\mathcal A_0^{(2)})\otimes\mathcal I_M)
\rho_{SM}(t_0)\right]\\
={}&\lambda\mathfrak T_{2:0}[\mathcal A_1,\mathcal A_0^{(1)}]
+\mu\mathfrak T_{2:0}[\mathcal A_1,\mathcal A_0^{(2)}].
\end{align*}
对每一个槽位同理，因此 $\mathfrak T_{N:0}$ 是多线性的高阶量子映射。

为把量子操作转化成普通正半定算符，取输入基 $\{|i\rangle\}$ 并定义 Choi 算符
\[
A
=\sum_{ij}|i\rangle\langle j|\otimes
\mathcal A(|i\rangle\langle j|).
\]
任意输入算符写成
$X=\sum_{mn}X_{mn}|m\rangle\langle n|$。于是
\begin{align*}
\operatorname{Tr}_{\rm in}
\left[A(X^{\mathsf T}\otimes\mathbb I)\right]
&=\sum_{ijmn}X_{nm}
\operatorname{Tr}
\left[|i\rangle\langle j|m\rangle\langle n|\right]
\mathcal A(|i\rangle\langle j|)\\
&=\sum_{ij}X_{ij}\mathcal A(|i\rangle\langle j|)\\
&=\mathcal A(X).
\end{align*}
因此一个映射的所有输入--输出作用都被压缩进一个正算符 $A$ 中。

现在把这条恢复公式逐次代入正文\eqnrefs{eq:process_tensor_multitime_definition}。每个 $\mathcal A_k$ 都引入一对“输入--输出” Hilbert 空间；联合传播与初态中不属于实验操作的全部张量指标可以预先收缩成一个固定算符 $\Upsilon_{N:0}$。剩余的每个操作只通过
\[
\operatorname{Tr}_{k^{\rm in}}
\left[A_k(X^{\mathsf T}\otimes\mathbb I)\right]
\]
与它相连。把“在共享空间取转置并偏迹”定义成 link product $\star$，便得到
\[
\rho_S(t_N)
=\Upsilon_{N:0}\star A_{N-1}\star\cdots\star A_0,
\]
即正文\eqnrefs{eq:process_tensor_choi_link}。

$\Upsilon_{N:0}$ 的正性来自完整实验网络是“初态 + CPTP 联合传播”的 Choi 算符偏迹：对任意一组完全正操作，其 link product 都必须给出正的未归一末态。因果归一化可逐时间向后检验。若最后一个操作取任意保迹信道，其 Choi 算符满足
\[
\operatorname{Tr}_{\rm out}A_k=\mathbb I_{\rm in},
\]
于是对 $\Upsilon_{N:0}$ 最后一个输出空间取迹必须退化成前一个时刻的过程张量乘输入单位算符。递归执行得到一组线性因果约束；它们正保证未来操作不能改变过去记录的概率。

最后看无记忆极限。若每个区间只与一个新鲜辅助系统 $M_k$ 相互作用，且
\[
\rho_{M_0M_1\cdots}
=\rho_{M_0}\otimes\rho_{M_1}\otimes\cdots,
\]
每个 $M_k$ 在本步结束后就被迹掉，那么第 $k$ 步产生独立 CPTP 映射
$\mathscr E_{k+1:k}$。联合网络中不存在连接非相邻时间片的记忆指标，所以 Choi 网络沿时间方向断开成一步映射的链，过程张量退化为普通信道复合。若同一个记忆 $M$ 被连续多步复用，则内部 $M$ 指标贯穿多个时间片，$\Upsilon_{N:0}$ 一般不能这样因子化；这就是重复电子碰撞中环境记忆的多时刻数学来源。
\end{derivation}

\begin{derivation}[从\eqnrefs{eq:process_tensor_temporal_schmidt}到\eqnrefs{eq:process_tensor_lowrank_identifiability}]
先证明有限记忆空间给出的 bond-dimension 上界。把第 $k$ 个时间切割以前的所有系统操作和联合演化收缩成一个从初始记忆算符到当前记忆算符的超算符。若
\[
\dim\mathcal H_M=d_M,
\]
则记忆算符空间
$L(\mathcal H_M)$ 的维数为 $d_M^2$。任选一组正交算符基
$\{M_\alpha\}_{\alpha=1}^{d_M^2}$，过去网络对未来的全部影响都可展开为
\[
X_M
=
\sum_{\alpha=1}^{d_M^2}
x_\alpha M_\alpha.
\]
因此跨越该时间切割需要传递的独立线性系数不超过 $d_M^2$。把过去所有指标收缩成 $A_\alpha^{(\le k)}$、未来网络收缩成 $B_\alpha^{(>k)}$，便有
\[
\Upsilon_{N:0}
=
\sum_{\alpha=1}^{d_M^2}
A_\alpha^{(\le k)}
\otimes
B_\alpha^{(>k)}.
\]
对这组系数再作奇异值分解，只会降低项数，所以
\[
R_k\le d_M^2.
\]
对每个时间切割重复这一分解，就得到时间方向的 MPO，最大 bond dimension 满足正文\eqnrefs{eq:process_tensor_memory_bond_bound}。

再看截断。把\eqnrefs{eq:process_tensor_temporal_schmidt} 看成 Hilbert--Schmidt 空间中的普通 Schmidt/SVD：
\[
\Upsilon
=
\sum_r s_r A_r\otimes B_r,
\]
并取
\[
\Tr(A_r^\dagger A_{r'})=\delta_{rr'},
\qquad
\Tr(B_r^\dagger B_{r'})=\delta_{rr'}.
\]
保留前 $D$ 项后
\[
\Upsilon_D
=
\sum_{r\le D}s_r A_r\otimes B_r.
\]
正交性给
\begin{align*}
\|\Upsilon-\Upsilon_D\|_2^2
&=
\Tr[(\Upsilon-\Upsilon_D)^\dagger
(\Upsilon-\Upsilon_D)]\\
&=
\sum_{r>D}s_r^2,
\end{align*}
从而得到正文\eqnrefs{eq:process_tensor_schmidt_truncation_error}。Eckart--Young 定理说明，在固定 rank $D$ 的 Hilbert--Schmidt 近似中，这个截断是最优的。

最后考虑层析。前文一般线性模型为
\[
\mathbf p=X\boldsymbol\theta.
\]
若物理约束和低秩 MPO 参数化只允许
\[
\boldsymbol\theta
=
\boldsymbol\theta_0
+
B\boldsymbol\vartheta
+O(\|\boldsymbol\vartheta\|^2),
\]
则局部概率变化为
\[
\delta\mathbf p
=
XB\,\delta\boldsymbol\vartheta.
\]
两个不同的微小参数方向在一阶实验数据中可区分，当且仅当它们不落入
$\ker(XB)$；因此局部可辨识条件正是
\[
\operatorname{rank}(XB)=d_{\rm eff}.
\]
对独立计数数据，若原参数的 Fisher 权矩阵写成 $W$，链式法则给
\[
\mathcal F_{\rm eff}
=
\left(\frac{\partial\mathbf p}{\partial\boldsymbol\vartheta}\right)^{\mathsf T}
W
\left(\frac{\partial\mathbf p}{\partial\boldsymbol\vartheta}\right)
=
B^{\mathsf T}X^{\mathsf T}WXB,
\]
得到正文\eqnrefs{eq:process_tensor_lowrank_identifiability}。
\end{derivation}

\begin{derivation}[从\eqnrefs{eq:process_tester_probability_linear_functional}到\eqnrefs{eq:process_mpo_total_variation_error_bound}]
令
\[
\Delta\Upsilon\equiv\Upsilon-\widetilde\Upsilon.
\]
真实过程与近似过程在同一过程测试算符元 $T_{r,y}$ 上给出的概率之差为
\begin{align*}
\delta p_{r,y}
&=p_{r,y}-\widetilde p_{r,y}\\
&=\Tr(T_{r,y}\Upsilon)-\Tr(T_{r,y}\widetilde\Upsilon)\\
&=\Tr(T_{r,y}\Delta\Upsilon).
\end{align*}
把最后一式看成 Hilbert--Schmidt 内积。由于
\[
\langle A,B\rangle_{\rm HS}=\Tr(A^\dagger B),
\qquad
\|A\|_2^2=\Tr(A^\dagger A),
\]
而过程测试算符元为厄米正算符，$T_{r,y}^\dagger=T_{r,y}$，Cauchy--Schwarz 不等式逐步给出
\begin{align*}
|\delta p_{r,y}|
&=|\langle T_{r,y},\Delta\Upsilon\rangle_{\rm HS}|\\
&\le
\sqrt{\langle T_{r,y},T_{r,y}\rangle_{\rm HS}}
\sqrt{\langle\Delta\Upsilon,\Delta\Upsilon\rangle_{\rm HS}}\\
&=\|T_{r,y}\|_2\|\Delta\Upsilon\|_2,
\end{align*}
得到正文\eqnrefs{eq:process_probability_hs_error_bound}。

再把多次 MPO 压缩的误差写成望远镜和。若
\[
\Upsilon^{(0)}=\Upsilon,
\qquad
\Upsilon^{(K)}=\widetilde\Upsilon,
\]
则严格恒等式为
\begin{align*}
\Delta\Upsilon
&=\Upsilon^{(0)}-\Upsilon^{(K)}\\
&=\sum_{k=1}^{K}
\left(\Upsilon^{(k-1)}-\Upsilon^{(k)}\right).
\end{align*}
定义每一次压缩对\emph{当前} MPO 的实际 Hilbert--Schmidt 改变量
\[
\epsilon_k
\equiv
\|\Upsilon^{(k-1)}-\Upsilon^{(k)}\|_2,
\]
连续使用三角不等式便得到
\begin{align*}
\|\Delta\Upsilon\|_2
&\le
\sum_{k=1}^{K}
\|\Upsilon^{(k-1)}-\Upsilon^{(k)}\|_2\\
&=\sum_{k=1}^{K}\epsilon_k.
\end{align*}
若第 $k$ 次操作确实是在当前切割上作正交 Schmidt/SVD 截断，设该步当前算符的奇异值为 $s_j^{(k,\mathrm{cur})}$，则正交性还给出
\[
\epsilon_k^2
=\sum_{j>D_k}\bigl(s_j^{(k,\mathrm{cur})}\bigr)^2.
\]
这里必须使用每一步\emph{当前} MPO 的尾和；不同切割上对原始过程分别计算的尾和并不组成同一个望远镜分解，因而不能直接相加。

对固定实验设置 $r$，记录 $y$ 互斥。总变差距离定义为
\[
D_{\rm TV}^{(r)}
\equiv
\frac12\sum_y|p_{r,y}-\widetilde p_{r,y}|.
\]
逐项代入单个记录的界，再使用上面的全局 MPO 误差界，得到
\begin{align*}
D_{\rm TV}^{(r)}
&\le\frac12\sum_y
\|T_{r,y}\|_2\|\Delta\Upsilon\|_2\\
&=\frac12
\left(\sum_y\|T_{r,y}\|_2\right)
\|\Delta\Upsilon\|_2\\
&\le\frac12
\left(\sum_y\|T_{r,y}\|_2\right)
\sum_{k=1}^{K}\epsilon_k,
\end{align*}
即正文\eqnrefs{eq:process_mpo_total_variation_error_bound}。
\end{derivation}



\section{量子测量与统计信息}

已经求得探测面上的密度算符后，接下来计算探测器会记录什么。理想能量测量给出某条边带的概率，实际谱仪则把有限能量区间分配给同一个读数，还可能漏掉部分电子。我们先从熟悉的投影测量出发，建立这些读出过程的概率公式，再讨论数据能够确定哪些态分量，以及计数有限时参数能测多准。

\subsection{量子仪器与正算符值测度（POVM）}

回忆投影测量。若一组正交态 $|x\rangle$ 对应可区分的离散结果，投影算符为 $P_x=|x\rangle\langle x|$，结果概率为 $\operatorname{Tr}(P_x\rho_S)$。真实谱仪的一个读数可能来自几个不同能量，因而该结果需要由更一般的正算符 $\Pi_x^{(S)}$ 描述。这里“正”指对任意态 $|\psi\rangle$，都有 $\langle\psi|\Pi_x^{(S)}|\psi\rangle\ge0$。一组和为恒等算符的这种效应算符称为\emph{正算符值测度}（positive-operator-valued measure, POVM）；它们保证各结果概率非负且总和为一。

如果还要知道读到 $x$ 后系统处于什么态，就需要规定状态更新。用线性映射 $\mathcal I_x$ 表示这个更新，但暂不把输出归一化：它的迹保留结果发生的概率。所有结果对应的一族映射称为\emph{量子仪器}（quantum instrument）。每个 $\mathcal I_x$ 必须完全正，即只对更大纠缠系统中的 $S$ 部分作此更新，联合态仍保持非负；它还必须迹不增，因为单个结果的概率不能超过一。把所有结果加起来则保迹。因此
\begin{equation}
\boxed{
p(x)=\operatorname{Tr}[\mathcal I_x(\rho_S)],
\qquad
\rho_{S|x}=\frac{\mathcal I_x(\rho_S)}{p(x)},
\qquad
\sum_x\mathcal I_x\ \text{保迹}.}
\label{eq:quantum_instrument_probability_definition}
\end{equation}
当 $p(x)>0$ 时，除以迹就得到归一化条件态 $\rho_{S|x}$；零概率结果不定义条件态。为了只计算概率，可以把映射从态移到观测算符上。定义\emph{对偶映射} $\mathcal I_x^*$ 满足
\[
\operatorname{Tr}[A\mathcal I_x(\rho)]
=\operatorname{Tr}[\mathcal I_x^*(A)\rho]
\]
对所考虑的任意 $A,\rho$ 成立。取 $A=\mathbb I_S$，就得到效应算符 $\Pi_x^{(S)}=\mathcal I_x^*(\mathbb I_S)$，以及
\begin{equation}
\boxed{
\Pi_x^{(S)}\succeq0,
\qquad
\sum_x\Pi_x^{(S)}=\mathbb I_S,
\qquad
p(x)=\operatorname{Tr}[\Pi_x^{(S)}\rho_S].}
\label{eq:abstract_instrument_to_povm}
\end{equation}
同一个 $\Pi_x^{(S)}$ 可以来自不同的 $\mathcal I_x$，所以相同统计并不唯一决定测量反作用。

抽象的 $\mathcal I_x$ 并不是额外假设出来的“测量公设”。实验上更具体的图像是：待测系统先与一个可控测量仪发生普通量子力学的联合幺正演化，然后只读取测量仪的一部分信息。忽略测量仪对应偏迹，选择某个测量仪结果对应条件化；Kraus 算符和 POVM 都从这两个操作中产生。物理实验常通过“系统 $S$ + 测量仪 $M$ + 联合幺正 + 测量仪读出”实现量子仪器。假设制备时系统与测量仪因子化，$\rho_{SM}^{(i)}=\rho_S\otimes\rho_M$，并把
$\rho_M=\sum_\nu p_\nu|\nu\rangle\langle\nu|$。若联合演化为 $U$，而实验完全忽略测量仪，系统先由联合态的偏迹定义无条件演化
\begin{equation}
\boxed{
\mathscr E(\rho_S)=\operatorname{Tr}_M\!\left[U(\rho_S\otimes\rho_M)U^\dagger\right].}
\label{eq:unitary_dilation_partial_trace_definition}
\end{equation}
对测量仪初态与末态基展开以后，它化为完全正保迹（CPTP）的 Kraus 形式
\begin{equation}
\boxed{
\mathscr E(\rho_S)
=\sum_{\mu,\nu}K_{\mu\nu}\rho_SK_{\mu\nu}^\dagger,
\qquad
K_{\mu\nu}=\sqrt{p_\nu}\langle\mu|U|\nu\rangle,
\qquad
\sum_{\mu,\nu}K_{\mu\nu}^\dagger K_{\mu\nu}=\mathbb I_S.}
\label{eq:general_unitary_dilation_kraus_operator}
\end{equation}
初始因子化是得到统一系统映射的物理条件；若 $S$ 与 $M$ 预先相关，一般不存在作用于任意 $\rho_S$ 的同一个 CPTP 映射，必须保留初始联合态，或改用过程张量/赋值映射描述。

读出测量仪时，令结果 $x$ 对应测量仪侧效应 $\Pi_x^{(M)}\succeq0$ 且 $\sum_x\Pi_x^{(M)}=\mathbb I_M$。对测量仪取偏迹后，系统的未归一化条件态为
\begin{equation}
\boxed{
\mathcal I_x(\rho_S)
=\operatorname{Tr}_M\!\left[
(\mathbb I_S\otimes [\Pi_x^{(M)}]^{1/2})
U(\rho_S\otimes\rho_M)U^\dagger
(\mathbb I_S\otimes [\Pi_x^{(M)}]^{1/2})
\right].}
\label{eq:general_indirect_measurement_instrument}
\end{equation}
由于最终对测量仪取偏迹，两个 $[\Pi_x^{(M)}]^{1/2}$ 也可以利用偏迹的循环性合成一个 $\Pi_x^{(M)}$。对 $\Pi_x^{(M)}$ 作谱分解后，\eqnrefs{eq:general_indirect_measurement_instrument} 化为 Kraus 求和，因此每个分支完全正。系统对应的 POVM 效应算符为
\[
\Pi_x^{(S)}=\mathcal I_x^*(\mathbb I_S),
\]
而不是测量仪上的 $\Pi_x^{(M)}$ 本身。

实验探测器通常还在理想量子结果之后加入有限分辨率和效率。对理想电子边带 $|\ell\rangle$，若能谱仪输出连续读数 $E_d$，定义归一化条件响应密度 $R(E_d|\mathcal E_\ell)$。当仪器只做经典能量模糊而不相干混合边带时，实际 POVM 密度为
\begin{equation}
\boxed{
\Pi_{\rm e}(E_d)
=\sum_\ell R(E_d|\mathcal E_\ell)|\ell\rangle\langle\ell|.}
\label{eq:electron_detector_response_povm}
\end{equation}
于是 $p(E_d)=\operatorname{Tr}[\Pi_{\rm e}(E_d)\rho_e]$。只有当 $R(E_d|\mathcal E_\ell)$ 还具有平移不变形式时，这个映射才退化为通常的卷积；能量依赖效率、角度分辨或相干混合都应直接进入更一般的 POVM，而不是修改相互作用哈密顿量。




\begin{exbox}[弱测量与弱值]
\eqnrefs{eq:general_indirect_measurement_instrument} 已经给出一般“系统--探针--读出”测量。弱测量只是在其中选择一个特定的脉冲耦合并对耦合强度作受控展开。令系统的无量纲厄米可观测量为 $\hat A$，探针具有正则变量
\[
[\hat Q,\hat P]=\ii\hbar,
\]
并取脉冲相互作用
\begin{equation}
\hat H_{\rm int}(t)=\lambda(t)\hat A\otimes\hat P,
\qquad
g\equiv\int\dd t\,\lambda(t),
\qquad [g]={\rm m}.
\label{eq:weak_measurement_von_neumann_coupling}
\end{equation}
忽略脉冲期间的自由演化时，联合幺正为
\[
\hat U_w=\exp\!\left(-\frac{\ii g}{\hbar}\hat A\otimes\hat P\right).
\]
系统预选为 $|i\rangle$，探针初态为 $|\phi\rangle$，随后把系统后选择到 $|f\rangle$。当 $\langle f|i\rangle\neq0$ 时，定义无量纲弱值
\begin{equation}
\boxed{
A_w\equiv\frac{\langle f|\hat A|i\rangle}{\langle f|i\rangle}.}
\label{eq:weak_value_definition}
\end{equation}
若探针动量均方根宽度为 $\sigma_P$，一阶弱展开至少要求
\[
\epsilon_w^{(1)}
\equiv\frac{|g|\sigma_P}{\hbar}|A_w|\ll1,
\qquad
\epsilon_w^{(2)}
\equiv\frac{g^2\sigma_P^2}{2\hbar^2}
\left|\frac{\langle f|\hat A^2|i\rangle}{\langle f|i\rangle}\right|\ll1.
\]
第二个条件显式检查接近正交后选择时被 $\langle f|i\rangle^{-1}$ 放大的二阶项。对 $\langle Q\rangle_i=\langle P\rangle_i=0$、实且无 $Q$--$P$ 相关的最小不确定 Gaussian 探针，
$\sigma_P=\hbar/(2\sigma_Q)$，后选择后的平均读数为
\begin{equation}
\boxed{
\Delta\langle Q\rangle_f
=g\,\operatorname{Re}A_w+O(g^2),
\qquad
\Delta\langle P\rangle_f
=\frac{2g\sigma_P^2}{\hbar}\operatorname{Im}A_w+O(g^2)
=\frac{g\hbar}{2\sigma_Q^2}\operatorname{Im}A_w+O(g^2).}
\label{eq:weak_value_gaussian_meter_shifts}
\end{equation}
因此 $\operatorname{Re}A_w$ 与 $\operatorname{Im}A_w$ 分别控制共轭探针变量的一阶条件位移。$A_w$ 可以落在 $\hat A$ 的本征值区间之外，但它是带后选择条件的响应系数；后选择概率和高阶修正必须与位移同时计算。若 $\epsilon_w^{(1)}$ 或 $\epsilon_w^{(2)}$ 不小，应直接使用\eqnrefs{eq:general_indirect_measurement_instrument} 的完整量子仪器，而不使用弱值截断。


\medskip

后选择后的未归一化探针态由系统偏矩阵元给出：
\[
|\widetilde\phi_f\rangle
=\langle f|\hat U_w|i\rangle|\phi\rangle.
\]
对 $g$ 展开联合幺正，
\begin{align*}
\langle f|\hat U_w|i\rangle
&=\langle f|i\rangle
-\frac{\ii g}{\hbar}\langle f|\hat A|i\rangle\hat P
+O(g^2)\\
&=\langle f|i\rangle
\left(1-\frac{\ii g}{\hbar}A_w\hat P\right)+O(g^2).
\end{align*}
初始探针满足 $\langle P\rangle_i=0$，所以归一化因子在一阶不变。对任意探针算符 $\hat O$，条件平均值的一阶修正为
\[
\delta\langle O\rangle
=\frac{\ii g}{\hbar}
\left(A_w^*\langle P O\rangle_i
-A_w\langle O P\rangle_i\right)+O(g^2).
\]
对实、居中的 Gaussian 波包，
\[
\langle\{Q,P\}\rangle_i=0,
\qquad
\langle QP\rangle_i=\frac{\ii\hbar}{2},
\qquad
\langle PQ\rangle_i=-\frac{\ii\hbar}{2}.
\]
因此取 $O=Q$ 得
\begin{align*}
\delta\langle Q\rangle
&=\frac{\ii g}{\hbar}
\left[-\frac{\ii\hbar}{2}A_w^*
-\frac{\ii\hbar}{2}A_w\right]+O(g^2)\\
&=g\,\operatorname{Re}A_w+O(g^2).
\end{align*}
取 $O=P$ 并使用 $\langle P^2\rangle_i=\sigma_P^2$，
\begin{align*}
\delta\langle P\rangle
&=\frac{\ii g}{\hbar}(A_w^*-A_w)\sigma_P^2+O(g^2)\\
&=\frac{2g\sigma_P^2}{\hbar}\operatorname{Im}A_w+O(g^2).
\end{align*}
最小不确定 Gaussian 满足 $\sigma_Q\sigma_P=\hbar/2$，代入即得到\eqnrefs{eq:weak_value_gaussian_meter_shifts} 的最后一个等号。
\end{exbox}

\begin{derivation}[从\eqnrefs{eq:quantum_instrument_probability_definition}到\eqnrefs{eq:abstract_instrument_to_povm}]
量子仪器的每个结果分支 $\mathcal I_x$ 是完全正、迹不增映射，而无条件映射
\[
\mathscr E=\sum_x\mathcal I_x
\]
保迹。定义对偶映射 $\mathcal I_x^*$ 使任意系统算符 $X$ 与迹类算符 $\rho$ 满足
\[
\Tr[X\,\mathcal I_x(\rho)]
=\Tr[\mathcal I_x^*(X)\rho].
\]
取 $X=\mathbb I_S$，定义
\[
\Pi_x^{(S)}\equiv\mathcal I_x^*(\mathbb I_S).
\]
于是结果概率逐步变为
\begin{align*}
p(x)
&=\Tr[\mathcal I_x(\rho_S)]\\
&=\Tr[\mathbb I_S\mathcal I_x(\rho_S)]\\
&=\Tr[\mathcal I_x^*(\mathbb I_S)\rho_S]\\
&=\Tr[\Pi_x^{(S)}\rho_S].
\end{align*}

再验证 $\Pi_x^{(S)}$ 是正算符。任取 $|\psi\rangle$，令
$\rho_\psi=|\psi\rangle\langle\psi|\succeq0$。完全正性蕴含正性，因此
$\mathcal I_x(\rho_\psi)\succeq0$，从而
\begin{align*}
\langle\psi|\Pi_x^{(S)}|\psi\rangle
&=\Tr[\Pi_x^{(S)}\rho_\psi]\\
&=\Tr[\mathcal I_x(\rho_\psi)]\\
&\ge0.
\end{align*}
该式对任意 $|\psi\rangle$ 成立，所以
\[
\Pi_x^{(S)}\succeq0.
\]
迹不增还给出 $\Pi_x^{(S)}\preceq \mathbb I_S$。因为对任意密度算符
\[
0\le\Tr[\mathcal I_x(\rho)]\le\Tr\rho,
\]
所以
\[
0\le\Tr[(\mathbb I_S-\Pi_x^{(S)})\rho]
\]
对任意 $\rho\succeq0$ 成立，故 $\mathbb I_S-\Pi_x^{(S)}\succeq0$。

最后利用无条件映射保迹：
\begin{align*}
\Tr\rho
&=\Tr[\mathscr E(\rho)]\\
&=\sum_x\Tr[\mathcal I_x(\rho)]\\
&=\Tr\!\left[\left(\sum_x\Pi_x^{(S)}\right)\rho\right].
\end{align*}
因此
\[
\Tr\!\left[\left(\mathbb I_S-\sum_x\Pi_x^{(S)}\right)\rho\right]=0
\]
对任意迹类 $\rho$ 成立，只能有
\[
\sum_x\Pi_x^{(S)}=\mathbb I_S.
\]
于是 $\{\Pi_x^{(S)}\}$ 满足正性、完备性和 Born 概率公式，正是正文\eqnrefs{eq:abstract_instrument_to_povm}。
\end{derivation}

\begin{derivation}[从\eqnrefs{eq:unitary_dilation_partial_trace_definition}到\eqnrefs{eq:general_unitary_dilation_kraus_operator}]
设测量仪初态作谱分解
\[
\rho_M=\sum_\nu p_\nu|\nu\rangle\langle\nu|,
\qquad
p_\nu\ge0,
\qquad
\sum_\nu p_\nu=1,
\]
并在测量仪末态空间选取任意完备正交基 $\{|\mu\rangle\}$。偏迹的基展开为
\[
\Tr_M X=\sum_\mu\langle\mu|X|\mu\rangle_M.
\]
因此正文\eqnrefs{eq:unitary_dilation_partial_trace_definition} 逐步化为
\begin{align*}
\mathscr E(\rho_S)
&=\sum_\mu
\langle\mu|U(\rho_S\otimes\rho_M)U^\dagger|\mu\rangle\\
&=\sum_{\mu,\nu}p_\nu
\langle\mu|U(\rho_S\otimes|\nu\rangle\langle\nu|)U^\dagger|\mu\rangle.
\end{align*}
系统算符 $\rho_S$ 不作用于测量仪空间，因此
\begin{align*}
&\langle\mu|U(\rho_S\otimes|\nu\rangle\langle\nu|)U^\dagger|\mu\rangle\\
&\qquad=
\langle\mu|U|\nu\rangle\,
\rho_S\,
\langle\nu|U^\dagger|\mu\rangle.
\end{align*}
定义只作用于系统 Hilbert 空间的 Kraus 算符
\[
K_{\mu\nu}
\equiv\sqrt{p_\nu}\,\langle\mu|U|\nu\rangle,
\]
便有
\[
\mathscr E(\rho_S)
=\sum_{\mu\nu}K_{\mu\nu}\rho_SK_{\mu\nu}^\dagger.
\]

完全正性不能只检验系统自身的正性，还要允许系统与任意参考空间 $R$ 纠缠。对任意
$X_{SR}\succeq0$，扩张映射为
\[
(\mathscr E\otimes\mathbb I_R)(X_{SR})
=\sum_{\mu\nu}
(K_{\mu\nu}\otimes I_R)X_{SR}
(K_{\mu\nu}^\dagger\otimes I_R).
\]
每一项都具有 $YXY^\dagger$ 形式，故为正算符；正算符之和仍为正。因此 $\mathscr E$ 对任意辅助系统都保持正性，即完全正。

再验证迹保持。先计算 Kraus 完备关系：
\begin{align*}
\sum_{\mu\nu}K_{\mu\nu}^\dagger K_{\mu\nu}
&=\sum_{\mu\nu}p_\nu
\langle\nu|U^\dagger|\mu\rangle
\langle\mu|U|\nu\rangle\\
&=\sum_\nu p_\nu
\langle\nu|U^\dagger
\left(\sum_\mu|\mu\rangle\langle\mu|\right)
U|\nu\rangle\\
&=\sum_\nu p_\nu\langle\nu|U^\dagger U|\nu\rangle\\
&=\sum_\nu p_\nu \mathbb I_S\\
&=\mathbb I_S.
\end{align*}
于是
\begin{align*}
\Tr\mathscr E(\rho_S)
&=\sum_{\mu\nu}\Tr(K_{\mu\nu}\rho_SK_{\mu\nu}^\dagger)\\
&=\Tr\!\left[\left(\sum_{\mu\nu}K_{\mu\nu}^\dagger K_{\mu\nu}\right)\rho_S\right]\\
&=\Tr\rho_S.
\end{align*}
所以该映射同时完全正且保迹，得到正文\eqnrefs{eq:general_unitary_dilation_kraus_operator}。
\end{derivation}

\begin{derivation}[从\eqnrefs{eq:general_unitary_dilation_kraus_operator}到\eqnrefs{eq:general_indirect_measurement_instrument}]
联合幺正结束以后，系统与测量仪的联合态为
\[
\rho'_{SM}=U(\rho_S\otimes\rho_M)U^\dagger.
\]
测量仪结果 $x$ 由效应 $\Pi_x^{(M)}\succeq0$ 描述，且
\[
\sum_x\Pi_x^{(M)}=\mathbb I_M.
\]
取正平方根
\[
M_x=[\Pi_x^{(M)}]^{1/2},
\qquad
M_x^\dagger M_x=\Pi_x^{(M)}.
\]
在测量仪上实施该结果对应的 L\"uders 操作后，联合未归一化分支为
\[
\widetilde\rho'_{SM|x}
=(\mathbb I_S\otimes M_x)\rho'_{SM}
(\mathbb I_S\otimes M_x^\dagger).
\]
只保留系统意味着对 $M$ 取偏迹，因此定义
\[
\mathcal I_x(\rho_S)=\Tr_M\widetilde\rho'_{SM|x}.
\]
把 $\rho'_{SM}$ 代回：
\begin{align*}
\mathcal I_x(\rho_S)
&=\Tr_M\!\left[
(\mathbb I_S\otimes M_x)
U(\rho_S\otimes\rho_M)U^\dagger
(\mathbb I_S\otimes M_x^\dagger)
\right]\\
&=\Tr_M\!\left[
(\mathbb I_S\otimes [\Pi_x^{(M)}]^{1/2})
U(\rho_S\otimes\rho_M)U^\dagger
(\mathbb I_S\otimes [\Pi_x^{(M)}]^{1/2})
\right],
\end{align*}
得到正文\eqnrefs{eq:general_indirect_measurement_instrument}。

该分支的迹就是结果概率。利用偏迹与全迹的关系以及循环性，
\begin{align*}
p(x)
&=\Tr_S\mathcal I_x(\rho_S)\\
&=\Tr_{SM}\!\left[
(\mathbb I_S\otimes M_x^\dagger M_x)
\rho'_{SM}
\right]\\
&=\Tr_{SM}\!\left[
(\mathbb I_S\otimes\Pi_x^{(M)})
U(\rho_S\otimes\rho_M)U^\dagger
\right].
\end{align*}
对所有结果求和时，
\begin{align*}
\sum_x\mathcal I_x(\rho_S)
&=\Tr_M\!\left[
\sum_x(\mathbb I_S\otimes M_x)
\rho'_{SM}
(\mathbb I_S\otimes M_x^\dagger)
\right],
\end{align*}
其迹满足
\begin{align*}
\Tr_S\sum_x\mathcal I_x(\rho_S)
&=\Tr_{SM}\!\left[
(\mathbb I_S\otimes\sum_x\Pi_x^{(M)})\rho'_{SM}
\right]\\
&=\Tr_{SM}\rho'_{SM}=1.
\end{align*}
因此每个 $\mathcal I_x$ 是迹不增加的完全正分支，而 $\sum_x\mathcal I_x$ 保迹，正是量子仪器的归一化结构。
\end{derivation}
\subsection{前向模型与信息界}

测量算符给出了从态到概率的规则。把相互作用和传播所得的态代入，就能得到概率对实验参数的依赖。设 $\boldsymbol\theta$ 是这些待估参数组成的向量，并用 $\mathscr C_{\rm int}$、$\mathscr C_{\rm prop}$、$\mathscr C_{\rm env}$ 分别表示相互作用、传播及环境阶段的完全正保迹映射。在这三个阶段可以按顺序分开描述的条件下，探测器前的态为
\begin{equation}
\boxed{
\rho_{\rm det}(\boldsymbol\theta)
=(\mathscr C_{\rm env}\circ\mathscr C_{\rm prop}\circ\mathscr C_{\rm int})(\boldsymbol\theta)[\rho_0].}
\label{eq:detector_pre_state_channel_composition}
\end{equation}
若理想量子结果 $x$ 的 POVM 效应算符是 $\Pi_x$，其后还有纯经典响应 $R(y|x)$，则完整实验似然函数是
\begin{equation}
\boxed{
p_{\rm meas}(y|\boldsymbol\theta)
=\sum_xR(y|x)
\operatorname{Tr}\!\left[
\Pi_x(\boldsymbol\theta_{\rm det})
(\mathscr C_{\rm env}\circ\mathscr C_{\rm prop}\circ\mathscr C_{\rm int})(\boldsymbol\theta)[\rho_0]
\right].}
\label{eq:general_experimental_forward_model}
\end{equation}
连续结果时把求和替换成对应测度积分。这个分层保证“相互作用改变状态”“未观测环境导致量子信道”“探测器把量子结果映射成数据”三件事不会被混在同一个唯象展宽中。

若目标是重建有限维子空间中的未知态，选择 Hilbert--Schmidt 正交的厄米基 $\{B_a\}$。未知态和每个设置下的 Born 概率先写成
\begin{equation}
\boxed{
\rho=\sum_a r_aB_a,
\qquad
p(y|s)=\operatorname{Tr}(E_{y|s}\rho).}
\label{eq:finite_dimensional_state_operator_expansion}
\end{equation}
其中全部已知控制都吸收到有效效应算符 $E_{y|s}$。代入基展开后，Born 规则成为线性算符框架
\begin{equation}
\boxed{
\mathbf p=\mathsf A\mathbf r,
\qquad
A_{(y,s),a}=\operatorname{Tr}(E_{y|s}B_a).}
\label{eq:measurement_operator_frame}
\end{equation}
在迹已固定为一的无迹厄米子空间上，信息完备性等价于 $\mathsf A$ 没有非零零方向：
\[
\ker\mathsf A\cap\mathsf{Herm}_0(\mathcal H_S)=\{0\}.
\]
这只是唯一性条件；实际稳定性由框架算符 $\mathsf F=\mathsf A^{\mathsf T}\mathsf W\mathsf A$ 的小奇异值控制。正则化可以抑制噪声，但不能创造测量从未编码的信息。

若目标只是估计参数 $\theta$，则由真实前向模型产生的概率 $p_x(\theta)$ 定义经典 Fisher 信息
\begin{equation}
\boxed{
F_C(\theta)=\sum_x\frac{[\partial_\theta p_x(\theta)]^2}{p_x(\theta)}.}
\label{eq:classical_fisher_information_definition}
\end{equation}
概率对参数的相对敏感度
\begin{equation*}
s_\theta(x)\equiv\partial_\theta\ln p_x(\theta)
\end{equation*}
称为该结果的\emph{得分函数}（score）。它的均值为零，而其平方的概率平均正是 $F_C=\sum_xp_xs_\theta^2$；因此 Fisher 信息衡量小幅改变参数时，整套观测概率分布平均能够改变多快。对多参数情形，相应得分向量为 $\mathbf s(x)=\nabla_{\boldsymbol\theta}\ln p_x$。

$N$ 次独立重复的局域无偏估计满足
\begin{equation}
\boxed{
\operatorname{Var}(\hat\theta)\ge\frac1{N F_C(\theta)}.}
\label{eq:classical_cramer_rao_general}
\end{equation}
探测器损失与模糊必须先进入 $p_x$ 再计算 Fisher 信息，因此它们不能在信息分析以后通过一个事后“校正因子”恢复。

真实反演通常同时包含多个未知参数。对参数向量 $\boldsymbol\theta=(\theta_1,\ldots,\theta_d)$，经典 Fisher 信息推广为矩阵
\begin{equation}
 \boxed{
 [\mathsf I_C(\boldsymbol\theta)]_{ab}
 =\sum_x\frac{\partial_a p_x\,\partial_b p_x}{p_x}.}
 \label{eq:multiparameter-fisher-matrix-dp31}
\end{equation}
若参数分成目标参数 $\psi$ 与干扰参数向量 $\boldsymbol\theta_{\rm n}$，并且干扰参数块 $\mathsf I_{{\rm n}{\rm n}}$ 可逆，则在同时拟合 $\boldsymbol\theta_{\rm n}$ 时，$\psi$ 的局部有效信息是 Schur 补
\begin{equation}
 \boxed{
 I_{\rm eff}(\psi)
 =I_{\psi\psi}
 -\mathsf I_{\psi{\rm n}}
 \mathsf I_{{\rm n}{\rm n}}^{-1}
 \mathsf I_{{\rm n}\psi}.}
 \label{eq:fisher-schur-nuisance-dp31}
\end{equation}
第二项非负，因此不知道干扰参数只会降低目标参数的可辨识性。若该 Schur 补为零，即使 $I_{\psi\psi}$ 本身很大，目标参数也会因为与干扰方向局部退化而不可辨识。

纯经典、且与待估参数无关的探测器响应也满足信息数据处理不等式。若理想结果 $x$ 经过随机核 $R(y|x)$ 变成记录 $y$，粗粒化概率满足
\begin{equation}
\boxed{
\bar p_y(\boldsymbol\theta)=\sum_xR(y|x)p_x(\boldsymbol\theta),
\qquad
\sum_yR(y|x)=1.}
\label{eq:classical_detector_coarse_graining}
\end{equation}
对这一概率重新计算 Fisher 矩阵，得到
\begin{equation}
 \boxed{
 \mathsf I_C^{(y)}\preceq\mathsf I_C^{(x)}.}
 \label{eq:fisher-data-processing-dp31}
\end{equation}
矩阵差为半正定。能量展宽、像素合并、漏计数后的条件重标记等只要可表示成参数无关的经典粗粒化，就不能创造新的 Fisher 信息；若探测器本身含未知参数，则这些参数必须并入 $\boldsymbol\theta$，再用上式的多参数结构共同估计。

量子 Fisher 信息则是给定态族 $\rho_\theta$ 时对所有 POVM 的共同上界。对称对数导数 $L_\theta$ 与量子 Fisher 信息定义为
\begin{equation}
\boxed{
\partial_\theta\rho_\theta
=\frac12(L_\theta\rho_\theta+\rho_\theta L_\theta),
\qquad
F_Q=\operatorname{Tr}(\rho_\theta L_\theta^2).}
\label{eq:qfi-sld-definition-dp68}
\end{equation}
对任意 POVM，经典 Fisher 信息满足
\begin{equation}
\boxed{F_C\le F_Q.}
\label{eq:qfi-measurement-upper-bound-dp68}
\end{equation}
这个不等式把“态本身携带的局域参数信息”与“某个具体测量实际读出的信息”区分开来。对纯态，
\begin{equation}
\boxed{
F_Q
=4\left[
\langle\partial_\theta\Psi|\partial_\theta\Psi\rangle
-|\langle\Psi|\partial_\theta\Psi\rangle|^2
\right].}
\label{eq:pure_state_quantum_fisher_general}
\end{equation}
\eqnrefs{eq:pure_state_quantum_fisher_general} 中第二项恰好扣除了 $|\Psi_\theta\rangle$ 方向上的纯相位变化，因此纯态 QFI 只保留真正改变量子态射线的正交分量。若参数由无量纲厄米生成元 $G$ 编码，$|\Psi_\theta\rangle=\ee^{-\ii\theta G}|\Psi_0\rangle$，则
\begin{equation}
\boxed{F_Q=4(\Delta G)^2.}
\label{eq:unitary_encoding_qfi_variance}
\end{equation}
若实际参数为带量纲的 $\lambda$ 且编码算符为 $\exp[-\ii\lambda K/\hbar]$，则 $G=K/\hbar$，从而 $F_Q^{(\lambda)}=4(\Delta K)^2/\hbar^2$；若生成元本身依赖参数，则必须对指数算符使用有序积分形式的导数，不能直接套用方差公式。因此增强耦合并不自动等价于更高可辨识性；真正的问题是态族怎样随参数改变、实际 POVM 读出了其中多少信息，以及噪声信道是否把这些变化投影到了不可观测方向。



如果实验目标不是重建单一时刻的电子态，而是识别多电子连续碰撞中的环境记忆，那么未知对象应从密度矩阵升级为前文定义的过程张量 $\Upsilon_{N:0}$。这时“数据足够多”必须有一个精确的线性代数含义。选取满足因果约束的过程算符基 $\{F_\alpha\}_{\alpha=1}^{d_{\rm proc}}$，写
\begin{equation*}
\Upsilon_{N:0}=\sum_{\alpha=1}^{d_{\rm proc}}\theta_\alpha F_\alpha.
\end{equation*}
一次实验设置 $r$ 由各时刻插入的电子制备/控制 Choi 算符
$A_{r,0},\ldots,A_{r,N-1}$ 和末端 POVM 元 $M_{r,y}$ 决定。link product 对每个槽位都是线性的，因此结果概率必可写成
\begin{equation}
\boxed{
p_{r,y}
=\sum_{\alpha=1}^{d_{\rm proc}}
X_{(r,y),\alpha}\theta_\alpha,
\qquad
X_{(r,y),\alpha}
\equiv
\Tr\!\left[
M_{r,y}
\left(F_\alpha\star A_{r,N-1}\star\cdots\star A_{r,0}\right)
\right].}
\label{eq:process_tensor_experimental_design_matrix}
\end{equation}
所以过程层析能否唯一完成，不取决于重复次数本身，而取决于实验设计矩阵 $X$ 是否覆盖所有允许的过程方向：
\begin{equation}
\boxed{
\operatorname{rank}X=d_{\rm proc}.}
\label{eq:process_tensor_identifiability_rank}
\end{equation}
若秩亏，增加同一设置的重复计数只能降低已可见方向的统计噪声，不能恢复落在 $\ker X$ 中的记忆参数。

对自由电子实验，这个结论具有直接操作意义。若每个电子都以完全相同的波包入射、使用同一相互作用强度、相同电子间隔，并且最终只测一个固定能谱，那么实验只探测过程张量的一个很低维切片。要区分“电子自身相干”“环境跨碰撞记忆”“材料慢漂移”和“探测器卷积”，通常必须改变至少一类控制：入射电子相位/能量基、电子间延迟、耦合强度、后选择基或末端探测基。过程张量形式把这些变化视为不同的槽位操作，而不是不同的理论模型。

在大样本局部估计中，若计数服从独立多项式/Poisson 统计，过程参数的经典 Fisher 信息仍由概率导数组成：
\begin{equation}
\boxed{
[\mathcal F_{\rm proc}]_{\alpha\beta}
=\sum_{r,y}
\frac{N_r}{p_{r,y}}
X_{(r,y),\alpha}X_{(r,y),\beta},}
\label{eq:process_tensor_fisher_information}
\end{equation}
其中 $N_r$ 是设置 $r$ 的重复次数。若所有被使用的 $p_{r,y}>0$，则
$\mathcal F_{\rm proc}$ 在允许参数空间上正定与 $X$ 满列秩等价。因此“过程可辨识性”和“有限方差估计”由同一个设计矩阵控制；实验优化可以转化为改善 $X$ 的奇异值谱，而不是只追求更多电子计数。


\begin{derivation}[从\eqnrefs{eq:process_tensor_choi_link}到\eqnrefs{eq:process_tensor_fisher_information}]
把因果约束已经吸收到基 $F_\alpha$ 中，过程 Choi 算符展开为
\[
\Upsilon=\sum_\alpha\theta_\alpha F_\alpha.
\]
固定一个实验设置 $r$。多时刻操作的 link product 定义一个从过程算符到末态的线性映射
\[
\mathcal L_r(F)
\equiv
F\star A_{r,N-1}\star\cdots\star A_{r,0}.
\]
末端 Born 规则给
\begin{align*}
p_{r,y}
&=\Tr[M_{r,y}\mathcal L_r(\Upsilon)]\\
&=\sum_\alpha\theta_\alpha
\Tr[M_{r,y}\mathcal L_r(F_\alpha)]\\
&\equiv\sum_\alpha X_{(r,y),\alpha}\theta_\alpha,
\end{align*}
即正文\eqnrefs{eq:process_tensor_experimental_design_matrix}。

若存在非零向量 $v\in\ker X$，则对任意足够小的实数 $\epsilon$，参数
$\theta$ 与 $\theta+\epsilon v$ 产生完全相同的全部一阶概率变化；若二者都保持在物理过程集合内，数据无法区分它们。因此 $\ker X\neq\{0\}$ 意味着局部不可辨识。反之，若 $X$ 满列秩，则线性方程组在无噪声极限中至多有一个参数解，所以局部线性层析唯一；得到正文\eqnrefs{eq:process_tensor_identifiability_rank}。

对独立计数，设置 $r$ 重复 $N_r$ 次。参数导数为
\[
\partial_{\theta_\alpha}p_{r,y}=X_{(r,y),\alpha}.
\]
经典 Fisher 信息定义
\[
[\mathcal F]_{\alpha\beta}
=\sum_{r,y}N_r p_{r,y}
(\partial_\alpha\ln p_{r,y})
(\partial_\beta\ln p_{r,y}),
\]
于是
\[
[\mathcal F]_{\alpha\beta}
=\sum_{r,y}
\frac{N_r}{p_{r,y}}
X_{(r,y),\alpha}X_{(r,y),\beta},
\]
得到正文\eqnrefs{eq:process_tensor_fisher_information}。对任意参数方向 $v$，
\[
v^{\mathsf T}\mathcal Fv
=\sum_{r,y}\frac{N_r}{p_{r,y}}
|(Xv)_{r,y}|^2.
\]
在 $p_{r,y}>0$ 的采样支撑上，这个量为零当且仅当 $Xv=0$。所以 Fisher 信息的零模正是设计矩阵不可见的过程方向。增大 $N_r$ 只放大已有非零奇异方向的权重，不能改变 $\ker X$；这严格证明了为什么“无限重复同一种电子实验”仍不能完成一般非 Markov 过程层析。
\end{derivation}
\begin{derivation}[从\eqnrefs{eq:finite_dimensional_state_operator_expansion}到\eqnrefs{eq:measurement_operator_frame}]
设系统 Hilbert 空间维数为 $d$，选取 Hilbert--Schmidt 正交的厄米算符基
\[
\Tr(B_aB_b)=\delta_{ab},
\qquad a,b=0,\ldots,d^2-1.
\]
取
\[
B_0=\frac{\mathbb I}{\sqrt d},
\]
其余 $B_a$ 为无迹厄米算符。正文的状态展开
\[
\rho=\sum_a r_aB_a
\]
两边左乘 $B_b$ 并取迹：
\begin{align*}
\Tr(B_b\rho)
&=\sum_a r_a\Tr(B_bB_a)\\
&=\sum_a r_a\delta_{ba}\\
&=r_b.
\end{align*}
因此 $r_a=\Tr(B_a\rho)$，且迹条件固定
\[
r_0=\Tr(B_0\rho)=\frac1{\sqrt d}.
\]

对实验设置 $s$ 和结果 $y$，全部已知控制、传播与探测已经包含在系统侧效应 $E_{y|s}$ 中。Born 规则逐步给出
\begin{align*}
p(y|s)
&=\Tr(E_{y|s}\rho)\\
&=\Tr\!\left[E_{y|s}\sum_a r_aB_a\right]\\
&=\sum_a r_a\Tr(E_{y|s}B_a).
\end{align*}
把复合数据指标记为 $\mu=(y,s)$，定义矩阵元
\[
A_{\mu a}\equiv\Tr(E_\mu B_a),
\qquad
p_\mu\equiv p(y|s).
\]
逐个结果的标量等式
\[
p_\mu=\sum_aA_{\mu a}r_a
\]
合并成矩阵形式
\[
\mathbf p=\mathsf A\mathbf r,
\qquad
A_{(y,s),a}=\Tr(E_{y|s}B_a),
\]
即正文\eqnrefs{eq:measurement_operator_frame}。
\end{derivation}
\begin{derivation}[从\eqnrefs{eq:classical_fisher_information_definition}到\eqnrefs{eq:classical_cramer_rao_general}]
设单次观测 $x$ 的概率为 $p_x(\theta)$，并定义得分函数
\[
s_\theta(x)=\partial_\theta\ln p_x(\theta)
=\frac{\partial_\theta p_x}{p_x}.
\]
先验证得分的均值为零。由概率归一化
\[
\sum_xp_x(\theta)=1
\]
对 $\theta$ 求导：
\begin{align*}
0
&=\partial_\theta\sum_xp_x(\theta)\\
&=\sum_x\partial_\theta p_x(\theta)\\
&=\sum_xp_x(\theta)s_\theta(x)\\
&=\mathbb E_\theta[s_\theta].
\end{align*}
于是 Fisher 信息就是得分方差：
\[
F_C(\theta)
=\mathbb E_\theta[s_\theta^2]
=\sum_x\frac{[\partial_\theta p_x]^2}{p_x}.
\]

令 $\hat\theta(x)$ 是局域无偏估计量，满足
\[
\mathbb E_\theta[\hat\theta]=\theta.
\]
对该等式求导：
\begin{align*}
1
&=\partial_\theta\sum_xp_x(\theta)\hat\theta(x)\\
&=\sum_x\hat\theta(x)\partial_\theta p_x(\theta)\\
&=\sum_xp_x(\theta)\hat\theta(x)s_\theta(x).
\end{align*}
因为 $\mathbb E[s_\theta]=0$，可减去任意常数，特别取真参数 $\theta$：
\begin{align*}
1
&=\sum_xp_x(\theta)[\hat\theta(x)-\theta]s_\theta(x)\\
&=\mathbb E_\theta[(\hat\theta-\theta)s_\theta].
\end{align*}
对随机变量 $\hat\theta-\theta$ 与 $s_\theta$ 使用 Cauchy--Schwarz 不等式：
\begin{align*}
1
&\le\sqrt{\mathbb E[(\hat\theta-\theta)^2]}
\sqrt{\mathbb E[s_\theta^2]}\\
&=\sqrt{\operatorname{Var}(\hat\theta)F_C(\theta)}.
\end{align*}
平方并整理得到单次观测的界
\[
\operatorname{Var}(\hat\theta)\ge\frac1{F_C(\theta)}.
\]

对 $N$ 次相互独立、同分布观测
$\mathbf x=(x_1,\ldots,x_N)$，联合概率为
\[
p(\mathbf x|\theta)=\prod_{j=1}^Np_{x_j}(\theta),
\]
所以联合得分是各次得分之和：
\[
s_\theta^{(N)}(\mathbf x)
=\partial_\theta\ln p(\mathbf x|\theta)
=\sum_{j=1}^Ns_\theta(x_j).
\]
独立性和零均值使交叉项消失：
\begin{align*}
F_C^{(N)}
&=\mathbb E[(s_\theta^{(N)})^2]\\
&=\sum_{j=1}^N\mathbb E[s_\theta^2(x_j)]
+\sum_{j\ne k}\mathbb E[s_\theta(x_j)]\mathbb E[s_\theta(x_k)]\\
&=NF_C(\theta).
\end{align*}
代回单次形式即得
\[
\operatorname{Var}(\hat\theta)
\ge\frac1{NF_C(\theta)},
\]
即正文\eqnrefs{eq:classical_cramer_rao_general}。
\end{derivation}

\begin{derivation}[从\eqnrefs{eq:qfi-sld-definition-dp68}到\eqnrefs{eq:qfi-measurement-upper-bound-dp68}]
设 POVM 元为 $E_x\ge0$、$\sum_xE_x=I$，测量概率为
$p_x=\operatorname{Tr}(E_x\rho_\theta)$。SLD 定义给
\begin{align*}
 \partial_\theta p_x
 &=\operatorname{Tr}(E_x\partial_\theta\rho_\theta)\\
 &=\frac12\operatorname{Tr}[E_x(L\rho+\rho L)]\\
 &=\operatorname{Re}\operatorname{Tr}(E_xL\rho).
\end{align*}
最后一步使用 $E_x,L,\rho$ 的厄米性：两项互为复共轭。

为了把每个概率导数与 QFI 中的 $L^2$ 联系起来，在 Hilbert--Schmidt 内积
$\langle A,B\rangle_{\rm HS}=\operatorname{Tr}(A^\dagger B)$ 下定义
\begin{equation*}
 A_x=E_x^{1/2}\rho^{1/2},
 \qquad
 B_x=E_x^{1/2}L\rho^{1/2}.
\end{equation*}
于是
\begin{align*}
 \operatorname{Tr}(A_x^\dagger B_x)
 &=\operatorname{Tr}(\rho^{1/2}E_xL\rho^{1/2})
 =\operatorname{Tr}(E_xL\rho),\\
 \operatorname{Tr}(A_x^\dagger A_x)
 &=\operatorname{Tr}(E_x\rho)=p_x,\\
 \operatorname{Tr}(B_x^\dagger B_x)
 &=\operatorname{Tr}(\rho L E_xL).
\end{align*}
因此先用 $|\operatorname{Re}z|^2\le|z|^2$，再用 Cauchy--Schwarz，得到
\begin{equation*}
 |\partial_\theta p_x|^2
 \le |\operatorname{Tr}(A_x^\dagger B_x)|^2
 \le p_x\operatorname{Tr}(\rho L E_xL).
\end{equation*}
先处理 $p_x=0$ 的边界情形。因为 $E_x\ge0$ 与 $\rho\ge0$，
\begin{equation*}
p_x=\operatorname{Tr}(E_x\rho)
=\operatorname{Tr}\!\left[
(E_x^{1/2}\rho^{1/2})^\dagger
(E_x^{1/2}\rho^{1/2})
\right]=0
\end{equation*}
蕴含
\begin{equation*}
E_x^{1/2}\rho^{1/2}=A_x=0.
\end{equation*}
于是
$\partial_\theta p_x=\operatorname{Re}\operatorname{Tr}(A_x^\dagger B_x)=0$；在参数点附近可微的正则模型中，这个结果对 Fisher 和式不产生奇异贡献。因而只需对 $p_x>0$ 的支撑除以 $p_x$：
\begin{equation*}
\frac{(\partial_\theta p_x)^2}{p_x}
\le\operatorname{Tr}(\rho L E_xL).
\end{equation*}
随后对所有测量结果求和，并使用 POVM 完备性：
\begin{align*}
 F_C
 &=\sum_{x:p_x>0}\frac{(\partial_\theta p_x)^2}{p_x}\\
 &\le\sum_x\operatorname{Tr}(\rho L E_xL)\\
 &=\operatorname{Tr}\!\left[
\rho L\left(\sum_xE_x\right)L
\right]\\
 &=\operatorname{Tr}(\rho L^2)\\
 &=F_Q.
\end{align*}
从概率导数到最后一行只用了 SLD 定义、Hilbert--Schmidt Cauchy--Schwarz 不等式和 $\sum_xE_x=I$，因此该界与具体探测器实现无关。这就是正文\eqnrefs{eq:qfi-measurement-upper-bound-dp68}。
\end{derivation}
等号要求 Cauchy--Schwarz 与取实部两步同时饱和；局域测量 SLD 本征投影可满足这一结构。因此 QFI 是对允许 POVM 优化后的局域上界，而不是固定探测器自动达到的灵敏度。
% DeepPass22: detailed pure-state QFI and SLD algebra.

\begin{derivation}[从\eqnrefs{eq:qfi-sld-definition-dp68}到\eqnrefs{eq:unitary_encoding_qfi_variance}]
令
\begin{equation*}
\rho=|\Psi\rangle\langle\Psi|,
\qquad
|\dot\Psi\rangle\equiv\partial_\theta|\Psi_\theta\rangle.
\end{equation*}
归一化条件给
\begin{equation*}
0=\partial_\theta\langle\Psi|\Psi\rangle
=\langle\dot\Psi|\Psi\rangle+\langle\Psi|\dot\Psi\rangle.
\end{equation*}
记 $a=\langle\Psi|\dot\Psi\rangle$，并分解
\begin{equation*}
|\dot\Psi\rangle=a|\Psi\rangle+|\dot\Psi_\perp\rangle,
\qquad
\langle\Psi|\dot\Psi_\perp\rangle=0.
\end{equation*}
由于 $a+a^*=0$，
\begin{equation*}
\dot\rho
=|\dot\Psi_\perp\rangle\langle\Psi|
+|\Psi\rangle\langle\dot\Psi_\perp|.
\end{equation*}
取
\begin{equation*}
L=2\left(
|\dot\Psi_\perp\rangle\langle\Psi|
+|\Psi\rangle\langle\dot\Psi_\perp|
\right),
\end{equation*}
则利用 $\rho|\dot\Psi_\perp\rangle=0$ 可验证
\begin{equation*}
\frac12(L\rho+\rho L)=\dot\rho.
\end{equation*}
再令 $|v\rangle=|\dot\Psi_\perp\rangle$，有
\begin{equation*}
\frac{L^2}{4}
=|v\rangle\langle v|+\langle v|v\rangle|\Psi\rangle\langle\Psi|,
\end{equation*}
因此
\begin{align*}
F_Q
&=\operatorname{Tr}(\rho L^2)
=4\langle v|v\rangle\\
&=4\left[
\langle\dot\Psi|\dot\Psi\rangle
-|\langle\Psi|\dot\Psi\rangle|^2
\right],
\end{align*}
即正文\eqnrefs{eq:pure_state_quantum_fisher_general}。

若 $|\Psi_\theta\rangle=\ee^{-\ii\theta G}|\Psi_0\rangle$ 且 $G=G^\dagger$ 与 $\theta$ 无关，则
\begin{equation*}
|\dot\Psi_\theta\rangle=-\ii G|\Psi_\theta\rangle,
\end{equation*}
从而
\begin{equation*}
\langle\dot\Psi|\dot\Psi\rangle=\langle G^2\rangle,
\qquad
|\langle\Psi|\dot\Psi\rangle|^2=\langle G\rangle^2.
\end{equation*}
代入上一式得到
\begin{equation*}
F_Q=4\left(\langle G^2\rangle-\langle G\rangle^2\right)
=4(\Delta G)^2,
\end{equation*}
即正文\eqnrefs{eq:unitary_encoding_qfi_variance}。
\end{derivation}


\begin{derivation}[从\eqnrefs{eq:multiparameter-fisher-matrix-dp31}到\eqnrefs{eq:fisher-schur-nuisance-dp31}]
把待估参数拆成目标标量 $\psi$ 与干扰参数向量 $\boldsymbol\theta_{\rm n}$。单次实验的 Fisher 矩阵写成
\[
\mathsf I_C=
\begin{pmatrix}
 a & \mathbf b^{\mathsf T}\\
 \mathbf b & \mathsf C
\end{pmatrix},
\qquad
 a\equiv I_{\psi\psi},\quad
 \mathbf b\equiv\mathsf I_{{\rm n}\psi},\quad
 \mathsf C\equiv\mathsf I_{{\rm n}{\rm n}}.
\]
先假设 $\mathsf C$ 在局部可辨识的干扰子空间上可逆。与其直接引用分块逆公式，不妨解线性方程
\[
\mathsf I_C
\begin{pmatrix}x\\ \mathbf y\end{pmatrix}
=
\begin{pmatrix}1\\ \mathbf 0\end{pmatrix}.
\]
第二行给
\[
\mathbf b\,x+\mathsf C\mathbf y=0
\quad\Longrightarrow\quad
\mathbf y=-\mathsf C^{-1}\mathbf b\,x.
\]
代回第一行：
\begin{align*}
ax+\mathbf b^{\mathsf T}\mathbf y&=1,\\
\left(a-\mathbf b^{\mathsf T}\mathsf C^{-1}\mathbf b\right)x&=1.
\end{align*}
而 $x$ 正是 $\mathsf I_C^{-1}$ 的左上角元素，因此
\[
[\mathsf I_C^{-1}]_{\psi\psi}
=\frac{1}{a-\mathbf b^{\mathsf T}\mathsf C^{-1}\mathbf b}.
\]
定义
\[
\boxed{
I_{\rm eff}(\psi)
\equiv
I_{\psi\psi}
-\mathsf I_{\psi{\rm n}}
\mathsf I_{{\rm n}{\rm n}}^{-1}
\mathsf I_{{\rm n}\psi}},
\]
多参数 Cram\'er--Rao 界
\[
\operatorname{Cov}(\hat{\boldsymbol\theta})
\succeq\frac1N\mathsf I_C^{-1}
\]
就在目标参数方向给出
\[
\boxed{
\operatorname{Var}(\hat\psi)
\ge\frac1{N I_{\rm eff}(\psi)}}.
\]

扣除项的来源也可从 score 几何直接看出。记单次数据的目标 score 为 $s_\psi$，干扰 score 向量为 $\mathbf s_{\rm n}$，则
\[
a=\mathbb E[s_\psi^2],\qquad
\mathbf b=\mathbb E[\mathbf s_{\rm n}s_\psi],\qquad
\mathsf C=\mathbb E[\mathbf s_{\rm n}\mathbf s_{\rm n}^{\mathsf T}].
\]
在线性组合 $s_\psi-\mathbf c^{\mathsf T}\mathbf s_{\rm n}$ 中选择 $\mathbf c$ 使均方最小：
\[
\frac{\partial}{\partial\mathbf c}
\mathbb E[(s_\psi-\mathbf c^{\mathsf T}\mathbf s_{\rm n})^2]
=-2\mathbf b+2\mathsf C\mathbf c=0,
\]
故 $\mathbf c_*=\mathsf C^{-1}\mathbf b$，最小剩余范数恰为
\[
\mathbb E[(s_\psi-\mathbf c_*^{\mathsf T}\mathbf s_{\rm n})^2]
=a-\mathbf b^{\mathsf T}\mathsf C^{-1}\mathbf b
=I_{\rm eff}.
\]
所以 Schur 补不是抽象的矩阵技巧：它把目标 score 中能够由干扰参数 score 线性模仿的方向投影掉，只保留实验真正能独立归因于 $\psi$ 的信息。若 $\mathsf C$ 奇异，应限制到其可辨识支撑空间或使用 Moore--Penrose 逆；零模对应数据无法区分的真实参数退化。
\end{derivation}

\begin{derivation}[从\eqnrefs{eq:classical_detector_coarse_graining}到\eqnrefs{eq:fisher-data-processing-dp31}]
设理想量子结果为 $x$，实际记录为 $y$。纯经典探测器后处理由与待估参数 $\boldsymbol\theta$ 无关的 Markov 核 $R(y|x)$ 描述：
\[
R(y|x)\ge0,
\qquad
\sum_yR(y|x)=1,
\qquad
\bar p_y=\sum_xR(y|x)p_x.
\]
对参数求导时，$R$ 本身没有导数，因此
\begin{align*}
\nabla_{\boldsymbol\theta}\bar p_y
&=\sum_xR(y|x)\nabla_{\boldsymbol\theta}p_x\\
&=\sum_xR(y|x)p_x\,\mathbf s(x),
\end{align*}
其中理想数据的 score 向量定义为
\[
\mathbf s(x)=\nabla_{\boldsymbol\theta}\ln p_x.
\]
实际记录 $y$ 的 score 为
\begin{align*}
\bar{\mathbf s}(y)
&=\nabla_{\boldsymbol\theta}\ln\bar p_y\\
&=\frac1{\bar p_y}
\sum_xR(y|x)p_x\mathbf s(x).
\end{align*}
定义给定 $y$ 后理想结果 $x$ 的后验权重
\[
w(x|y)=\frac{R(y|x)p_x}{\bar p_y},
\qquad
\sum_xw(x|y)=1.
\]
于是
\[
\bar{\mathbf s}(y)=\sum_xw(x|y)\mathbf s(x)
=\mathbb E[\mathbf s(X)|Y=y].
\]
为了证明 Fisher 矩阵的半正定序关系，取任意实参数方向 $\mathbf v$，并定义标量投影
\[
s_x=\mathbf v^{\mathsf T}\mathbf s(x),
\qquad
t_y=\mathbf v^{\mathsf T}\bar{\mathbf s}(y).
\]
由上式
\[
t_y=\mathbb E[s_X|Y=y].
\]
对凸函数 $f(s)=s^2$ 使用条件 Jensen 不等式：
\[
t_y^2
=\bigl(\mathbb E[s_X|Y=y]\bigr)^2
\le\mathbb E[s_X^2|Y=y].
\]
乘以 $\bar p_y$ 并对 $y$ 求和：
\begin{align*}
\sum_y\bar p_yt_y^2
&\le\sum_y\bar p_y
\sum_xw(x|y)s_x^2\\
&=\sum_{x,y}R(y|x)p_xs_x^2\\
&=\sum_xp_xs_x^2\sum_yR(y|x)\\
&=\sum_xp_xs_x^2.
\end{align*}
两端分别就是实际记录和理想结果的 Fisher 二次型：
\begin{align*}
\mathbf v^{\mathsf T}\mathsf I_C^{(y)}\mathbf v
&=\sum_y\bar p_y
[\mathbf v^{\mathsf T}\bar{\mathbf s}(y)]^2,\\
\mathbf v^{\mathsf T}\mathsf I_C^{(x)}\mathbf v
&=\sum_xp_x
[\mathbf v^{\mathsf T}\mathbf s(x)]^2.
\end{align*}
因此对任意 $\mathbf v$ 都有
\[
\mathbf v^{\mathsf T}
(\mathsf I_C^{(x)}-\mathsf I_C^{(y)})\mathbf v\ge0,
\]
这正是矩阵半正定序
\[
\mathsf I_C^{(y)}\preceq\mathsf I_C^{(x)},
\]
即正文\eqnrefs{eq:fisher-data-processing-dp31}。
\end{derivation}

\subsection{可扫描参考操作与层析可辨识性}

在能量基底中，能谱仪给出的概率只含密度矩阵的对角元。若还要确定非对角相干元，可以在探测前施加一个已知操作，让不同通道先发生干涉，再测各条边带的概率。改变参考操作，便得到关于同一未知密度矩阵的多组线性方程。从这些概率恢复量子态的过程称为\emph{量子态层析}。本节用线性代数判断这些方程是否足以确定所需态分量。

\begin{exbox}[能量格与纵向相空间层析]
电子能谱仪只直接读取电子密度算符在能量基底中的对角元；若未知态含有不同能量格点之间的相干，单次能量测量不足以确定这些非对角元。层析因此需要一个\emph{已知、可重复且可扫描}的参考操作，把不同初始能量路径相干地汇合到同一个末态能量通道。实验中参考操作先改变待测电子态，再测量能量；数学上也可把这一过程等价地写成原始未知态上的旋转测量算符。扫描参考相位后，不同密度矩阵副对角线在相位傅里叶空间中被分离。

在有限能量窗口内，用整数 $\ell,n\in\mathbb Z$ 标记电子能量格基 $|\ell\rangle$。这些基矢无量纲且正交，未知电子态写成
\begin{equation*}
\hat\rho_e
=\sum_{\ell,n}\rho_{\ell n}|\ell\rangle\langle n|,
\qquad
\rho_{\ell n}\equiv\langle\ell|\hat\rho_e|n\rangle,
\end{equation*}
其中 $\rho_{\ell n}$ 无量纲，并满足 $\hat\rho_e\succeq0$、$\rho_{n\ell}=\rho_{\ell n}^*$ 与 $\sum_\ell\rho_{\ell\ell}=1$。对角元 $\rho_{\ell\ell}$ 是能量格占据概率，非对角元 $\rho_{\ell n}$ 则给出相隔 $|\ell-n|$ 个能量量子的相干幅度。

取一个已标定的平移协变参考幺正算符 $\hat U_R(\varphi_R)$。扫描参数 $\varphi_R\in[0,2\pi)$ 是无量纲相位；平移协变性意味着矩阵元只依赖末态与初态格点之差，因此定义无量纲参考振幅 $u_r$：
\begin{equation}
\boxed{
\langle m|\hat U_R(\varphi_R)|\ell\rangle
=u_{m-\ell}\,\ee^{\ii(m-\ell)\varphi_R}.}
\label{eq:reference-pinem-matrix}
\end{equation}
这里 $r=m-\ell\in\mathbb Z$ 是参考操作引起的净能量格位移。由于 $\hat U_R$ 幺正，这些系数不能任意选择，而必须满足
\begin{equation*}
\sum_m u_{m-\ell}u_{m-n}^*=\delta_{\ell n}.
\end{equation*}
实验上，$u_r$ 由已知参考场的耦合强度、相互作用长度和电子速度等标定参数决定；层析时它们被视为前向模型中已知的复数系数。

参考操作后理想测量末态能量格点 $m$。对应 POVM 元素为 $|m\rangle\langle m|$，Born 规则给出无量纲概率
\begin{equation}
\boxed{
P_m(\varphi_R)
=\sum_{\ell,n}
 u_{m-\ell}u_{m-n}^*\rho_{\ell n}
\ee^{\ii(n-\ell)\varphi_R}.}
\label{eq:electron_tomography_linear}
\end{equation}
\eqnrefs{eq:electron_tomography_linear} 已经是未知密度矩阵到可测能谱的线性前向模型。相位 $\varphi_R$ 的作用是给不同副对角线乘上不同 Fourier 谐波，因此对 $P_m(\varphi_R)$ 作相位 Fourier 分解可以逐条抽取副对角线。

定义第 $\Delta\ell\in\mathbb Z$ 阶 Fourier 系数
\begin{equation*}
\widetilde P_m^{(\Delta\ell)}
\equiv\frac{1}{2\pi}
\int_0^{2\pi}\dd\varphi_R\,
\ee^{-\ii (\Delta\ell)\varphi_R}P_m(\varphi_R).
\end{equation*}
$\widetilde P_m^{(\Delta\ell)}$ 仍无量纲。相位正交性只保留 $n-\ell=\Delta\ell$，于是
\begin{equation}
\boxed{
\widetilde P_m^{(\Delta\ell)}
=\sum_\ell \mathsf A_{m\ell}^{(\Delta\ell)}\,
\rho_{\ell,\ell+\Delta\ell},
\qquad
\mathsf A_{m\ell}^{(\Delta\ell)}
\equiv u_{m-\ell}u_{m-\ell-\Delta\ell}^*.}
\label{eq:electron_tomography_side_diagonal}
\end{equation}
$\mathsf A^{(\Delta\ell)}$ 是第 $\Delta\ell$ 条副对角线的\emph{层析设计矩阵}，其矩阵元无量纲。它由实验参考操作完全决定，而向量 $\rho_{\ell,\ell+\Delta\ell}$ 才是需要反演的未知量。

要判断某条副对角线是否真的可以由实验数据确定，不能只比较测量方程数目；必须检查 $\mathsf A^{(\Delta\ell)}$ 在目标参数子空间中的奇异值谱。写奇异值分解
\begin{equation}
\boxed{
\mathsf A^{(\Delta\ell)}
=\mathsf U_{\Delta\ell}\,
\operatorname{diag}(s_1^{(\Delta\ell)},s_2^{(\Delta\ell)},\ldots)
\mathsf V_{\Delta\ell}^\dagger,}
\label{eq:electron_tomography_svd}
\end{equation}
其中 $\mathsf U_{\Delta\ell},\mathsf V_{\Delta\ell}$ 为无量纲酉矩阵，$s_a^{(\Delta\ell)}\ge0$ 为无量纲奇异值。若 $s_a^{(\Delta\ell)}=0$，对应的右奇异向量位于测量映射零空间中：沿该密度矩阵方向改变未知态不会改变任何当前设置下的理想数据，因此该方向物理上不可辨识。若 $0<s_a^{(\Delta\ell)}\ll1$，伪逆会把数据噪声沿该方向放大约 $1/s_a^{(\Delta\ell)}$；这时改变参考耦合、增加传播设置或改变读出基可以改变设计矩阵并抬高最小有效奇异值，而单纯增加同一设置的重复计数只能降低统计噪声，不能移除零空间。

离散能量格层析适用于有限 PINEM/QPINEM 边带窗口。若实验通过二次电子光学连续扫描纵向相空间，则同一测量思想可写成 Wigner--Radon 层析。这里必须先把有量纲的纵向位置 $z$ 与动量 $p$ 化为一对正则无量纲变量。任选一个固定长度标尺 $z_0>0$，定义
\begin{equation}
\boxed{
X\equiv\frac{z}{z_0},
\qquad
P\equiv\frac{z_0p}{\hbar},
\qquad
[\hat X,\hat P]=\ii.}
\label{eq:dimensionless_longitudinal_phase_space}
\end{equation}
$z_0$ 的单位为米，用来选择相空间坐标尺度；可测结论对这一任意尺度的选择保持不变。有量纲 Wigner 函数 $W(z,p)$ 由\eqnrefs{eq:electron_wigner} 定义，并满足
\[
[W]=({\rm J\,s})^{-1},
\qquad
\int\dd z\dd p\,W(z,p)=1,
\]
因为 $\dd z\dd p$ 的单位正是作用量。对应的无量纲 Wigner 函数定义为
\begin{equation*}
W_{XP}(X,P)
\equiv\hbar\,W\!\left(z_0X,\frac{\hbar P}{z_0}\right),
\qquad
\int\dd X\dd P\,W_{XP}=1.
\end{equation*}

已标定的二次电子光学对 $(X,P)$ 实现线性辛旋转。用无量纲旋转角 $\theta$ 表示该设置，并定义被读出的旋转象限
\begin{equation*}
X_\theta=X\cos\theta+P\sin\theta.
\end{equation*}
若探测器理想测量 $X_\theta$，其结果记为无量纲变量 $x$，概率密度 $\mathcal P_\theta(x)$ 对 $x$ 归一化。Wigner 边缘性质给出
\begin{equation}
\boxed{
\mathcal P_\theta(x)
=\int\dd X\dd P\,
W_{XP}(X,P)
\delta[x-X\cos\theta-P\sin\theta].}
\label{eq:wigner_radon_transform}
\end{equation}
这就是 $W_{XP}$ 沿方向 $\theta$ 的 Radon 投影。

令 $s_X$ 为与无量纲读出变量 $x$ Fourier 共轭的无量纲 Fourier 变量，定义
\begin{equation*}
\widetilde{\mathcal P}_\theta(s_X)
\equiv\int\dd x\,\ee^{-\ii s_Xx}\mathcal P_\theta(x).
\end{equation*}
二维 Fourier 切片定理给出
\begin{equation}
\boxed{
\widetilde{\mathcal P}_\theta(s_X)
=\widetilde W(s_X\cos\theta,s_X\sin\theta),}
\label{eq:wigner_fourier_slice_theorem}
\end{equation}
其中 $\widetilde W(K_X,K_P)$ 是 $W_{XP}$ 的二维 Fourier 变换。于是扫描 $\theta$ 等价于在二维 Fourier 平面上采集不同径向切片。理想连续数据可由
\begin{equation}
\boxed{
W_{XP}(X,P)
=\frac{1}{(2\pi)^2}
\int_0^\pi\dd\theta
\int_{-\infty}^{\infty}\dd s_X\,|s_X|
\widetilde{\mathcal P}_\theta(s_X)
\ee^{\ii s_X(X\cos\theta+P\sin\theta)}}
\label{eq:wigner_filtered_backprojection}
\end{equation}
重建；其中 $|s_X|$ 来自二维 Fourier 平面改用有符号径向坐标后的雅可比行列式，因此高 $|s_X|$ 分量在反投影中被放大。

有限空间分辨率必须进入前向模型。令 $x_d$ 为探测器记录的无量纲读出，并定义条件线形 $r_X(x_d|x)$，满足
\begin{equation*}
r_X(x_d|x)\ge0,
\qquad
\int\dd x_d\,r_X(x_d|x)=1.
\end{equation*}
若响应平移不变，可写 $r_X(x_d|x)=r_X(x_d-x)$。观测投影为卷积，Fourier 空间中有
\begin{equation}
\boxed{
\widetilde{\mathcal P}^{\rm obs}_\theta(s_X)
=\widetilde r_X(s_X)\,
\widetilde{\mathcal P}_\theta(s_X).}
\label{eq:wigner_tomography_detector_transfer}
\end{equation}
当 $|\widetilde r_X(s_X)|$ 在高频趋近零时，形式反卷积 $1/\widetilde r_X(s_X)$ 会放大计数噪声；这与离散层析中小奇异值导致的噪声放大是同一个线性逆问题结构。


\medskip

自由漂移对纵向 Wigner 函数实施相空间剪切；不同漂移距离的投影等价于不同 Radon 角，从而允许层析。 建立无量纲相空间变量，推出 Radon 投影、Fourier 切片定理和带探测器传递函数的滤波反投影~\eqnrange{eq:dimensionless_longitudinal_phase_space}{eq:wigner_tomography_detector_transfer}；这里采用二次自由传播使映射线性；扫描角度覆盖所需范围；探测器传递函数在可反演频带内非零或用正则化处理。由 $X=z/z_0$、$P=z_0p/\hbar$ 与 $[\hat z,\hat p]=\ii\hbar$，
\[
[\hat X,\hat P]
=\frac1{z_0}\frac{z_0}{\hbar}[\hat z,\hat p]
=\ii.
\]
坐标变换的雅可比行列式为
\[
\left|\frac{\partial(z,p)}{\partial(X,P)}\right|
=z_0\frac{\hbar}{z_0}=\hbar,
\]
因此
\begin{align*}
\int\dd X\dd P\,W_{XP}(X,P)
&=\int\dd X\dd P\,
\hbar W\!\left(z_0X,\frac{\hbar P}{z_0}\right)\\
&=\int\dd z\dd p\,W(z,p)=1.
\end{align*}

辛旋转写成
\[
\begin{pmatrix}X_\theta\\P_\theta\end{pmatrix}
=\begin{pmatrix}
\cos\theta&\sin\theta\\
-\sin\theta&\cos\theta
\end{pmatrix}
\begin{pmatrix}X\\P\end{pmatrix}.
\]
旋转矩阵行列式为一，因此 $\dd X_\theta\dd P_\theta=\dd X\dd P$。测量 $X_\theta=x$ 而不读取 $P_\theta$ 时，
\begin{align*}
\mathcal P_\theta(x)
&=\int\dd P_\theta\,W_\theta(x,P_\theta)\\
&=\int\dd X_\theta\dd P_\theta\,
W_\theta(X_\theta,P_\theta)\delta(x-X_\theta)\\
&=\int\dd X\dd P\,W_{XP}(X,P)
\delta[x-X\cos\theta-P\sin\theta],
\end{align*}
得到\eqnrefs{eq:wigner_radon_transform}。对该式作一维 Fourier 变换：
\begin{align*}
\widetilde{\mathcal P}_\theta(s_X)
&=\int\dd x\,\ee^{-\ii s_Xx}\mathcal P_\theta(x)\\
&=\int\dd X\dd P\,W_{XP}(X,P)
\int\dd x\,\ee^{-\ii s_Xx}
\delta[x-X\cos\theta-P\sin\theta]\\
&=\int\dd X\dd P\,W_{XP}(X,P)
\ee^{-\ii s_X(X\cos\theta+P\sin\theta)}.
\end{align*}
若二维 Fourier 变换定义为
\[
\widetilde W(K_X,K_P)
=\int\dd X\dd P\,
\ee^{-\ii(K_XX+K_PP)}W_{XP}(X,P),
\]
逐项比较指数即得到
\[
K_X=s_X\cos\theta,
\qquad K_P=s_X\sin\theta,
\]
所以\eqnrefs{eq:wigner_fourier_slice_theorem} 成立。二维逆 Fourier 变换为
\[
W_{XP}(X,P)
=\frac1{(2\pi)^2}
\int\dd K_X\dd K_P\,
\widetilde W(K_X,K_P)
\ee^{\ii(K_XX+K_PP)}.
\]
使用带符号径向坐标
\[
K_X=s_X\cos\theta,
\qquad K_P=s_X\sin\theta,
\qquad 0\le\theta<\pi,
\qquad -\infty<s_X<\infty,
\]
雅可比行列式为
\begin{align*}
\det\frac{\partial(K_X,K_P)}{\partial(s_X,\theta)}
&=\det\begin{pmatrix}
\cos\theta&-s_X\sin\theta\\
\sin\theta&s_X\cos\theta
\end{pmatrix}\\
&=s_X.
\end{align*}
积分测度取绝对值，因此 $\dd K_X\dd K_P=|s_X|\dd s_X\dd\theta$。负 $s_X$ 已覆盖相反径向方向，所以 $0\le\theta<\pi$ 恰好覆盖 Fourier 平面一次。代入切片定理即得到\eqnrefs{eq:wigner_filtered_backprojection}。最后，若平移不变探测器线形满足
\[
\mathcal P_\theta^{\rm obs}(x_d)
=\int\dd x\,r_X(x_d-x)\mathcal P_\theta(x),
\]
则
\begin{align*}
\widetilde{\mathcal P}_\theta^{\rm obs}(s_X)
&=\int\dd x_d\,\ee^{-\ii s_Xx_d}
\int\dd x\,r_X(x_d-x)\mathcal P_\theta(x)\\
&=\int\dd x\,\mathcal P_\theta(x)\ee^{-\ii s_Xx}
\int\dd y\,r_X(y)\ee^{-\ii s_Xy}\\
&=\widetilde r_X(s_X)\widetilde{\mathcal P}_\theta(s_X),
\end{align*}
其中第二行令 $y=x_d-x$，得到\eqnrefs{eq:wigner_tomography_detector_transfer}。
\end{exbox}



\subsection{四个参考相位如何重建一个电子相干元}
\label{sec:electron-reference-tomography-worked}

上一节说明可以在能谱测量前加入参考操作，将相干信息变为概率变化。现在具体计算：未知电子态只占据两个通道时，扫描四个参考相位能确定哪些矩阵元。先保留一般的参考幺正算符 $U_R(\varphi_R)$，其中 $\varphi_R$ 是可控相位；将探测结果 $m$ 的效应算符记为 $E_m^{\rm det}$。测得概率为
\begin{equation}
 P_m(\varphi_R)=\operatorname{Tr}[\rho_e U_R^\dagger(\varphi_R)
 E_m^{\rm det}U_R(\varphi_R)].
 \label{eq:electron-reference-effect-rotation}
\end{equation}
$U_R$ 在实验中实际作用于电子；把它移到测量算符两侧，是同一过程的 Heisenberg 表述。\emph{层析}的任务是利用一组这样的已知效应算符，从概率反推出初始 $\rho_e$。自由电子的光学参考层析已有实验实现\cite{Priebe2017NatPhoton}；下面用一个有限支持模型算清最少需要读哪些数。

假设已独立验证未知态只占据 $|0\rangle,|1\rangle$，写
\begin{equation*}
 \rho_e=\begin{pmatrix}p&c_{01}\\c_{01}^*&1-p\end{pmatrix},\qquad
 0\le p\le1,\qquad |c_{01}|^2\le p(1-p).
\end{equation*}
按 $|0\rangle,|1\rangle$ 的顺序排列基矢，$p$ 是通道 $0$ 的占据概率，可以直接测得。厄米性要求两个非对角元互为共轭，所以只剩 $c_{01}=x_{01}+\ii y_{01}$ 的实部和虚部未知。由 $\det\rho_e=p(1-p)-|c_{01}|^2\ge0$，它们必须位于半径为 $\sqrt{p(1-p)}$ 的圆盘内；这个条件保证重建结果仍是合法量子态。

令参考为经典 PINEM 脉冲，幅度 $b_R=|\beta_R|$ 已知、相位 $\varphi_R$ 可扫描。在\emph{完整双边格}上有
\begin{equation*}
 \langle m|U_R|\ell\rangle
 =J_{m-\ell}(2b_R)\ee^{\ii(m-\ell)\varphi_R}.
\end{equation*}
未知态的两通道支持并不意味着参考操作也只有两通道：参考会把电子送到其他 $m$，不能把其 $2\times2$ 子块当作幺正算符。以下只统计末态 $m=0$ 的概率，同时把所有其他输出视为互补结果。

取 $u_0=J_0(2b_R)$、$u_1=J_1(2b_R)$，利用 $J_{-1}=-J_1$，由\eqnrefs{eq:electron_tomography_linear} 得
\begin{equation}
 \begin{aligned}
 P_0(\varphi_R)&=A_R-2u_0u_1\Re[c_{01}\ee^{\ii\varphi_R}],\\
 A_R&=u_0^2p+u_1^2(1-p).
 \end{aligned}
 \label{eq:electron-reference-two-channel-fringe}
\end{equation}
相位扫描将 $x_{01}$ 乘在余弦上、$y_{01}$ 乘在正弦上。四个相位 $0,\pi/2,\pi,3\pi/2$ 给出
\begin{equation}
 x_{01}=\frac{P_0(\pi)-P_0(0)}{4u_0u_1},\qquad
 y_{01}=\frac{P_0(\pi/2)-P_0(3\pi/2)}{4u_0u_1}.
 \label{eq:electron-reference-four-phase-inversion}
\end{equation}
这些式子要求 $u_0u_1\neq0$。若参考强度恰落在相关 Bessel 零点，即使计数无限多也无法用这个输出通道确定相干；接近零点时误差按 $1/|u_0u_1|$ 放大。这就是设计矩阵的小奇异值在一个具体实验中的含义。

\begin{exbox}[四相位数据的解析重建与正性检查]
将四相位测得概率记为 $p_0,p_{\pi/2},p_\pi,p_{3\pi/2}$，已知输入占据 $p$ 和参考系数 $u_0,u_1$。求相干，并给出无需对角化的物理一致性条件。

先把条纹写成三个实系数：
\begin{equation*}
 P_0(\varphi_R)=A_R-C_R\cos\varphi_R+S_R\sin\varphi_R,
 \quad C_R=2u_0u_1x_{01},\quad S_R=2u_0u_1y_{01}.
\end{equation*}
四个设置依次读出 $A_R-C_R,A_R+S_R,A_R+C_R,A_R-S_R$。两两作差与作和即可解出
\begin{equation*}
 C_R=\frac{p_\pi-p_0}{2},\qquad
 S_R=\frac{p_{\pi/2}-p_{3\pi/2}}2,\qquad
 c_{01}=\frac{C_R+\ii S_R}{2u_0u_1}.
\end{equation*}
同一模型还要求
\begin{equation*}
 p_0+p_\pi=p_{\pi/2}+p_{3\pi/2}=2A_R,
 \qquad A_R=u_0^2p+u_1^2(1-p).
\end{equation*}
这是独立于未知相干的背景检查。正半定条件 $\det\rho_e\ge0$ 进一步给
\begin{equation*}
 C_R^2+S_R^2\le4u_0^2u_1^2p(1-p).
\end{equation*}
最后，由矩阵平方直接得到纯度
\begin{equation*}
 \operatorname{Tr}\rho_e^2=p^2+(1-p)^2
 +\frac{C_R^2+S_R^2}{2u_0^2u_1^2}.
\end{equation*}
等号边界对应纯态；$C_R=S_R=0$ 对应具有同样普通能谱的非相干混合。有限计数会使一致性关系近似成立，统计检验必须结合下述计数误差，不能把一次偏离直接当成模型失效。
\end{exbox}

有限计数时，设每个相位使用 $N$ 个独立电子，$\widehat P_0=N_0/N$。在理想完整记录条件下，$N_0$ 为二项计数，其方差为 $NP_0(1-P_0)$。四组电子独立，故
\begin{equation}
 \operatorname{Var}\widehat x_{01}
 =\frac{P_0(\pi)[1-P_0(\pi)]+P_0(0)[1-P_0(0)]}
 {16u_0^2u_1^2N},
 \label{eq:electron-reference-counting-uncertainty}
\end{equation}
$\widehat y_{01}$ 的式子把两相位换成 $\pi/2,3\pi/2$。用例题中的条纹系数，上式分子还可写成 $2A_R(1-A_R)-2C_R^2$；虚部分子把 $C_R$ 换成 $S_R$。若要求两个分量的标准差均不超过 $\varepsilon_c$，充分的每相位计数为
\begin{equation*}
 N\ge\frac{\max\{2A_R(1-A_R)-2C_R^2,\,
 2A_R(1-A_R)-2S_R^2\}}{16u_0^2u_1^2\varepsilon_c^2}.
\end{equation*}
总剂量为四组共 $4N$，另需用于确定 $p$ 和标定参考的计数。实际设计时可用 $P(1-P)\le1/4$ 给出与未知态无关的保守上界。

若存在能量误分配或漏检，应把 $E_m^{\rm det}$ 换成标定的效应算符重新构造设计矩阵。把只落在某个窗口内的计数重新归一化会引入依赖未知态的分母，一般不再服从上述线性公式。噪声导致线性反演略越出正性圆盘时，应在同一前向概率模型中作正半定约束的似然拟合，而不是把负本征值无说明地删掉。

\figref{fig:electron-reference-tomography} 显示同样的能谱如何对应不同参考条纹。
\begin{figure}[htbp]
 \centering
 \includegraphics[width=0.76\linewidth]{electron_reference_tomography.pdf}
 \caption{两通道电子态的参考读出。数值图取 $p=0.65$、$b_R=0.40$；实线使用 $c_{01}=0.20+0.15\ii$，虚线使用具有相同占据的非相干混合；四个标记点是解析反演所需的相位。误差棒表示每相位 $10^4$ 个电子的二项计数标准差，并非模拟测量数据。}
 \label{fig:electron-reference-tomography}
\end{figure}

\begin{derivation}[参考场如何让两条路径相遇]
末态 $m=0$ 可以由初态 $0$ 不交换到达，振幅为 $u_0$；也可以由初态 $1$ 损失一阶到达，振幅为 $-u_1\ee^{-\ii\varphi_R}$。两条路径进入同一个末态，故交叉项是
\begin{equation*}
 u_0(-u_1\ee^{-\ii\varphi_R})^*\rho_{01}
 +(-u_1\ee^{-\ii\varphi_R})u_0\rho_{10}
 =-2u_0u_1\Re[c_{01}\ee^{\ii\varphi_R}].
\end{equation*}
加上各路径的对角概率得到\eqnrefs{eq:electron-reference-two-channel-fringe}。相差 $\pi$ 的两次读出具有相同背景 $A_R$、相反交叉项，作差后便分别留下 $x_{01}$ 和 $y_{01}$。二项方差通过同一个线性作差公式传播，给出\eqnrefs{eq:electron-reference-counting-uncertainty}。
\end{derivation}


\section{条件测量与统计反演}

电子与光相互作用后通常存在相关。若测量电子并只保留某一个结果，与这些电子对应的光场统计也会改变。这与条件概率的思想相同，但光场还具有相干，因此需要更新整个密度算符。本节由联合散射和电子读出求条件光场态，并讨论怎样从多次电子计数推断入射光子统计。

\subsection{条件光场仪器}

现在把光场视为待保留的系统，把电子视为探针。先由联合散射算符计算电子与光的关联，再对电子施加前节定义的测量效应。对电子求偏迹后，得到光场的未归一化条件态；它的迹是电子结果的概率，除以该概率才是相应的光场态。这就从已有的联合动力学构造出作用在光场上的量子仪器。

设电子入射态在能量格基底中为
\begin{equation*}
\hat\rho_e
=\sum_{j,j'}r_{jj'}|j\rangle\langle j'|,
\qquad
\sum_j r_{jj}=1.
\end{equation*}
$r_{jj'}$ 无量纲；$r_{jj}$ 是电子初始能量格占据概率，$r_{jj'}$ 的非对角元描述电子能量相干。设光场初态为 $\hat\rho_{\rm ph}$，联合散射算符为无量纲幺正算符 $\hat S_Q$。

若电子理想测得末态格点 $\ell$，对每一个初始电子格点 $j$ 定义作用在光场希尔伯特空间上的\emph{条件偏矩阵元}
\begin{equation*}
\mathscr M_{\ell j}
\equiv {}_e\!\langle\ell|\hat S_Q|j\rangle_e.
\end{equation*}
$\mathscr M_{\ell j}$ 是无量纲光场算符；电子指标 $\ell,j$ 只标记取了哪个电子矩阵元，并不是光场算符的矩阵指标。理想结果 $\ell$ 对应的未归一化条件光场态为
\begin{equation}
\boxed{
\widetilde\rho_{\rm ph|\ell}
=\sum_{j,j'}r_{jj'}
\mathscr M_{\ell j}\hat\rho_{\rm ph}
\mathscr M_{\ell j'}^\dagger,
\qquad
P_\ell
=\operatorname{Tr}_{\rm ph}\widetilde\rho_{\rm ph|\ell}.}
\label{eq:general_electron_meter_conditional_photon_map}
\end{equation}
$\widetilde\rho_{\rm ph|\ell}$ 无量纲但迹一般小于一；其迹 $P_\ell$ 是无量纲测量概率。若 $P_\ell>0$，归一化条件态为 $\rho_{\rm ph|\ell}=\widetilde\rho_{\rm ph|\ell}/P_\ell$。

电子测量在光场上诱导的效应算符为
\begin{equation*}
\hat\Pi_\ell^{({\rm ph})}
\equiv\sum_{j,j'}r_{jj'}
\mathscr M_{\ell j'}^\dagger\mathscr M_{\ell j},
\qquad
P_\ell=\operatorname{Tr}_{\rm ph}
[\hat\Pi_\ell^{({\rm ph})}\rho_{\rm ph}].
\end{equation*}
联合幺正性与电子末态完备性保证
\begin{equation*}
\sum_\ell \hat\Pi_\ell^{({\rm ph})}=\mathbb I_{\rm ph}.
\end{equation*}
因此 $\{\hat\Pi_\ell^{({\rm ph})}\}$ 是作用在光场上的 POVM。这里的 $\hat\Pi_\ell^{({\rm ph})}$ 只是测量效应算符的统一记号，并不要求它满足投影条件 $\hat\Pi^2=\hat\Pi$；只有理想投影测量才具有这一额外性质。

QPINEM 单模交换还满足守恒律\eqnrefs{eq:qpinem_conservation}。用已经冻结的总交换量
\[
\mathcal N_{\rm ex}=\Lambda_e+N_{\rm ph}
\]
表示电子能量格阶数与光子数之和，则电子从 $j$ 变到 $\ell$ 时，光子数必须由 $n$ 变到 $n+j-\ell$。因此 $r_{jj'}$ 的第 $j-j'$ 阶电子相干会把光场相隔 $j-j'$ 个 Fock 格点的相干带入同一个电子计数概率。单能电子入射时 $r_{jj'}=\delta_{j0}\delta_{j'0}$，这些交叉项消失，光场 POVM 才退化为光子数对角形式。

真实能谱仪还具有有限效率和有限能量分辨率。前面定义的探测器响应密度
\[
R_{\rm det}(E_d|\mathcal E_\ell)
=\mathcal Q_\ell^{\rm det} r(E_d|\mathcal E_\ell)
\]
单位为 $\mathrm{J^{-1}}$，并满足 $\int\dd E_d\,R_{\rm det}=\mathcal Q_\ell^{\rm det}$。因此记录连续能量结果 $E_d$ 后的\emph{算符值概率密度}为
\begin{equation}
\boxed{
\widetilde\rho_{\rm ph|E_d}
=\sum_\ell R_{\rm det}(E_d|\mathcal E_\ell)
\sum_{j,j'}r_{jj'}
\mathscr M_{\ell j}\hat\rho_{\rm ph}
\mathscr M_{\ell j'}^\dagger.}
\label{eq:conditional_photon_detector_coarse_graining}
\end{equation}
因为 $R_{\rm det}$ 带 $\mathrm{J^{-1}}$，所以
\[
[\widetilde\rho_{\rm ph|E_d}]=\mathrm{J^{-1}},
\]
并且 $\operatorname{Tr}\widetilde\rho_{\rm ph|E_d}\,\dd E_d$ 是在区间 $[E_d,E_d+\dd E_d]$ 内得到点击的概率。若实验只接受能窗 $\mathcal W$，定义无量纲权重
\[
w_\ell(\mathcal W)
\equiv\int_{\mathcal W}\dd E_d\,
R_{\rm det}(E_d|\mathcal E_\ell),
\qquad 0\le w_\ell\le\mathcal Q_\ell^{\rm det},
\]
即可得到窗口后选择的条件量子仪器。


\subsection{单模计数核与条件仪器}\label{subsec:single-mode-number-instrument-dp19}

在单模 QPINEM 中，前面已经求出每个光子数输入对应的联合散射振幅。若只统计电子落在哪条边带，对不同的未测光子末态求和即可得到概率。若希望在电子后选择后继续使用光场，则要保留这些振幅的相位，并用 Kraus 算符表示条件更新。下面从同一组散射矩阵元写出这两种结果。

\begin{exbox}[QPINEM 单模数统计]
联合散射给出每个初始光子数 $n$ 到电子边带 $\ell$ 的振幅。若只读取电子边带而不保留末态光场相位，所需信息只是这些振幅的模平方，因此形成从光子数概率到电子计数概率的随机核；若进一步按某个电子结果后选择，光场振幅本身仍被保留，于是得到作用在光场上的条件算符。二者都由同一联合散射矩阵元导出。

取入射电子为单一能量格点 $|0\rangle_e$，光场密度算符写成
\begin{equation*}
\hat\rho_{\rm ph}
=\sum_{n,n'\ge0}\rho_{nn'}|n\rangle\langle n'|,
\qquad
p_n\equiv\rho_{nn}.
\end{equation*}
$p_n$ 是无量纲光子数概率，满足 $p_n\ge0$ 与 $\sum_n p_n=1$。

电子测得边带 $\ell$ 时，守恒量把初始光子数 $n$ 与末态光子数 $n-\ell$ 绑定。对末态光场求迹后，不同 Fock 态正交，因此电子边带概率只依赖光子数分布。定义无量纲随机核
\begin{equation}
\boxed{
P_\ell
=\sum_{n=0}^{\infty}\mathsf K_{\ell n}(|\gq|)p_n,
\qquad
\mathsf K_{\ell n}(|\gq|)
\equiv|\mathcal A_{\ell|n}(\gq)|^2.}
\label{eq:electron_spectrum_linear_kernel}
\end{equation}
$\mathsf K_{\ell n}$ 是条件概率 $P(\ell|n)$，因此
\begin{equation*}
\mathsf K_{\ell n}\ge0,
\qquad
\sum_\ell\mathsf K_{\ell n}=1
\end{equation*}
对每个固定 $n$ 成立。这里 $\gq$ 是 QPINEM 单模无量纲交换振幅。由于 $\mathcal A_{\ell|n}(\gq)$ 来自该单模有效模型的完整位移算符矩阵元，$\mathsf K_{\ell n}$ 包含 $|\gq|$ 的全部阶次；弱耦合响应是这一模型内全阶核在 $|\gq|\ll1$ 时的 Taylor 展开。

由位移算符矩阵元~\eqnrefs{eq:A_l_n_displacement}，增益边带 $\ell\ge0$ 的模型内全阶核为
\begin{equation}
\boxed{
\mathsf K_{\ell n}
=\ee^{-|\gq|^2}\frac{(n-\ell)!}{n!}|\gq|^{2\ell}
\left[L_{n-\ell}^{(\ell)}(|\gq|^2)\right]^2,
\qquad n\ge\ell,}
\label{eq:kernel_gain}
\end{equation}
而电子损失 $r_{\rm loss}=-\ell\ge0$ 时
\begin{equation}
\boxed{
\mathsf K_{-r_{\rm loss},n}
=\ee^{-|\gq|^2}\frac{n!}{(n+r_{\rm loss})!}|\gq|^{2r_{\rm loss}}
\left[L_n^{(r_{\rm loss})}(|\gq|^2)\right]^2.}
\label{eq:kernel_loss}
\end{equation}

与随机核不同，理想电子结果 $\ell$ 对应的光场条件算符定义为
\begin{equation*}
\mathscr M_\ell
\equiv{}_e\!\langle\ell|\hat S_Q|0\rangle_e
=\sum_{n\ge\max(0,\ell)}
\mathcal A_{\ell|n}|n-\ell\rangle\langle n|.
\end{equation*}
$\mathscr M_\ell$ 无量纲，作用于光场 Hilbert 空间。相应 POVM 效应和归一化条件态为
\begin{equation}
\boxed{
\hat\Pi_\ell^{({\rm ph})}
=\mathscr M_\ell^\dagger\mathscr M_\ell
=\sum_n\mathsf K_{\ell n}|n\rangle\langle n|,
\qquad
\rho_{\rm ph|\ell}
=\frac{\mathscr M_\ell\rho_{\rm ph}\mathscr M_\ell^\dagger}{P_\ell}.}
\label{eq:single_electron_photon_instrument}
\end{equation}
联合幺正性给
\[
\sum_\ell \hat\Pi_\ell^{({\rm ph})}=\mathbb I_{\rm ph},
\]
所以 $\sum_\ell P_\ell=1$。

弱交换极限把上述模型内全阶核与光子下降阶乘矩连接起来。定义
\begin{equation*}
F_{r_{\rm ph}}
\equiv\left\langle
(a^\dagger)^{r_{\rm ph}}a^{r_{\rm ph}}
\right\rangle,
\qquad r_{\rm ph}\in\mathbb N_+,
\qquad
[F_{r_{\rm ph}}]=1.
\end{equation*}
$F_{r_{\rm ph}}$ 是无量纲第 $r_{\rm ph}$ 阶下降阶乘矩。若平均光子数 $\bar n=\langle N\rangle>0$，归一化零延迟关联定义为
\begin{equation*}
g_{\rm ph}^{(r_{\rm ph})}(0)
\equiv\frac{F_{r_{\rm ph}}}{\bar n^{r_{\rm ph}}}.
\end{equation*}
由弱交换展开，电子增益第 $r_{\rm ph}$ 阶边带满足\eqnrefs{eq:qpinem-weak-factorial-moment-dp19}，从而
\begin{equation}
\boxed{
g_{\rm ph}^{(r_{\rm ph})}(0)
=\lim_{|\gq|\to0}
\frac{(r_{\rm ph}!)^2P_{+r_{\rm ph}}}{P_{+1}^{r_{\rm ph}}}.}
\label{eq:electron_factorial_correlation_readout}
\end{equation}
有限 $|\gq|$ 时应直接使用 $\mathsf K_{\ell n}$，因为高阶 $|\gq|$ 修正会同时混合多个光子数矩。


\medskip

若只读出电子能谱，光场信息通过联合散射后电子结果的条件概率进入。该关系可写成从光子统计到电子计数的随机核。 从联合散射和 Born 规则推出电子谱线性核~\eqnrefs{eq:electron_spectrum_linear_kernel}，并重写为单电子对光场的 instrument~\eqnrefs{eq:single_electron_photon_instrument}。 这里重复电子事件独立；输入光场态在单次事件内固定；核的列归一化由概率守恒保证。对初态 $|0\rangle_e\langle0|\otimes\rho_{\rm ph}$，先展开光场密度算符：
\[
\rho_{\rm ph}=\sum_{n,n'}\rho_{nn'}|n\rangle\langle n'|.
\]
联合散射作用在右矢上：
\[
\hat S_Q|0,n\rangle
=\sum_\ell\mathcal A_{\ell|n}|\ell,n-\ell\rangle.
\]
左矢对应共轭展开。投影电子结果 $\ell$ 后，概率为
\begin{align*}
P_\ell
&=\sum_{n,n'}\rho_{nn'}
\mathcal A_{\ell|n}\mathcal A_{\ell|n'}^*
\langle n'-\ell|n-\ell\rangle\\
&=\sum_{n,n'}\rho_{nn'}
\mathcal A_{\ell|n}\mathcal A_{\ell|n'}^*\delta_{nn'}\\
&=\sum_n p_n|\mathcal A_{\ell|n}|^2.
\end{align*}
因此 $\mathsf K_{\ell n}=|\mathcal A_{\ell|n}|^2$。对固定初始 $n$，联合幺正性要求
\begin{align*}
1
&=\langle0,n|\hat S_Q^\dagger\hat S_Q|0,n\rangle\\
&=\sum_\ell|\mathcal A_{\ell|n}|^2
=\sum_\ell\mathsf K_{\ell n},
\end{align*}
所以每一列 $\mathsf K_{\ell n}$ 都是归一化条件概率分布。保留末态光场而只对电子取偏矩阵元：
\[
\mathscr M_\ell
={}_e\!\langle\ell|\hat S_Q|0\rangle_e.
\]
作用在 $|n\rangle$ 上得到
\[
\mathscr M_\ell|n\rangle
=\mathcal A_{\ell|n}|n-\ell\rangle,
\]
所以
\[
\mathscr M_\ell
=\sum_n\mathcal A_{\ell|n}|n-\ell\rangle\langle n|.
\]
于是
\begin{align*}
\mathscr M_\ell^\dagger\mathscr M_\ell
&=\sum_{n,n'}
\mathcal A_{\ell|n}^*\mathcal A_{\ell|n'}
|n\rangle\langle n-\ell|n'-\ell\rangle\langle n'|\\
&=\sum_n|\mathcal A_{\ell|n}|^2|n\rangle\langle n|\\
&=\sum_n\mathsf K_{\ell n}|n\rangle\langle n|.
\end{align*}
再用电子末态完备性：
\begin{align*}
\sum_\ell\mathscr M_\ell^\dagger\mathscr M_\ell
&={}_e\!\langle0|\hat S_Q^\dagger
\left(\sum_\ell|\ell\rangle\langle\ell|\right)
\hat S_Q|0\rangle_e\\
&={}_e\!\langle0|\hat S_Q^\dagger\hat S_Q|0\rangle_e\\
&=\mathbb I_{\rm ph}.
\end{align*}

最后，把精确 Laguerre 振幅\eqnrefs{eq:A_l_n_displacement} 取模平方。对于 $\ell\ge0$，相位因子模长为一，剩余阶乘与 Laguerre 多项式直接给\eqnrefs{eq:kernel_gain}；对于 $\ell=-r_{\rm loss}\le0$，对应另一支矩阵元给\eqnrefs{eq:kernel_loss}。
\end{exbox}



\subsection{有限计数核的反演与条件操作}\label{subsec:statistics-inversion-conditioning-dp19}

设只需重建有限个光子数概率，电子探测器也只读取有限个能量区间，那么前面的计数关系成为一个有限维矩阵方程。要从电子计数求出光子概率，矩阵的不同列必须提供独立信息；若两列几乎相同，少量计数误差就会造成很大的反演误差。我们先用秩和奇异值检查这一点，再计算选定电子结果后的光场更新。

\begin{exbox}[光子统计反演与电子后选择条件操作]
\eqnrefs{eq:electron_spectrum_linear_kernel} 把有限单模光子数统计映射为电子边带分布。Fock 态、相干态和热态共享同一个条件响应核 $\mathsf K_{\ell n}$，它们的差异由输入光子数分布 $p_n$ 进入。典型状态的下降阶乘矩为
\begin{equation*}
\begin{array}{c|c|c}
\text{光态}&F_{r_{\rm ph}}&g_{\rm ph}^{(r_{\rm ph})}(0)\\ \hline
|N\rangle&N(N-1)\cdots(N-r_{\rm ph}+1)&\dfrac{N!}{(N-r_{\rm ph})!N^{r_{\rm ph}}}\\[4pt]
|\alpha_c\rangle&\bar n^{r_{\rm ph}}&1\\[2pt]
\rho_{\rm th}&r_{\rm ph}!\,\bar n^{r_{\rm ph}}&r_{\rm ph}!
\end{array}
\end{equation*}
这里用 $\alpha_c$ 表示相干态振幅，避免与精细结构常数 $\alpha$ 混用；$\bar n=|\alpha_c|^2$ 或热平均光子数均无量纲。

热态还存在一个受控的经典标度极限。令
\begin{equation*}
|\gq|\to0,
\qquad
\bar n\to\infty,
\qquad
\bar N_{\rm ex}^{(\rm th)}\equiv|\gq|^2\bar n
\end{equation*}
保持有限，则 $\bar N_{\rm ex}^{(\rm th)}$ 是无量纲平均交换强度。此时电子边带趋于
\begin{equation}
\boxed{
P_\ell^{({\rm th,cl})}
=\ee^{-2\bar N_{\rm ex}^{(\rm th)}}I_{|\ell|}(2\bar N_{\rm ex}^{(\rm th)}),
\qquad
\operatorname{Var}(\ell)=2\bar N_{\rm ex}^{(\rm th)}.}
\label{eq:thermal_skellam}
\end{equation}
该分布是对称 Skellam 型离散分布；它描述弱单光子耦合与大热占据联合缩放下的极限，只描述上述联合标度；有限 $\gq$ 时应使用单模有效模型的全阶 Laguerre 核。

实验反演只能使用有限电子边带集合 $\mathcal L_{\rm det}$，并通常把待估光子数截断到 $0\le n\le N_{\max}$。把观测边带概率组成向量 $\mathbf P$，把未知光子数分布组成 $\mathbf p$，定义截断后的无量纲设计矩阵
\begin{equation}
\boxed{
\mathsf K_{\rm inv}
\equiv[\mathsf K_{\ell n}]_{\ell\in\mathcal L_{\rm det},\,0\le n\le N_{\max}},
\qquad
\mathbf P=\mathsf K_{\rm inv}\mathbf p.}
\label{eq:single_mode_kernel_inverse_problem_dp68}
\end{equation}
这里 $\mathsf K_{\rm inv}$ 是从单模有效模型的全阶前向核中抽取出的有限数据矩阵。其奇异值分解
\begin{equation}
\boxed{
\mathsf K_{\rm inv}
=\mathsf U_K\operatorname{diag}(s_1^{(K)},s_2^{(K)},\ldots)\mathsf V_K^\dagger}
\label{eq:single_mode_kernel_svd}
\end{equation}
直接给出反演可辨识性。由于概率还满足 $\sum_n p_n=1$，真正的零空间与条件数应限制在归一化约束的切空间上。若 $s_a^{(K)}=0$，相应光子数方向不改变当前电子数据；若 $s_a^{(K)}\ll1$，伪逆沿该方向把数据噪声放大约 $1/s_a^{(K)}$。改变 $|\gq|$、电子参考相位或预整形设置相当于增加新的前向矩阵块，可以改变联合奇异值谱。

同一个联合散射还可以用来制备条件光态。设入射电子为能量梳
\begin{equation*}
|\psi_e\rangle
=\sum_j c_j|j\rangle,
\qquad
\sum_j|c_j|^2=1,
\end{equation*}
其中 $c_j$ 是无量纲复振幅。后选择电子末态 $m$ 后，定义无量纲条件光场算符
\begin{equation}
\boxed{
\mathscr C_m[\psi_e]
\equiv{}_e\!\langle m|\hat S_Q|\psi_e\rangle_e
=\sum_j c_j\,\mathscr M_{mj}.}
\label{eq:comb_conditioned_kraus-dp19}
\end{equation}
若平移协变性使 $\mathscr M_{mj}$ 只依赖净电子位移 $m-j$，则可进一步写成 $\sum_jc_j\mathscr M_{m-j}$。因此电子梳的相对相位会相干叠加不同光子增减路径，条件操作取决于振幅而不只取决于电子能量概率。

有限宽电子相位参考态 $|\theta;M\rangle$ 由 $2M+1$ 个能量格点构成。其对双边平移算符的边界误差来源于有限格宽；随着 $M$ 增大，归一化边界振幅按 $M^{-1/2}$ 衰减，而相位主瓣宽度按 $M^{-1}$ 缩小。这两个尺度分别控制“近似平移本征态”误差和可用相位分辨率。


\medskip

热光场下吸收与发射事件在经典大占据数极限可近似为两个独立 Poisson 过程，其净边带阶数是二者之差。 由两个 Poisson 生成函数的乘积推出 Skellam 分布~\eqnrefs{eq:thermal_skellam}，并给出均值与方差；这里经典标度极限使单次跃迁稀疏而总事件数有限；正负跃迁独立。热态的 Glauber--Sudarshan 表示为
\[
\rho_{\rm th}
=\int\frac{\dd^2 \alpha_c}{\pi\bar n}
\exp\!\left(-\frac{|\alpha_c|^2}{\bar n}\right)
|\alpha_c\rangle\langle\alpha_c|.
\]
在联合标度 $|\gq|\to0$、$\bar n\to\infty$ 且 $\bar N_{\rm ex}^{(\rm th)}=|\gq|^2\bar n$ 固定时，给定相干振幅 $\alpha_c$ 的电子边带趋于经典 Bessel 核
\[
P_\ell(\alpha_c)=J_\ell^2(2|\gq\alpha_c|).
\]
写 $\alpha_c=re^{\ii\varphi}$。核只依赖 $r$，所以角变量积分给 $2\pi$：
\begin{align*}
P_\ell
&=\int_0^\infty\frac{2r\dd r}{\bar n}
\exp\!\left(-\frac{r^2}{\bar n}\right)
J_\ell^2(2|\gq|r).
\end{align*}
令
\[
x=\frac{r^2}{\bar n},
\qquad
r\dd r=\frac{\bar n}{2}\dd x,
\]
得到
\[
P_\ell
=\int_0^\infty\dd x\,\ee^{-x}
J_\ell^2(2\sqrt{\bar N_{\rm ex}^{(\rm th)}x}).
\]
与其直接调用积分表，先构造边带特征函数
\[
\chi_{\rm sb}(\vartheta)
\equiv\sum_{\ell=-\infty}^{\infty}
P_\ell\ee^{\ii\ell\vartheta}.
\]
利用 Bessel 加法恒等式
\[
\sum_{\ell=-\infty}^{\infty}
J_\ell^2(z)\ee^{\ii\ell\vartheta}
=J_0\!\left(2z\sin\frac{\vartheta}{2}\right),
\]
有
\[
\chi_{\rm sb}(\vartheta)
=\int_0^\infty\dd x\,\ee^{-x}
J_0\!\left(4\sqrt{\bar N_{\rm ex}^{(\rm th)}x}\sin\frac{\vartheta}{2}\right).
\]
把 $J_0$ 展开成幂级数
\[
J_0(2\sqrt{ax})
=\sum_{m=0}^\infty
\frac{(-a)^m x^m}{(m!)^2},
\]
其中此处
\[
a=4\bar N_{\rm ex}^{(\rm th)}\sin^2\frac{\vartheta}{2}.
\]
逐项积分并使用
\[
\int_0^\infty\dd x\,\ee^{-x}x^m=m!,
\]
得到
\begin{align*}
\chi_{\rm sb}(\vartheta)
&=\sum_{m=0}^\infty\frac{(-a)^m}{m!}\\
&=\ee^{-a}\\
&=\exp\!\left[-4\bar N_{\rm ex}^{(\rm th)}\sin^2\frac{\vartheta}{2}\right]\\
&=\exp[2\bar N_{\rm ex}^{(\rm th)}(\cos\vartheta-1)].
\end{align*}
修正 Bessel 函数生成式为
\[
\exp[x(t+t^{-1})]
=\sum_{\ell=-\infty}^{\infty}I_{|\ell|}(2x)t^\ell.
\]
取 $t=\ee^{\ii\vartheta}$，
\[
\exp[2x\cos\vartheta]
=\sum_\ell I_{|\ell|}(2x)\ee^{\ii\ell\vartheta}.
\]
因此
\[
\chi_{\rm sb}(\vartheta)
=\sum_\ell
\ee^{-2\bar N_{\rm ex}^{(\rm th)}}I_{|\ell|}(2\bar N_{\rm ex}^{(\rm th)})
\ee^{\ii\ell\vartheta},
\]
逐项比较 Fourier 系数即得到\eqnrefs{eq:thermal_skellam}。由于
\[
\langle\ell\rangle
=\frac1{\ii}\chi_{\rm sb}'(0)=0,
\]
并且
\[
\langle\ell^2\rangle
=-\chi_{\rm sb}''(0)=2\bar N_{\rm ex}^{(\rm th)},
\]
所以方差为 $2\bar N_{\rm ex}^{(\rm th)}$。


\medskip

\eqnrefs{eq:single_mode_kernel_inverse_problem_dp68} 给出有限维线性前向问题
\[
\mathbf P=\mathsf K_{\rm inv}\mathbf p.
\]
对实测矩阵作奇异值分解
\[
\mathsf K_{\rm inv}=\mathsf U_K\mathsf S_K\mathsf V_K^\dagger,
\qquad
\mathsf S_K=\operatorname{diag}(s_1^{(K)},s_2^{(K)},\ldots),
\]
其中 $\mathsf U_K,\mathsf V_K$ 酉。把未知分布的微小变化展开在右奇异向量 $\mathbf v_a$ 上，
\[
\delta\mathbf p=\sum_a c_a\mathbf v_a,
\qquad
\mathsf V_K^\dagger\mathbf v_a=\mathbf e_a,
\]
则数据变化为
\begin{align*}
\delta\mathbf P
&=\mathsf K_{\rm inv}\delta\mathbf p\\
&=\mathsf U_K\mathsf S_K\mathsf V_K^\dagger
\sum_a c_a\mathbf v_a\\
&=\sum_a c_a s_a^{(K)}\mathbf u_a.
\end{align*}
因此 $s_a^{(K)}=0$ 时，沿 $\mathbf v_a$ 的参数变化严格落在前向映射零空间中；若 $s_a^{(K)}>0$，形式伪逆对相应数据分量给
\[
\delta\mathbf p_a
=\frac{\mathbf v_a\mathbf u_a^\dagger\delta\mathbf P}{s_a^{(K)}}.
\]
所以该方向上的数据噪声被放大 $1/s_a^{(K)}$。概率归一化要求 $\mathbf 1^{\mathsf T}\delta\mathbf p=0$，实际可辨识性应在这一归一化切空间内检查，即正文对零空间和条件数的限定。


\medskip

\eqnrefs{eq:general_electron_meter_conditional_photon_map} 已说明任意入射电子密度矩阵如何通过联合散射的电子偏矩阵元条件化光场。若入射电子进一步取纯的能量梳态，条件映射在振幅层可以直接重写为这些偏矩阵元的相干叠加；这一步给出\eqnrefs{eq:comb_conditioned_kraus-dp19}，并在平移协变时化成卷积。

对电子能量梳，\begin{align*}
\mathscr C_m[\psi_e]
&={}_e\!\langle m|\hat S_Q|\psi_e\rangle_e\\
&={}_e\!\langle m|\hat S_Q
\sum_jc_j|j\rangle_e\\
&=\sum_jc_j\,{}_e\!\langle m|\hat S_Q|j\rangle_e\\
&=\sum_jc_j\mathscr M_{mj}.
\end{align*}
若电子格平移协变，存在单位平移算符 $T_e$ 使 $|j\rangle=T_e^j|0\rangle$，并且联合散射只依赖净电子格位移。于是
\[
\mathscr M_{mj}=\mathscr M_{m-j},
\]
从而得到卷积形式
\[
\mathscr C_m[\psi_e]
=\sum_jc_j\mathscr M_{m-j}.
\]
每一项保留 $c_j$ 的复相位，所以能量梳相位会在条件光场振幅层发生相干干涉，而不是在概率层作非相干求和。
\end{exbox}








\subsection{聚束最优传播、能散与退相干}\label{subsec:bunching-decoherence-closure-dp47}

\eqnrefs{eq:electron_bunching_harmonic_dp46} 已把散射后能量格密度矩阵映射为传播后的密度谐波。实验上真正需要的是在给定漂移距离、输入能散和环境噪声下的可见谐波。令 $t=L/v_0$，并把理想聚束因子记为 $B_{\Delta\ell}^{(0)}(L)$。若输入电子具有小的速度涨落 $\delta v$，则同一实验距离对应传播时间涨落
\begin{equation*}
 \delta t\simeq-\frac{L}{v_0^2}\delta v.
\end{equation*}
对窄带相对论电子，$\delta v=(\partial v/\partial E)_{E_0}\delta E$。定义
\begin{equation}
\boxed{
 \sigma_t(L)
 =\frac{L}{v_0^2}
 \left|\frac{\partial v}{\partial E}\right|_{E_0}\sigma_E .}
\label{eq:arrival_time_jitter_energy_spread_dp47}
\end{equation}
若第 $\Delta\ell$ 阶密度谐波的局部角频率为 $\Delta\ell\omega$，对高斯能量分布平均后得到
\begin{equation}
\boxed{
 B_{\Delta\ell}^{(E)}(L)
 =B_{\Delta\ell}^{(0)}(L)
 \exp\!\left[-\frac12\Delta\ell^2\omega^2\sigma_t^2(L)\right].}
\label{eq:bunching_energy_spread_damping_dp47}
\end{equation}
因此高阶谐波比基频更快被有限能散压低，其衰减指数严格按 $\Delta\ell^2$ 增长。

开放系统退相干可以在同一对象上表示。若能量格的 Markov 纯退相干满足
\begin{equation}
\boxed{
 \dot\rho_{\ell\ell'}\big|_{\rm deph}
 =-\Gamma_\varphi(\ell-\ell')^2\rho_{\ell\ell'}.}
\label{eq:bunching_pure_dephasing_master_dp68}
\end{equation}
则第 $\Delta\ell$ 条副对角线统一获得因子
\begin{equation}
\boxed{
 B_{\Delta\ell}^{(\mathrm{deph})}(L)
 =B_{\Delta\ell}^{(0)}(L)
 \exp\!\left[-\Delta\ell^2\Gamma_\varphi\frac{L}{v_0}\right].}
\label{eq:bunching_dephasing_damping_dp47}
\end{equation}
能散平均与环境退相干来自不同物理机制，但都作用于同一个 $B_{\Delta\ell}$。若二者独立，实验前向模型可写成
\begin{equation}
\boxed{
 B_{\Delta\ell}^{\rm obs}(L)
 =B_{\Delta\ell}^{(0)}(L)
 \exp\!\left[-\frac12\Delta\ell^2\omega^2\sigma_t^2(L)
 -\Delta\ell^2\Gamma_\varphi\frac{L}{v_0}\right].}
\label{eq:bunching_observed_master_dp47}
\end{equation}
这一定量式把电子源品质、漂移长度和环境相干时间直接连接到可观测密度谐波。

对给定调制强度，最优漂移长度应由最大化目标谐波而不是引用固定 Talbot 分数来确定。定义
\begin{equation}
\boxed{
 L_{\Delta\ell}^*\equiv\underset{L>0}{\operatorname{argmax}}
 \left|B_{\Delta\ell}^{\rm obs}(L)\right| .}
\label{eq:optimal_bunching_length_dp47}
\end{equation}
在无噪声二阶色散模型中，$B_{\Delta\ell}^{(0)}(L)$ 对 $L_T$ 周期；加入\eqnrefs{eq:bunching_observed_master_dp47} 的指数包络后，不同 Talbot 周期不再等价，通常应选择最早达到局域极大的漂移距离。

若最终实验是用整形电子驱动目标跃迁，则\secref{subsec:shaped-electron-drive-dp19}的关系直接给
\begin{equation}
\boxed{
 \Omega_{3i}^{(e)}(L)
 =\frac{2G_{3i}}{T}B_{\Delta\ell_i}^{\rm obs}(L),
 \qquad
 \omega_{3i}=\Delta\ell_\ii\omega .}
\label{eq:target_rabi_with_decoherence_dp47}
\end{equation}
于是电子源能散和电子退相干不再是独立的“束流质量指标”，而是直接进入目标 Rabi 频率的乘法衰减。



